
\Data\ID{1003124001}
\Updated{2020-07-15 03:07:21}
\Level{A}
\SubArea{Analytic geometry in a space}
\LANGen{
\Question{
We are given a straight line \( q=\left\{[3t;2-2t;1+t]\text{, }t\in\mathbb{R}\right\} \) and four points \( A=[-6;6;-1] \), \( B=[-3;0;0] \), \( C=[0;2;1] \) and \( D=[3;0;2] \). Out of the  given points select all that lie on the straight line \( q \). (Choose the corresponding option.)}
{\( A \), \( C \), \( D \)}{\( B \), \( C \), \( D \)}{\( B \), \( C \)}{\( A \), \( B \), \( C \)}{}{}}
\LANGpl{
\Question{
Dana jest prosta  \( q=\left\{[3t;2-2t;1+t]\text{, }t\in\mathbb{R}\right\} \) oraz cztery punkty \( A=[-6;6;-1] \), \( B=[-3;0;0] \), \( C=[0;2;1] \) i \( D=[3;0;2] \). Z podanych punktów wybierz wszystkie punkty leżące na prostej  \( q \). (Wybierz właściwą opcje)}
{\( A \), \( C \), \( D \)}{\( B \), \( C \), \( D \)}{\( B \), \( C \)}{\( A \), \( B \), \( C \)}{}{}}
\LANGcs{
\Question{
Je dána přímka \( q=\left\{[3t;2-2t;1+t]\text{, }t\in\mathbb{R}\right\} \) a body \( A=[-6;6;-1] \), \( B=[-3;0;0] \), \( C=[0;2;1] \) a \( D=[3;0;2] \). Vyberte z těchto čtyř bodů všechny, které leží na přímce \( q \), a tuto možnost označte.}
{\( A \), \( C \), \( D \)}{\( B \), \( C \), \( D \)}{\( B \), \( C \)}{\( A \), \( B \), \( C \)}{}{}}
\LANGsk{
\Question{
Je daná priamka \( q=\left\{[3t;2-2t;1+t]\text{, }t\in\mathbb{R}\right\} \) a body \( A=[-6;6;-1] \), \( B=[-3;0;0] \), \( C=[0;2;1] \) a \( D=[3;0;2] \). Vyberte z týchto štyroch bodov všetky, ktoré ležia na priamke \( q \), a túto možnosť označte.}
{\( A \), \( C \), \( D \)}{\( B \), \( C \), \( D \)}{\( B \), \( C \)}{\( A \), \( B \), \( C \)}{}{}}
\LANGes{
\Question{
Dada la recta \( q=\left\{[3t;2-2t;1+t]\text{, }t\in\mathbb{R}\right\} \) y los puntos  \( A=[-6;6;-1] \), \( B=[-3;0;0] \), \( C=[0;2;1] \) y \( D=[3;0;2] \). Determina cuáles de estos cuatro puntos pertenecen a la recta q. Elige la opción correcta.}
{\( A \), \( C \), \( D \)}{\( B \), \( C \), \( D \)}{\( B \), \( C \)}{\( A \), \( B \), \( C \)}{}{}}
\End

\Data\ID{1003124002}
\Updated{2020-07-15 03:07:21}
\Level{A}
\SubArea{Analytic geometry in a space}
\LANGen{
\Question{
From the given options choose the parametric equations which describe a straight line \( p \) passing through the points \( A=[-2;0;1] \) and \( B=[2;0;-3] \).}
{\( \begin{aligned}
p\colon x&=2-t, \\
y&=0, \\
z&=-3+t;\ t\in\mathbb{R}
\end{aligned} \)}{\( \begin{aligned}
p\colon x&=2+4t, \\
y&=0, \\
z&=-3+4t;\ t\in\mathbb{R}
\end{aligned} \)}{\( \begin{aligned}
p\colon x&=2, \\
y&=0, \\
z&=-3+t;\ t\in\mathbb{R}
\end{aligned} \)}{\( \begin{aligned}
p\colon x&=2-2t, \\
y&=0, \\
z&=-3+t;\ t\in\mathbb{R}
\end{aligned} \)}{}{}}
\LANGpl{
\Question{
Wybierz równanie parametryczne prostej \( p \) przechodzącej przez punkty  \( A=[-2;0;1] \) i \( B=[2;0;-3] \).}
{\( \begin{aligned}
p\colon x&=2-t, \\
y&=0, \\
z&=-3+t;\ t\in\mathbb{R}
\end{aligned} \)}{\( \begin{aligned}
p\colon x&=2+4t, \\
y&=0, \\
z&=-3+4t;\ t\in\mathbb{R}
\end{aligned} \)}{\( \begin{aligned}
p\colon x&=2, \\
y&=0, \\
z&=-3+t;\ t\in\mathbb{R}
\end{aligned} \)}{\( \begin{aligned}
p\colon x&=2-2t, \\
y&=0, \\
z&=-3+t;\ t\in\mathbb{R}
\end{aligned} \)}{}{}}
\LANGcs{
\Question{
Z nabízených možností vyberte parametrické rovnice, které určují přímku \( p \) procházející body \( A=[-2;0;1] \) a \( B=[2;0;-3] \).}
{\( \begin{aligned}
p\colon x&=2-t, \\
y&=0, \\
z&=-3+t;\ t\in\mathbb{R}
\end{aligned} \)}{\( \begin{aligned}
p\colon x&=2+4t, \\
y&=0, \\
z&=-3+4t;\ t\in\mathbb{R}
\end{aligned} \)}{\( \begin{aligned}
p\colon x&=2, \\
y&=0, \\
z&=-3+t;\ t\in\mathbb{R}
\end{aligned} \)}{\( \begin{aligned}
p\colon x&=2-2t, \\
y&=0, \\
z&=-3+t;\ t\in\mathbb{R}
\end{aligned} \)}{}{}}
\LANGsk{
\Question{
Z ponúknutých možností vyberte parametrické rovnice, ktoré určujú priamku \( p \) prechádzajúcu bodmi \( A=[-2;0;1] \) a \( B=[2;0;-3] \).}
{\( \begin{aligned}
p\colon x&=2-t, \\
y&=0, \\
z&=-3+t;\ t\in\mathbb{R}
\end{aligned} \)}{\( \begin{aligned}
p\colon x&=2+4t, \\
y&=0, \\
z&=-3+4t;\ t\in\mathbb{R}
\end{aligned} \)}{\( \begin{aligned}
p\colon x&=2, \\
y&=0, \\
z&=-3+t;\ t\in\mathbb{R}
\end{aligned} \)}{\( \begin{aligned}
p\colon x&=2-2t, \\
y&=0, \\
z&=-3+t;\ t\in\mathbb{R}
\end{aligned} \)}{}{}}
\LANGes{
\Question{
De las siguientes ecuaciones paramétricas que definen la recta  \( p \), determina cuáles de ellas pasan por los puntos \( A=[-2;0;1] \) y \( B=[2;0;-3] \).}
{\( \begin{aligned}
p\colon x&=2-t, \\
y&=0, \\
z&=-3+t;\ t\in\mathbb{R}
\end{aligned} \)}{\( \begin{aligned}
p\colon x&=2+4t, \\
y&=0, \\
z&=-3+4t;\ t\in\mathbb{R}
\end{aligned} \)}{\( \begin{aligned}
p\colon x&=2, \\
y&=0, \\
z&=-3+t;\ t\in\mathbb{R}
\end{aligned} \)}{\( \begin{aligned}
p\colon x&=2-2t, \\
y&=0, \\
z&=-3+t;\ t\in\mathbb{R}
\end{aligned} \)}{}{}}
\End

\Data\ID{1003124003}
\Updated{2020-07-15 03:07:21}
\Level{A}
\SubArea{Analytic geometry in a space}
\LANGen{
\Question{
Find the missing coordinates of the point \( B=[x_B; y_B;-3] \) lying on a straight line \( p \) defined by the parametric equations \[\begin{aligned} p\colon x&=-1+\frac14m,\\ y&=2+m,\\ z&=5-m;\ m\in\mathbb{R}.\end{aligned} \]}
{\( B=[1;10;-3] \)}{\( B=[-3;-6;-3] \)}{\( B=[1;3;-3] \)}{\( B=[-3;6;-3] \)}{}{}}
\LANGpl{
\Question{
Wskaż brakujące współrzędne punktu \( B=[x_B; y_B;-3] \) lezącego na prostej \( p \) określonej równaniem parametrycznym \[\begin{aligned}p\colon x&=-1+\frac14m,\\ y&=2+m,\\ z&=5-m;\ m\in\mathbb{R}.\end{aligned} \]}
{\( B=[1;10;-3] \)}{\( B=[-3;-6;-3] \)}{\( B=[1;3;-3] \)}{\( B=[-3;6;-3] \)}{}{}}
\LANGcs{
\Question{
Určete chybějící souřadnice bodu \( B=[x_B; y_B;-3] \) ležícího na přímce \( p \) s parametrickým vyjádřením: \[\begin{aligned} x&=-1+\frac14m,\\ y&=2+m,\\ z&=5-m;\ m\in\mathbb{R}\end{aligned} \]}
{\( B=[1;10;-3] \)}{\( B=[-3;-6;-3] \)}{\( B=[1;3;-3] \)}{\( B=[-3;6;-3] \)}{}{}}
\LANGsk{
\Question{
Určte chýbajúce súradnice bodov \( B=[x_B; y_B;-3] \) ležiaceho na priamke \( p \) s parametrickým vyjadrením: 
\[\begin{aligned} x&=-1+\frac14m,\\ y&=2+m,\\ z&=5-m;\ m\in\mathbb{R}\end{aligned} \]}
{\( B=[1;10;-3] \)}{\( B=[-3;-6;-3] \)}{\( B=[1;3;-3] \)}{\( B=[-3;6;-3] \)}{}{}}
\LANGes{
\Question{
Halla las coordenadas del punto \( B=[x_B; y_B;-3] \) que pertenece a la recta \( p \) definida por las siguientes ecuaciones paramétricas: \[\begin{aligned} p\colon x&=-1+\frac14m,\\ y&=2+m,\\ z&=5-m;\ m\in\mathbb{R}.\end{aligned} \]}
{\( B=[1;10;-3] \)}{\( B=[-3;-6;-3] \)}{\( B=[1;3;-3] \)}{\( B=[-3;6;-3] \)}{}{}}
\End

\Data\ID{1003124004}
\Updated{2020-07-15 03:07:21}
\Level{A}
\SubArea{Analytic geometry in a space}
\LANGen{
\Question{
Find the value of a parameter \( a\in\mathbb{R} \) so that the point \( B=[1;4;5] \) lies on the straight line \( p \) defined by the parametric equations \[\begin{aligned}p\colon x&=-1+m,\\ y&=2+am,\\ z&=3+m;\ m\in\mathbb{R}. \end{aligned}\]}
{\( a=1 \)}{\( a=-1 \)}{\( a=2 \)}{no such value of \( a \) exists}{}{}}
\LANGpl{
\Question{
Wskaż wartość parametru  \( a\in\mathbb{R} \) tak, aby punkt  \( B=[1;4;5] \) leżał na prostej  \( p \) określonej równaniem parametryczny \[\begin{aligned} p\colon x&=-1+m,\\ y&=2+am,\\ z&=3+m;\ m\in\mathbb{R}. \end{aligned}\]}
{\( a=1 \)}{\( a=-1 \)}{\( a=2 \)}{wartość parametru \( a \) nie istnieje}{}{}}
\LANGcs{
\Question{
Určete hodnotu parametru \( a\in\mathbb{R} \) tak, aby bod \( B=[1;4;5] \) ležel na přímce \( p \) s parametrickým vyjádřením: \[\begin{aligned} x&=-1+m,\\ y&=2+am,\\ z&=3+m;\ m\in\mathbb{R} \end{aligned}\]}
{\( a=1 \)}{\( a=-1 \)}{\( a=2 \)}{taková hodnota \( a \) neexistuje}{}{}}
\LANGsk{
\Question{
Určte hodnotu parametra \( a\in\mathbb{R} \) tak, aby bod \( B=[1;4;5] \) ležal na priamke \( p \) s parametrickým vyjadrením: \[\begin{aligned} x&=-1+m,\\ y&=2+am,\\ z&=3+m;\ m\in\mathbb{R} \end{aligned}\]}
{\( a=1 \)}{\( a=-1 \)}{\( a=2 \)}{taká hodnota \( a \) neexistuje}{}{}}
\LANGes{
\Question{
Halla el valor del parámetro  \( a\in\mathbb{R} \) de modo que el punto \( B=[1;4;5] \) pertenezca a la recta  \( p \), la cual tiene las siguientes ecuaciones paramétricas:  \[\begin{aligned}p\colon x&=-1+m,\\ y&=2+am,\\ z&=3+m;\ m\in\mathbb{R}. \end{aligned}\]}
{\( a=1 \)}{\( a=-1 \)}{\( a=2 \)}{este valor \( a \) no existe}{}{}}
\End

\Data\ID{1003124005}
\Updated{2020-07-15 03:07:21}
\Level{A}
\SubArea{Analytic geometry in a space}
\LANGen{
\Question{
Find the value of a parameter \( a\in\mathbb{R} \) so that the point \( C=[2;0;6] \) lies on the straight line \( p \) defined by the parametric equations \[\begin{aligned}p\colon x&=-1+m,\\ y&=a+m,\\ z&=3+m;\ m\in\mathbb{R}.\end{aligned}\]}
{\( a=-3 \)}{\( a=0 \)}{\( a=-1 \)}{no such values of \(a \) exists}{}{}}
\LANGpl{
\Question{
Wskaż wartość parametru  \( a\in\mathbb{R} \) tak, aby punkt \( C=[2;0;6] \) leżał na prostej  \( p \) określonej równaniem parametrycznym \[\begin{aligned}p\colon x&=-1+m,\\ y&=a+m,\\ z&=3+m;\ m\in\mathbb{R}.\end{aligned}\]}
{\( a=-3 \)}{\( a=0 \)}{\( a=-1 \)}{taka wartość parametru \(a \) nie istnieje}{}{}}
\LANGcs{
\Question{
Určete hodnotu parametru \( a\in\mathbb{R} \) tak, aby bod \( C=[2;0;6] \) ležel na přímce \( p \) s parametrickým vyjádřením: \[\begin{aligned} x&=-1+m,\\ y&=a+m,\\ z&=3+m;\ m\in\mathbb{R}\end{aligned}\]}
{\( a=-3 \)}{\( a=0 \)}{\( a=-1 \)}{taková hodnota \(a \) neexistuje}{}{}}
\LANGsk{
\Question{
Určte hodnotu parametra \( a\in\mathbb{R} \) tak, aby bod \( C=[2;0;6] \) ležal na priamke \( p \) s parametrickým vyjadrením: \[\begin{aligned} x&=-1+m,\\ y&=a+m,\\ z&=3+m;\ m\in\mathbb{R}\end{aligned}\]}
{\( a=-3 \)}{\( a=0 \)}{\( a=-1 \)}{taká hodnota \(a \) neexistuje}{}{}}
\LANGes{
\Question{
Halla el valor del parámetro  \( a\in\mathbb{R} \) para que el punto \( C=[2;0;6] \) pertenezca a la recta \( p \), la cual tiene las siguientes ecuaciones paramétricas: \[\begin{aligned}p\colon x&=-1+m,\\ y&=a+m,\\ z&=3+m;\ m\in\mathbb{R}.\end{aligned}\]}
{\( a=-3 \)}{\( a=0 \)}{\( a=-1 \)}{este valor \(a \) no existe}{}{}}
\End

\Data\ID{1003124006}
\Updated{2020-07-15 03:07:21}
\Level{A}
\SubArea{Analytic geometry in a space}
\LANGen{
\Question{
Find the value of a parameter \( a\in\mathbb{R} \) so that the point \( D=[-2;1;1] \) lies on the straight line \( p \) defined by the parametric equations \[\begin{aligned}p\colon x&=1+m,\\ y&=-2+m,\\ z&=a+m;\ m\in\mathbb{R}. \end{aligned}\]}
{no such values of \(a \) exists}{\( a=-1 \)}{\( a=0 \)}{\( a = 1\)}{}{}}
\LANGpl{
\Question{
Wskaż wartość parametru\( a\in\mathbb{R} \) tak, aby punkt \( D=[-2;1;1] \)leżał na prostej  \( p \) określonej równaniem parametrycznym \[\begin{aligned}p\colon x&=1+m,\\ y&=-2+m,\\ z&=a+m;\ m\in\mathbb{R}. \end{aligned}\]}
{taka wartość parametru  \(a \) nie istnieje}{\( a=-1 \)}{\( a=0 \)}{\( a = 1\)}{}{}}
\LANGcs{
\Question{
Určete hodnotu parametru \( a\in\mathbb{R} \) tak, aby bod \( D=[-2;1;1] \) ležel na přímce \( p \) s parametrickým vyjádřením: \[\begin{aligned} x&=1+m,\\ y&=-2+m,\\ z&=a+m;\ m\in\mathbb{R} \end{aligned}\]}
{taková hodnota \(a \) neexistuje}{\( a=-1 \)}{\( a=0 \)}{\( a = 1\)}{}{}}
\LANGsk{
\Question{
Určte hodnotu parametra \( a\in\mathbb{R} \) tak, aby bod \( D=[-2;1;1] \) ležal na priamke \( p \) s parametrickým vyjadrením: \[\begin{aligned} x&=1+m,\\ y&=-2+m,\\ z&=a+m;\ m\in\mathbb{R} \end{aligned}\]}
{taká hodnota \(a \) neexistuje}{\( a=-1 \)}{\( a=0 \)}{\( a = 1\)}{}{}}
\LANGes{
\Question{
Halla el valor del parámetro \( a\in\mathbb{R} \) para que el punto \( D=[-2;1;1] \) pertenezca a la recta \( p \) cuyas ecuaciones paramétricas son: \[\begin{aligned}p\colon x&=1+m,\\ y&=-2+m,\\ z&=a+m;\ m\in\mathbb{R}. \end{aligned}\]}
{este valor \(a \) no existe}{\( a=-1 \)}{\( a=0 \)}{\( a = 1\)}{}{}}
\End

\Data\ID{1003164401}
\Updated{2020-07-15 03:07:21}
\Level{A}
\SubArea{Analytic geometry in a space}
\LANGen{
\Question{
Let a straight line \( p \) be defined by parametric equations:
\begin{align*}
x&=-1+2t, \\
y&=2+t, \\
z&=5-t;\ t\in\mathbb{R}.
\end{align*} 
Find the coordinates of the intersection point \( M \) of the line \( p \) with the \( xy \)-coordinate plane.}
{\( M=[9;7;0] \)}{\( M=[0;0;5] \)}{\( M=[-1;2;0] \)}{\( M=[0;0;-1] \)}{}{}}
\LANGpl{
\Question{
Dana jest prosta  \( p \) określona równaniem parametrycznym:
\begin{align*}
x&=-1+2t, \\
y&=2+t, \\
z&=5-t;\ t\in\mathbb{R}.
\end{align*} 
Wskaż współrzędne punktu \( M \), który jest punktem przecięcia prostej  \( p \) z płaszczyzną w układzie współrzędnych \( xy \).}
{\( M=[9;7;0] \)}{\( M=[0;0;5] \)}{\( M=[-1;2;0] \)}{\( M=[0;0;-1] \)}{}{}}
\LANGcs{
\Question{
Přímka \( p \) je zadána svým parametrickým vyjádřením:
\begin{align*}
x&=-1+2t, \\
y&=2+t, \\
z&=5-t;\ t\in\mathbb{R}.
\end{align*} 
Určete souřadnice bodu \( M \), ve kterém přímka \( p \) protíná souřadnou rovinu \( (xy) \).}
{\( M=[9;7;0] \)}{\( M=[0;0;5] \)}{\( M=[-1;2;0] \)}{\( M=[0;0;-1] \)}{}{}}
\LANGsk{
\Question{
Priamka \( p \) je zadaná svojím parametrickým vyjadrením:
\begin{align*}
x&=-1+2t, \\
y&=2+t, \\
z&=5-t;\ t\in\mathbb{R}.
\end{align*} 
Určte súradnice bodu \( M \), v ktorom priamka \( p \) pretína súradnicovú rovinu \( (xy) \).}
{\( M=[9;7;0] \)}{\( M=[0;0;5] \)}{\( M=[-1;2;0] \)}{\( M=[0;0;-1] \)}{}{}}
\LANGes{
\Question{
Dada la recta \( p \) cuyas ecuaciones paramétricas son:
\begin{align*}
x&=-1+2t, \\
y&=2+t, \\
z&=5-t;\ t\in\mathbb{R}.
\end{align*} 
Encuentra las coordenadas del punto \( M \), intersección de la recta \( p \) el plano de coordenadas \( xy \).}
{\( M=[9;7;0] \)}{\( M=[0;0;5] \)}{\( M=[-1;2;0] \)}{\( M=[0;0;-1] \)}{}{}}
\End

\Data\ID{1003164402}
\Updated{2020-07-15 03:07:21}
\Level{A}
\SubArea{Analytic geometry in a space}
\LANGen{
\Question{
Let a straight line \( p \) be defined by parametric equations:
\begin{align*}
x&=-1+2t, \\
y&=2+t, \\
z&=5-t;\ t\in\mathbb{R}.
\end{align*} 
Find the coordinates of the intersection point \( M \) of the line \( p \) with the \( xz \)-coordinate plane.}
{\( M=[-5;0;7] \)}{\( M=[0;2;0] \)}{\( M=[-1;0;5] \)}{\( M=[2;0;-1] \)}{}{}}
\LANGpl{
\Question{
Dana jest prosta  \( p \) określona równaniem parametrycznym:
\begin{align*}
x&=-1+2t, \\
y&=2+t, \\
z&=5-t,\ t\in\mathbb{R}.
\end{align*} 
Wskaż współrzędne punktu  \( M \), który jest punktem przecięcia prostej \( p \) z płaszczyzną w układzie współrzędnych \( xz \).}
{\( M=[-5;0;7] \)}{\( M=[0;2;0] \)}{\( M=[-1;0;5] \)}{\( M=[2;0;-1] \)}{}{}}
\LANGcs{
\Question{
Přímka \( p \) je zadána svým parametrickým vyjádřením:
\begin{align*}
x&=-1+2t, \\
y&=2+t, \\
z&=5-t,\ t\in\mathbb{R}.
\end{align*} 
Určete souřadnice bodu \( M \), ve kterém přímka \( p \) protíná souřadnou rovinu \( (xz) \).}
{\( M=[-5;0;7] \)}{\( M=[0;2;0] \)}{\( M=[-1;0;5] \)}{\( M=[2;0;-1] \)}{}{}}
\LANGsk{
\Question{
Priamka \( p \) je zadaná svojím parametrickým vyjadrením:
\begin{align*}
x&=-1+2t, \\
y&=2+t, \\
z&=5-t,\ t\in\mathbb{R}.
\end{align*} 
Určte súradnice bodu \( M \), v ktorom priamka \( p \) pretína súradnicovú rovinu \( (xz) \).}
{\( M=[-5;0;7] \)}{\( M=[0;2;0] \)}{\( M=[-1;0;5] \)}{\( M=[2;0;-1] \)}{}{}}
\LANGes{
\Question{
Dada la recta \( p \) cuyas ecuaciones paramétricas son:
\begin{align*}
x&=-1+2t, \\
y&=2+t, \\
z&=5-t;\ t\in\mathbb{R}.
\end{align*} 
Encuentra las coordenadas del punto \( M \), intersección de la recta \( p \) y el plano de coordenadas \( xz \).}
{\( M=[-5;0;7] \)}{\( M=[0;2;0] \)}{\( M=[-1;0;5] \)}{\( M=[2;0;-1] \)}{}{}}
\End

\Data\ID{1003164403}
\Updated{2020-07-15 03:07:21}
\Level{A}
\SubArea{Analytic geometry in a space}
\LANGen{
\Question{
Let a straight line $p$ be defined by parametric equations:
\begin{align*}
x&=-1+t, \\
y&=2+3t, \\
z&=5-t;\ t\in\mathbb{R}.
\end{align*} 
Find the coordinates of the intersection point \( M \) of the line \( p \) with the \( yz \)-coordinate plane.}
{\( M=[0;5;4] \)}{\( M=[-1;0;0] \)}{\( M=[0;3;-1] \)}{\( M=[1;0;0] \)}{}{}}
\LANGpl{
\Question{
Dana jest prosta  \( p \) określona równaniem parametrycznym:
\begin{align*}
x&=-1+t, \\
y&=2+3t, \\
z&=5-t;\ t\in\mathbb{R}.
\end{align*} 
Wskaż współrzędne punktu  \( M \), który jest punktem przecięcia prostej \( p \) z płaszczyzną  w układzie współrzędnych \( yz \).}
{\( M=[0;5;4] \)}{\( M=[-1;0;0] \)}{\( M=[0;3;-1] \)}{\( M=[1;0;0] \)}{}{}}
\LANGcs{
\Question{
Přímka \( p \) je zadána svým parametrickým vyjádřením:
\begin{align*}
x&=-1+t, \\
y&=2+3t, \\
z&=5-t;\ t\in\mathbb{R}.
\end{align*} 
Určete souřadnice bodu \( M \), ve kterém  přímka \( p \) protíná souřadnou rovinu \( (yz) \).}
{\( M=[0;5;4] \)}{\( M=[-1;0;0] \)}{\( M=[0;3;-1] \)}{\( M=[1;0;0] \)}{}{}}
\LANGsk{
\Question{
Priamka \( p \) je zadaná svojím parametrickým vyjadrením:
\begin{align*}
x&=-1+t, \\
y&=2+3t, \\
z&=5-t;\ t\in\mathbb{R}.
\end{align*} 
Určte súradnice bodu \( M \), v ktorom  priamka \( p \) pretína súradnicovú rovinu \( (yz) \).}
{\( M=[0;5;4] \)}{\( M=[-1;0;0] \)}{\( M=[0;3;-1] \)}{\( M=[1;0;0] \)}{}{}}
\LANGes{
\Question{
Dada la recta $p$ cuyas ecuaciones paramétricas son:
\begin{align*}
x&=-1+t, \\
y&=2+3t, \\
z&=5-t;\ t\in\mathbb{R}.
\end{align*} 
Encuentra las coordenadas del punto \( M \), intersección de la recta \( p \) y el plano de coordenadas \( yz \).}
{\( M=[0;5;4] \)}{\( M=[-1;0;0] \)}{\( M=[0;3;-1] \)}{\( M=[1;0;0] \)}{}{}}
\End

\Data\ID{1003164404}
\Updated{2020-07-15 03:07:21}
\Level{A}
\SubArea{Analytic geometry in a space}
\LANGen{
\Question{
Let a straight line \( p \) be defined by parametric equations:
\begin{align*}
x&=3+t, \\
y&=2-t, \\
z&=4;\ t\in\mathbb{R}.
\end{align*} 
Find the coordinates of the intersection point \( M \) of the line \( p \) with the \( xy \)-coordinate plane.}
{There is no such point \( M \).}{\( M=[0;0;4] \)}{\( M=[-3;2;0] \)}{\( M=[1;-1;0] \)}{}{}}
\LANGpl{
\Question{
Dana jest prosta  \( p \) określona równaniem parametrycznym:
\begin{align*}
x&=3+t, \\
y&=2-t, \\
z&=4;\ t\in\mathbb{R}.
\end{align*} 
Wskaż współrzędne punktu \( M \), który jest punktem przecięcia prostej  \( p \) z płaszczyzną w układzie współrzędnych  \( xy \).}
{Punkt \( M \) nie istnieje.}{\( M=[0;0;4] \)}{\( M=[-3;2;0] \)}{\( M=[1;-1;0] \)}{}{}}
\LANGcs{
\Question{
Přímka \( p \) je zadána svým parametrickým vyjádřením:
\begin{align*}
x&=3+t, \\
y&=2-t, \\
z&=4;\ t\in\mathbb{R}.
\end{align*} 
Určete souřadnice bodu \( M \), ve kterém přímka \( p \) protíná souřadnou rovinu \( (xy) \).}
{Takový bod \( M \) neexistuje.}{\( M=[0;0;4] \)}{\( M=[-3;2;0] \)}{\( M=[1;-1;0] \)}{}{}}
\LANGsk{
\Question{
Priamka \( p \) je zadaná svojím parametrickým vyjadrením:
\begin{align*}
x&=3+t, \\
y&=2-t, \\
z&=4;\ t\in\mathbb{R}.
\end{align*} 
Určte súradnice bodu \( M \), v ktorom priamka \( p \) pretína súradnicovú rovinu \( (xy) \).}
{Taký bod \( M \) neexistuje.}{\( M=[0;0;4] \)}{\( M=[-3;2;0] \)}{\( M=[1;-1;0] \)}{}{}}
\LANGes{
\Question{
Dada la recta \( p \) cuyas ecuaciones paramétricas son: 
\begin{align*}
x&=3+t, \\
y&=2-t, \\
z&=4;\ t\in\mathbb{R}.
\end{align*} 
Encuentra las coordenadas del punto \( M \), intersección de la recta \( p \) y el plano de coordenadas \( xy \).}
{Este punto \( M \) no existe.}{\( M=[0;0;4] \)}{\( M=[-3;2;0] \)}{\( M=[1;-1;0] \)}{}{}}
\End

\Data\ID{1003164405}
\Updated{2020-07-15 03:07:21}
\Level{A}
\SubArea{Analytic geometry in a space}
\LANGen{
\Question{
Determine whether the line \( p \) defined by parametric equations:
\begin{align*}
x&=-2+2t, \\ 
y&=1+3t, \\ 
z&=-3+3t;\ t\in\mathbb{R}
\end{align*} 
intersects any of the coordinate axis.}
{Yes, it intersects the \( y \)-axis.}{Yes, it intersects the \( x \)-axis.}{Yes, it intersects the \( z \)-axis.}{It intersects no coordinate axis.}{}{}}
\LANGpl{
\Question{
Określ czy prosta \( p \) określona równaniem parametrycznym
\begin{align*}
x&=-2+2t, \\ 
y&=1+3t, \\ 
z&=-3+3t;\ t\in\mathbb{R}
\end{align*} 
przecina oś współrzędnych.}
{Przecina oś \( y \).}{Przecina oś \( x \).}{Przecina oś \( z \).}{Nie przecina żadnej z osi.}{}{}}
\LANGcs{
\Question{
Zjistěte, zda přímka \( p \) s parametrickým vyjádřením:
\begin{align*}
x&=-2+2t, \\ 
y&=1+3t, \\ 
z&=-3+3t;\ t\in\mathbb{R}
\end{align*} 
protíná některou souřadnicovou osu.}
{Ano, protíná osu \( y \).}{Ano, protíná osu \( x \).}{Ano, protíná osu \( z \).}{Neprotíná žádnou souřadnicovou osu.}{}{}}
\LANGsk{
\Question{
Zistite, či priamka \( p \) s parametrickým vyjadrením:
\begin{align*}
x&=-2+2t, \\ 
y&=1+3t, \\ 
z&=-3+3t;\ t\in\mathbb{R}
\end{align*} 
pretína niektorú súradnicovú os.}
{Áno, pretína os \( y \).}{Áno, pretína os \( x \).}{Áno, pretína os \( z \).}{Nepretína žiadnu súradnicovú os.}{}{}}
\LANGes{
\Question{
Determina si la recta \( p \) dada por las siguientes ecuaciones paramétricas:
\begin{align*}
x&=-2+2t, \\ 
y&=1+3t, \\ 
z&=-3+3t;\ t\in\mathbb{R}
\end{align*} 
corta algún eje de coordenadas.}
{Sí, corta el eje \( y \).}{Sí, corta el eje \( x \).}{Sí, corta el eje \( z \).}{No corta ninguno de los ejes.}{}{}}
\End

\Data\ID{1003164406}
\Updated{2020-07-15 03:07:21}
\Level{A}
\SubArea{Analytic geometry in a space}
\LANGen{
\Question{
Determine whether any of the lines \( p \), \( q \) or \( r \) defined by the parametric equations given below passes through the coordinate origin.
\begin{align*}
p\colon x&=-2+4t, & q\colon x&=-5-5s, & r\colon x&=3-6u, \\             
y&=1-2t, & y&=2-2s, & y&=-\frac12+u, \\ 
z&=-3+3t;\ t\in\mathbb R & z&=5+5s;\ s\in \mathbb R & z&=2-4u;\ u\in \mathbb R 
\end{align*}}
{Yes, it's the straight line \( r \).}{Yes, it's the straight line \( p \).}{Yes, it's the straight line \( q \).}{None of the lines passes through the coordinate origin.}{}{}}
\LANGpl{
\Question{
Określ czy proste \( p \), \( q \) lub \( r \) określone równaniami parametrycznymi przechodzą przez początek układu współrzędnych.
\begin{align*}
p\colon x&=-2+4t, & q\colon x&=-5-5s, & r\colon x&=3-6u, \\             
y&=1-2t, & y&=2-2s, & y&=-\frac12+u, \\ 
z&=-3+3t;\ t\in\mathbb R & z&=5+5s;\ s\in \mathbb R & z&=2-4u;\ u\in \mathbb R 
\end{align*}}
{Tak, prosta \( r \).}{Tak, prosta \( p \).}{Tak, prosta  \( q \).}{Żadna z podanych prostych.}{}{}}
\LANGcs{
\Question{
Zjistěte, zda některá z přímek \( p \), \( q \), \( r \) zadaných parametrickými vyjádřeními prochází počátkem souřadné soustavy.
\begin{align*}
p\colon x&=-2+4t, & q\colon x&=-5-5s, & r\colon x&=3-6u, \\             
y&=1-2t, & y&=2-2s, & y&=-\frac12+u, \\ 
z&=-3+3t;\ t\in\mathbb R & z&=5+5s;\ s\in \mathbb R & z&=2-4u;\ u\in \mathbb R 
\end{align*}}
{Ano, je to přímka \( r \).}{Ano, je to přímka \( p \).}{Ano, je to přímka \( q \).}{Žádná z přímek neprochází počátkem.}{}{}}
\LANGsk{
\Question{
Zistite, či niektorá z priamok \( p \), \( q \), \( r \) zadaných parametrickým vyjadrením prechádza začiatkom súradnicovej sústavy.
\begin{align*}
p\colon x&=-2+4t, & q\colon x&=-5-5s, & r\colon x&=3-6u, \\             
y&=1-2t, & y&=2-2s, & y&=-\frac12+u, \\ 
z&=-3+3t;\ t\in\mathbb R & z&=5+5s;\ s\in \mathbb R & z&=2-4u;\ u\in \mathbb R 
\end{align*}}
{Áno, je to priamka \( r \).}{Áno, je to priamka \( p \).}{Áno, je to priamka \( q \).}{Žiadna z priamok neprechádza začiatkom.}{}{}}
\LANGes{
\Question{
Determina si alguna de las rectas \( p \), \( q \) o \( r \), dadas por las ecuaciones paramétricas, corta el origen del sistema coordenado.
\begin{align*}
p\colon x&=-2+4t, & q\colon x&=-5-5s, & r\colon x&=3-6u, \\             
y&=1-2t, & y&=2-2s, & y&=-\frac12+u, \\ 
z&=-3+3t;\ t\in\mathbb R & z&=5+5s;\ s\in \mathbb R & z&=2-4u;\ u\in \mathbb R 
\end{align*}}
{Sí, es la recta \( r \).}{Sí, es la recta \( p \).}{Sí, es la recta \( q \).}{Ninguna de estas rectas corta el origen.}{}{}}
\End

\Data\ID{1003188701}
\Updated{2020-07-15 03:07:21}
\Level{A}
\SubArea{Analytic geometry in a space}
\LANGen{
\Question{
We are given points \( A=[-2;7;0] \) and \( B=[1;7;2] \). Find the missing coordinates of the point \( C = [m; n; 4] \) so that \( C \) lies on the straight line \( AB \).}
{\( C=[4;7;4] \)}{\( C=[4;0;4] \)}{\( C=[-2;7;4] \)}{\( C=[-2;0;4] \)}{}{}}
\LANGpl{
\Question{
Dane są punkty \( A=[-2;7;0] \) i \( B=[1;7;2] \). Wskaż brakujące współrzędne punktu \( C = [m; n; 4] \) tak, aby punkt \( C \) leżał na prostej \( AB \).}
{\( C=[4;7;4] \)}{\( C=[4;0;4] \)}{\( C=[-2;7;4] \)}{\( C=[-2;0;4] \)}{}{}}
\LANGcs{
\Question{
Jsou dány body \( A=[-2;7;0] \) a \( B=[1;7;2] \). Určete chybějící souřadnice bodu \( C = [m; n; 4] \) tak, aby ležel na přímce \( AB \).}
{\( C=[4;7;4] \)}{\( C=[4;0;4] \)}{\( C=[-2;7;4] \)}{\( C=[-2;0;4] \)}{}{}}
\LANGsk{
\Question{
Sú dané body \( A=[-2;7;0] \) a \( B=[1;7;2] \). Určte chýbajúce súradnice bodu \( C = [m; n; 4] \) tak, aby ležal na priamke \( AB \).}
{\( C=[4;7;4] \)}{\( C=[4;0;4] \)}{\( C=[-2;7;4] \)}{\( C=[-2;0;4] \)}{}{}}
\LANGes{
\Question{
Dados los puntos \( A=[-2;7;0] \) y \( B=[1;7;2] \). Define las coordenadas del punto \( C = [m; n; 4] \) para que pertenezca a la recta \( AB \).}
{\( C=[4;7;4] \)}{\( C=[4;0;4] \)}{\( C=[-2;7;4] \)}{\( C=[-2;0;4] \)}{}{}}
\End

\Data\ID{1003188702}
\Updated{2020-07-15 03:07:21}
\Level{A}
\SubArea{Analytic geometry in a space}
\LANGen{
\Question{
We are given points \( A=[-2;3;0] \), \( B=[6;1;6] \) and \( C=[1;0;4] \). Find the parametric equations of a line \( p \) that passes through the point \( C \) and through the midpoint of the line segment \( AB \).}
{$\begin{aligned}
p\colon x&=1+t, \\
y&=2t, \\
z&=4-t;\ t\in\mathbb{R}
\end{aligned}$}{$\begin{aligned}
p\colon x&=1+2t, \\
y&=-t, \\
z&=4-t;\ t\in\mathbb{R}
\end{aligned}$}{$\begin{aligned}
p\colon x&=1-t, \\
y&=2t, \\
z&=4+t;\ t\in\mathbb{R}
\end{aligned}$}{$\begin{aligned}
p\colon x&=1+2t, \\
y&=t, \\
z&=4+t;\ t\in\mathbb{R}
\end{aligned}$}{}{}}
\LANGpl{
\Question{
Dane są punkty \( A=[-2;3;0] \), \( B=[6;1;6] \) i \( C=[1;0;4] \). Wskaż równanie parametryczne prostej \( p \) przechodzącej przez punkt  \( C \) oraz środek odcinka  \( AB \).}
{$\begin{aligned}
p\colon x&=1+t, \\
y&=2t, \\
z&=4-t;\ t\in\mathbb{R}
\end{aligned}$}{$\begin{aligned}
p\colon x&=1+2t, \\
y&=-t, \\
z&=4-t;\ t\in\mathbb{R}
\end{aligned}$}{$\begin{aligned}
p\colon x&=1-t, \\
y&=2t, \\
z&=4+t;\ t\in\mathbb{R}
\end{aligned}$}{$\begin{aligned}
p\colon x&=1+2t, \\
y&=t, \\
z&=4+t;\ t\in\mathbb{R}
\end{aligned}$}{}{}}
\LANGcs{
\Question{
Jsou dány body \( A=[-2;3;0] \), \( B=[6;1;6] \) a \( C=[1;0;4] \). Určete parametrické vyjádření přímky \( p \), která prochází bodem \( C \) a středem úsečky \( AB \).}
{$\begin{aligned}
p\colon x&=1+t, \\
y&=2t, \\
z&=4-t;\ t\in\mathbb{R}
\end{aligned}$}{$\begin{aligned}
p\colon x&=1+2t, \\
y&=-t, \\
z&=4-t;\ t\in\mathbb{R}
\end{aligned}$}{$\begin{aligned}
p\colon x&=1-t, \\
y&=2t, \\
z&=4+t;\ t\in\mathbb{R}
\end{aligned}$}{$\begin{aligned}
p\colon x&=1+2t, \\
y&=t, \\
z&=4+t;\ t\in\mathbb{R}
\end{aligned}$}{}{}}
\LANGsk{
\Question{
Sú dané body \( A=[-2;3;0] \), \( B=[6;1;6] \) a \( C=[1;0;4] \). Určte parametrické vyjadrenie priamky \( p \), ktorá prechádza bodom \( C \) a stredom úsečky \( AB \).}
{$\begin{aligned}
p\colon x&=1+t, \\
y&=2t, \\
z&=4-t;\ t\in\mathbb{R}
\end{aligned}$}{$\begin{aligned}
p\colon x&=1+2t, \\
y&=-t, \\
z&=4-t;\ t\in\mathbb{R}
\end{aligned}$}{$\begin{aligned}
p\colon x&=1-t, \\
y&=2t, \\
z&=4+t;\ t\in\mathbb{R}
\end{aligned}$}{$\begin{aligned}
p\colon x&=1+2t, \\
y&=t, \\
z&=4+t;\ t\in\mathbb{R}
\end{aligned}$}{}{}}
\LANGes{
\Question{
Dados los puntos \( A=[-2;3;0] \), \( B=[6;1;6] \) y \( C=[1;0;4] \). Define las ecuaciones paramétricas de la recta \( p \) que pasa por el punto \( C \) y por el centro del segmento \( AB \).}
{$\begin{aligned}
p\colon x&=1+t, \\
y&=2t, \\
z&=4-t;\ t\in\mathbb{R}
\end{aligned}$}{$\begin{aligned}
p\colon x&=1+2t, \\
y&=-t, \\
z&=4-t;\ t\in\mathbb{R}
\end{aligned}$}{$\begin{aligned}
p\colon x&=1-t, \\
y&=2t, \\
z&=4+t;\ t\in\mathbb{R}
\end{aligned}$}{$\begin{aligned}
p\colon x&=1+2t, \\
y&=t, \\
z&=4+t;\ t\in\mathbb{R}
\end{aligned}$}{}{}}
\End

\Data\ID{1003188703}
\Updated{2020-07-15 03:07:21}
\Level{A}
\SubArea{Analytic geometry in a space}
\LANGen{
\Question{
Given points \( A=[-4;1;4] \) and \( B=[4;-3;0] \), determine which of the following parametric equations do not define the line segment \( AB \).}
{$\begin{aligned}
AB\colon x&=-4+8t, \\
y&=1+4t, \\
z&=4-4t,\ t\in[0;1]
\end{aligned}$}{$\begin{aligned}
AB\colon x&=-4+8t, \\
y&=1-4t, \\
z&=4-4t,\ t\in[0;1]
\end{aligned}$}{$\begin{aligned}
AB\colon x&=4+8t, \\
y&=-3-4t, \\
z&=-4t,\ t\in[-1;0]
\end{aligned}$}{$\begin{aligned}
AB\colon x&=-4+2t, \\
y&=1-t, \\
z&=4-t,\ t\in[0;4]
\end{aligned}$}{}{}}
\LANGpl{
\Question{
Dane są punkty \( A=[-4;1;4] \) i \( B=[4;-3;0] \), wskaż równanie parametryczne, które nie określa odcinka \( AB \).}
{$\begin{aligned}
AB\colon x&=-4+8t, \\
y&=1+4t, \\
z&=4-4t;\ t\in\langle0;1\rangle
\end{aligned}$}{$\begin{aligned}
AB\colon x&=-4+8t, \\
y&=1-4t, \\
z&=4-4t;\ t\in\langle0;1\rangle
\end{aligned}$}{$\begin{aligned}
AB\colon x&=4+8t, \\
y&=-3-4t, \\
z&=-4t;\ t\in\langle-1;0\rangle
\end{aligned}$}{$\begin{aligned}
AB\colon x&=-4+2t, \\
y&=1-t, \\
z&=4-t;\ t\in\langle0;4\rangle
\end{aligned}$}{}{}}
\LANGcs{
\Question{
Jsou dány body \( A=[-4;1;4] \) a \( B=[4;-3;0] \). Zjistěte, které z následujících parametrických vyjádření není parametrickým vyjádřením úsečky \( AB \).}
{$\begin{aligned}
AB\colon x&=-4+8t, \\
y&=1+4t, \\
z&=4-4t;\ t\in\langle0;1\rangle
\end{aligned}$}{$\begin{aligned}
AB\colon x&=-4+8t, \\
y&=1-4t, \\
z&=4-4t;\ t\in\langle0;1\rangle
\end{aligned}$}{$\begin{aligned}
AB\colon x&=4+8t, \\
y&=-3-4t, \\
z&=-4t;\ t\in\langle-1;0\rangle
\end{aligned}$}{$\begin{aligned}
AB\colon x&=-4+2t, \\
y&=1-t, \\
z&=4-t;\ t\in\langle0;4\rangle
\end{aligned}$}{}{}}
\LANGsk{
\Question{
Sú dané body \( A=[-4;1;4] \) a \( B=[4;-3;0] \). Zistite, ktoré z nasledujúcich parametrických vyjadrení nie je parametrickým vyjadrením úsečky \( AB \).}
{$\begin{aligned}
AB\colon x&=-4+8t, \\
y&=1+4t, \\
z&=4-4t;\ t\in\langle0;1\rangle
\end{aligned}$}{$\begin{aligned}
AB\colon x&=-4+8t, \\
y&=1-4t, \\
z&=4-4t;\ t\in\langle0;1\rangle
\end{aligned}$}{$\begin{aligned}
AB\colon x&=4+8t, \\
y&=-3-4t, \\
z&=-4t;\ t\in\langle-1;0\rangle
\end{aligned}$}{$\begin{aligned}
AB\colon x&=-4+2t, \\
y&=1-t, \\
z&=4-t;\ t\in\langle0;4\rangle
\end{aligned}$}{}{}}
\LANGes{
\Question{
Dados los puntos \( A=[-4;1;4] \) y \( B=[4;-3;0] \). Define cuáles de las ecuaciones paramétricas no son las ecuaciones paramétricas para el segmento \( AB \).}
{$\begin{aligned}
AB\colon x&=-4+8t, \\
y&=1+4t, \\
z&=4-4t,\ t\in[0;1]
\end{aligned}$}{$\begin{aligned}
AB\colon x&=-4+8t, \\
y&=1-4t, \\
z&=4-4t,\ t\in[0;1]
\end{aligned}$}{$\begin{aligned}
AB\colon x&=4+8t, \\
y&=-3-4t, \\
z&=-4t,\ t\in[-1;0]
\end{aligned}$}{$\begin{aligned}
AB\colon x&=-4+2t, \\
y&=1-t, \\
z&=4-t,\ t\in[0;4]
\end{aligned}$}{}{}}
\End

\Data\ID{1003188704}
\Updated{2020-09-05 01:09:12}
\Level{A}
\SubArea{Analytic geometry in a space}
\LANGen{
\Question{
Given points \( A=[-4;1;4] \) and \( B=[4;-3;0] \), determine which of the following parametric equations does not define the ray \( AB \).}
{$\begin{aligned}
\mapsto AB\colon x&=-4+8t, \\
y&=1-4t, \\
z&=4-4t;\ t\in(-\infty;0]
\end{aligned}$}{$\begin{aligned}
\mapsto AB\colon x&=-4+8t, \\
y&=1-4t, \\
z&=4-4t;\ t\in[0;\infty)
\end{aligned}$}{$\begin{aligned}
\mapsto AB\colon x&=-4+2t, \\
y&=1-t, \\
z&=4-t;\ t\in[0;\infty)
\end{aligned}$}{$\begin{aligned}
\mapsto AB\colon x&=-4-8t, \\
y&=1+4t, \\
z&=4+4t;\ t\in(-\infty;0]
\end{aligned}$}{}{}}
\LANGpl{
\Question{
Dane są punkty  \( A=[-4;1;4] \) i \( B=[4;-3;0] \), wskaż równanie parametryczne, które nie określa półprostej \( AB \).}
{$\begin{aligned}
\mapsto AB\colon x&=-4+8t, \\
y&=1-4t, \\
z&=4-4t;\ t\in(-\infty;0\rangle
\end{aligned}$}{$\begin{aligned}
\mapsto AB\colon x&=-4+8t, \\
y&=1-4t, \\
z&=4-4t;\ t\in\langle0;\infty)
\end{aligned}$}{$\begin{aligned}
\mapsto AB\colon x&=-4+2t, \\
y&=1-t, \\
z&=4-t;\ t\in\langle0;\infty)
\end{aligned}$}{$\begin{aligned}
\mapsto AB\colon x&=-4-8t, \\
y&=1+4t, \\
z&=4+4t;\ t\in(-\infty;0\rangle
\end{aligned}$}{}{}}
\LANGcs{
\Question{
Jsou dány body \( A=[-4;1;4] \) a \( B=[4;-3;0] \). Zjistěte, které z následujících parametrických vyjádření není vyjádřením polopřímky \( AB \).}
{$\begin{aligned}
\mapsto AB\colon x&=-4+8t, \\
y&=1-4t, \\
z&=4-4t;\ t\in(-\infty;0\rangle
\end{aligned}$}{$\begin{aligned}
\mapsto AB\colon x&=-4+8t, \\
y&=1-4t, \\
z&=4-4t;\ t\in\langle0;\infty)
\end{aligned}$}{$\begin{aligned}
\mapsto AB\colon x&=-4+2t, \\
y&=1-t, \\
z&=4-t;\ t\in\langle0;\infty)
\end{aligned}$}{$\begin{aligned}
\mapsto AB\colon x&=-4-8t, \\
y&=1+4t, \\
z&=4+4t;\ t\in(-\infty;0\rangle
\end{aligned}$}{}{}}
\LANGsk{
\Question{
Sú dané body \( A=[-4;1;4] \) a \( B=[4;-3;0] \). Zistite, ktoré z nasledujúcich parametrických vyjadrení nie je vyjadrením polpriamky \( AB \).}
{$\begin{aligned}
\mapsto AB\colon x&=-4+8t, \\
y&=1-4t, \\
z&=4-4t;\ t\in(-\infty;0\rangle
\end{aligned}$}{$\begin{aligned}
\mapsto AB\colon x&=-4+8t, \\
y&=1-4t, \\
z&=4-4t;\ t\in\langle0;\infty)
\end{aligned}$}{$\begin{aligned}
\mapsto AB\colon x&=-4+2t, \\
y&=1-t, \\
z&=4-t;\ t\in\langle0;\infty)
\end{aligned}$}{$\begin{aligned}
\mapsto AB\colon x&=-4-8t, \\
y&=1+4t, \\
z&=4+4t;\ t\in(-\infty;0\rangle
\end{aligned}$}{}{}}
\LANGes{
\Question{
Dados los puntos \( A=[-4;1;4] \) y \( B=[4;-3;0] \). Determina cuáles de las siguientes ecuaciones paramétricas no son las ecuaciones de semirecta \( AB \).}
{$\begin{aligned}
\mapsto AB\colon x&=-4+8t, \\
y&=1-4t, \\
z&=4-4t;\ t\in(-\infty;0]
\end{aligned}$}{$\begin{aligned}
\mapsto AB\colon x&=-4+8t, \\
y&=1-4t, \\
z&=4-4t;\ t\in[0;\infty)
\end{aligned}$}{$\begin{aligned}
\mapsto AB\colon x&=-4+2t, \\
y&=1-t, \\
z&=4-t;\ t\in[0;\infty)
\end{aligned}$}{$\begin{aligned}
\mapsto AB\colon x&=-4-8t, \\
y&=1+4t, \\
z&=4+4t;\ t\in(-\infty;0]
\end{aligned}$}{}{}}
\End

\Data\ID{1003188801}
\Updated{2020-07-15 03:07:21}
\Level{A}
\SubArea{Analytic geometry in a space}
\LANGen{
\Question{
We are given points \( A=[2;4;0] \), \( B=[4;-1;1] \) and \( C=[0;1;1] \). From the following list, choose the parametric equations which represent a plane \( \rho \) defined by the points \( A \), \( B \), and \( C \).}
{$\begin{aligned}
\rho\colon x&=4+2t+2s, \\
y&=-1-t-5s, \\
z&=1+s;\ t,s\in\mathbb{R}
\end{aligned}$}{$\begin{aligned}
\rho\colon x&=4+4t+2s, \\
y&=-1-2t-5s, \\
z&=1+t+s;\ t,s\in\mathbb{R}
\end{aligned}$}{$\begin{aligned}
\rho\colon x&=2t+4s, \\
y&=1-t-2s, \\
z&=1;\ t,s\in\mathbb{R}
\end{aligned}$}{$\begin{aligned}
\rho\colon x&=2t-2s, \\
y&=1-5t+5s, \\
z&=1+t-s;\ t,s\in\mathbb{R}
\end{aligned}$}{}{}}
\LANGpl{
\Question{
Dane są punkty \( A=[2;4;0] \), \( B=[4;-1;1] \) i \( C=[0;1;1] \). Wskaż równanie parametryczne płaszczyzny  \( \rho \) określonej punktami  \( A \), \( B \), i \( C \).}
{$\begin{aligned}
\rho\colon x&=4+2t+2s, \\
y&=-1-t-5s, \\
z&=1+s;\ t,s\in\mathbb{R}
\end{aligned}$}{$\begin{aligned}
\rho\colon x&=4+4t+2s, \\
y&=-1-2t-5s, \\
z&=1+t+s;\ t,s\in\mathbb{R}
\end{aligned}$}{$\begin{aligned}
\rho\colon x&=2t+4s, \\
y&=1-t-2s, \\
z&=1;\ t,s\in\mathbb{R}
\end{aligned}$}{$\begin{aligned}
\rho\colon x&=2t-2s, \\
y&=1-5t+5s, \\
z&=1+t-s;\ t,s\in\mathbb{R}
\end{aligned}$}{}{}}
\LANGcs{
\Question{
Jsou dány body \( A=[2;4;0] \), \( B=[4;-1;1] \) a \( C=[0;1;1] \). Z nabízených možností vyberte parametrické rovnice, které vyjadřují rovinu \( \rho \) danou body \( A \), \( B \) a \( C \).}
{$\begin{aligned}
\rho\colon x&=4+2t+2s, \\
y&=-1-t-5s, \\
z&=1+s;\ t,s\in\mathbb{R}
\end{aligned}$}{$\begin{aligned}
\rho\colon x&=4+4t+2s, \\
y&=-1-2t-5s, \\
z&=1+t+s;\ t,s\in\mathbb{R}
\end{aligned}$}{$\begin{aligned}
\rho\colon x&=2t+4s, \\
y&=1-t-2s, \\
z&=1;\ t,s\in\mathbb{R}
\end{aligned}$}{$\begin{aligned}
\rho\colon x&=2t-2s, \\
y&=1-5t+5s, \\
z&=1+t-s;\ t,s\in\mathbb{R}
\end{aligned}$}{}{}}
\LANGsk{
\Question{
Sú dané body \( A=[2;4;0] \), \( B=[4;-1;1] \) a \( C=[0;1;1] \). Z ponúknutých možností vyberte parametrické rovnice, ktoré vyjadrujú rovinu \( \rho \) danú bodmi \( A \), \( B \) a \( C \).}
{$\begin{aligned}
\rho\colon x&=4+2t+2s, \\
y&=-1-t-5s, \\
z&=1+s;\ t,s\in\mathbb{R}
\end{aligned}$}{$\begin{aligned}
\rho\colon x&=4+4t+2s, \\
y&=-1-2t-5s, \\
z&=1+t+s;\ t,s\in\mathbb{R}
\end{aligned}$}{$\begin{aligned}
\rho\colon x&=2t+4s, \\
y&=1-t-2s, \\
z&=1;\ t,s\in\mathbb{R}
\end{aligned}$}{$\begin{aligned}
\rho\colon x&=2t-2s, \\
y&=1-5t+5s, \\
z&=1+t-s;\ t,s\in\mathbb{R}
\end{aligned}$}{}{}}
\LANGes{
\Question{
Dados los puntos \( A=[2;4;0] \), \( B=[4;-1;1] \) y \( C=[0;1;1] \). De las siguientes posibilidades elige las ecuaciones paramétricas que representan el plano \( \rho \) dado por los puntos \( A \), \( B \), y \( C \).}
{$\begin{aligned}
\rho\colon x&=4+2t+2s, \\
y&=-1-t-5s, \\
z&=1+s;\ t,s\in\mathbb{R}
\end{aligned}$}{$\begin{aligned}
\rho\colon x&=4+4t+2s, \\
y&=-1-2t-5s, \\
z&=1+t+s;\ t,s\in\mathbb{R}
\end{aligned}$}{$\begin{aligned}
\rho\colon x&=2t+4s, \\
y&=1-t-2s, \\
z&=1;\ t,s\in\mathbb{R}
\end{aligned}$}{$\begin{aligned}
\rho\colon x&=2t-2s, \\
y&=1-5t+5s, \\
z&=1+t-s;\ t,s\in\mathbb{R}
\end{aligned}$}{}{}}
\End

\Data\ID{1003188802}
\Updated{2020-07-15 03:07:21}
\Level{A}
\SubArea{Analytic geometry in a space}
\LANGen{
\Question{
Find the missing coordinates of the points\( M=[2;m;0] \) and \( N=[0;3;n] \) so that they lie on a plane \( \rho \) defined by the following parametric equations:
\begin{align*}
\rho\colon x&=4+2s, \\
y&=-1-2t, \\
z&=1+t+s;\ t,s\in\mathbb{R}
\end{align*}
Choose the option in which values of both \( m \) and \( n \) are correct.}
{\( m=-1 \), \( n=-3 \)}{\( m=-1 \), \( n=3 \)}{\( m=1 \), \( n=-3 \)}{\( m=1 \), \( n=3 \)}{}{}}
\LANGpl{
\Question{
Wskaż brakujące współrzędne punktów \( M=[2;m;0] \) i  \( N=[0;3;n] \) tak, aby leżały na płaszczyźnie  \( \rho \) określonej równaniem parametrycznym:
\begin{align*}
\rho\colon x&=4+2s, \\
y&=-1-2t, \\
z&=1+t+s;\ t,s\in\mathbb{R}
\end{align*}
Wybierz opcję, w której wartość \( m \) i  \( n \) są poprawne.}
{\( m=-1 \), \( n=-3 \)}{\( m=-1 \), \( n=3 \)}{\( m=1 \), \( n=-3 \)}{\( m=1 \), \( n=3 \)}{}{}}
\LANGcs{
\Question{
Určete chybějící souřadnice bodů \( M=[2;m;0] \) a \( N=[0;3;n] \) tak, aby tyto body ležely v rovině \( \rho \) dané parametrickým vyjádřením:
\begin{align*}
\rho\colon x&=4+2s, \\
y&=-1-2t, \\
z&=1+t+s;\ t,s\in\mathbb{R}
\end{align*}
Vyberte možnost, ve které jsou obě hodnoty \( m \) a \( n \) správně.}
{\( m=-1 \), \( n=-3 \)}{\( m=-1 \), \( n=3 \)}{\( m=1 \), \( n=-3 \)}{\( m=1 \), \( n=3 \)}{}{}}
\LANGsk{
\Question{
Určte chýbajúce súradnice bodov \( M=[2;m;0] \) a \( N=[0;3;n] \) tak, aby tieto body ležali v rovine \( \rho \) danej parametrickým vyjadrením:
\begin{align*}
\rho\colon x&=4+2s, \\
y&=-1-2t, \\
z&=1+t+s;\ t,s\in\mathbb{R}
\end{align*}
Vyberte možnosť, v ktorom sú obidve hodnoty \( m \) a \( n \) správne.}
{\( m=-1 \), \( n=-3 \)}{\( m=-1 \), \( n=3 \)}{\( m=1 \), \( n=-3 \)}{\( m=1 \), \( n=3 \)}{}{}}
\LANGes{
\Question{
Determina las coordenadas que les faltan a los puntos \( M=[2;m;0] \) y \( N=[0;3;n] \) para que éstos pertenezcan al plano \( \rho \) dado por la ecuación:
\begin{align*}
\rho\colon x&=4+2s, \\
y&=-1-2t, \\
z&=1+t+s;\ t,s\in\mathbb{R}
\end{align*}
Elige cuáles de los valores \( m \) y \( n \) son correctos.}
{\( m=-1 \), \( n=-3 \)}{\( m=-1 \), \( n=3 \)}{\( m=1 \), \( n=-3 \)}{\( m=1 \), \( n=3 \)}{}{}}
\End

\Data\ID{1003188803}
\Updated{2020-07-15 03:07:21}
\Level{A}
\SubArea{Analytic geometry in a space}
\LANGen{
\Question{
A plane \( \rho \) is defined by the point \( A=[3;1;1] \) and a straight line \( p \) defined by the following parametric equations:
\begin{align*}
p\colon x&=4+4t, \\
y&=-1-2t, \\
z&=1+t;\ t\in\mathbb{R}
\end{align*}
Find the parametric equations of the plane \( \rho \).}
{$\begin{aligned}
\rho\colon x&=4+4t+s, \\
y&=-1-2t-2s, \\
z&=1+t;\ t,s\in\mathbb{R}
\end{aligned}$}{$\begin{aligned}
\rho\colon x&=4+4t+3s, \\
y&=-1-2t+s, \\
z&=1+t+s;\ t,s\in\mathbb{R}
\end{aligned}$}{$\begin{aligned}
\rho\colon x&=3+4t+4s, \\
y&=1-2t-s, \\
z&=1+t+s;\ t,s\in\mathbb{R}
\end{aligned}$}{$\begin{aligned}
\rho\colon x&=3+4t-4s, \\
y&=1-2t+2s, \\
z&=1+t-s;\ t,s\in\mathbb{R}
\end{aligned}$}{}{}}
\LANGpl{
\Question{
Płaszczyzna  \( \rho \) jest określona punktem \( A=[3;1;1] \) i prostą \( p \) określoną równaniem parametrycznym::
\begin{align*}
p\colon x&=4+4t, \\
y&=-1-2t, \\
z&=1+t;\ t\in\mathbb{R}
\end{align*}
Wskaż równanie parametryczne płaszczyzny \( \rho \).}
{$\begin{aligned}
\rho\colon x&=4+4t+s, \\
y&=-1-2t-2s, \\
z&=1+t;\ t,s\in\mathbb{R}
\end{aligned}$}{$\begin{aligned}
\rho\colon x&=4+4t+3s, \\
y&=-1-2t+s, \\
z&=1+t+s;\ t,s\in\mathbb{R}
\end{aligned}$}{$\begin{aligned}
\rho\colon x&=3+4t+4s, \\
y&=1-2t-s, \\
z&=1+t+s;\ t,s\in\mathbb{R}
\end{aligned}$}{$\begin{aligned}
\rho\colon x&=3+4t-4s, \\
y&=1-2t+2s, \\
z&=1+t-s;\ t,s\in\mathbb{R}
\end{aligned}$}{}{}}
\LANGcs{
\Question{
Rovina \( \rho \) je určena bodem \( A=[3;1;1] \) a přímkou \( p \) danou parametrickým vyjádřením:
\begin{align*}
p\colon x&=4+4t, \\
y&=-1-2t, \\
z&=1+t;\ t\in\mathbb{R}
\end{align*}
Určete parametrické vyjádření roviny \( \rho \).}
{$\begin{aligned}
\rho\colon x&=4+4t+s, \\
y&=-1-2t-2s, \\
z&=1+t;\ t,s\in\mathbb{R}
\end{aligned}$}{$\begin{aligned}
\rho\colon x&=4+4t+3s, \\
y&=-1-2t+s, \\
z&=1+t+s;\ t,s\in\mathbb{R}
\end{aligned}$}{$\begin{aligned}
\rho\colon x&=3+4t+4s, \\
y&=1-2t-s, \\
z&=1+t+s;\ t,s\in\mathbb{R}
\end{aligned}$}{$\begin{aligned}
\rho\colon x&=3+4t-4s, \\
y&=1-2t+2s, \\
z&=1+t-s;\ t,s\in\mathbb{R}
\end{aligned}$}{}{}}
\LANGsk{
\Question{
Rovina \( \rho \) je určená bodom \( A=[3;1;1] \) a priamkou \( p \) danou parametrickým vyjadrením:
\begin{align*}
p\colon x&=4+4t, \\
y&=-1-2t, \\
z&=1+t;\ t\in\mathbb{R}
\end{align*}
Určte parametrické vyjadrenie roviny \( \rho \).}
{$\begin{aligned}
\rho\colon x&=4+4t+s, \\
y&=-1-2t-2s, \\
z&=1+t;\ t,s\in\mathbb{R}
\end{aligned}$}{$\begin{aligned}
\rho\colon x&=4+4t+3s, \\
y&=-1-2t+s, \\
z&=1+t+s;\ t,s\in\mathbb{R}
\end{aligned}$}{$\begin{aligned}
\rho\colon x&=3+4t+4s, \\
y&=1-2t-s, \\
z&=1+t+s;\ t,s\in\mathbb{R}
\end{aligned}$}{$\begin{aligned}
\rho\colon x&=3+4t-4s, \\
y&=1-2t+2s, \\
z&=1+t-s;\ t,s\in\mathbb{R}
\end{aligned}$}{}{}}
\LANGes{
\Question{
Dado el plano \( \rho \) por el punto \( A=[3;1;1] \) y por la recta \( p \) cuyas ecuaciones paramétricas son:
\begin{align*}
p\colon x&=4+4t, \\
y&=-1-2t, \\
z&=1+t;\ t\in\mathbb{R}
\end{align*}
Determina las ecuaciones paramétricas del plano \( \rho \).}
{$\begin{aligned}
\rho\colon x&=4+4t+s, \\
y&=-1-2t-2s, \\
z&=1+t;\ t,s\in\mathbb{R}
\end{aligned}$}{$\begin{aligned}
\rho\colon x&=4+4t+3s, \\
y&=-1-2t+s, \\
z&=1+t+s;\ t,s\in\mathbb{R}
\end{aligned}$}{$\begin{aligned}
\rho\colon x&=3+4t+4s, \\
y&=1-2t-s, \\
z&=1+t+s;\ t,s\in\mathbb{R}
\end{aligned}$}{$\begin{aligned}
\rho\colon x&=3+4t-4s, \\
y&=1-2t+2s, \\
z&=1+t-s;\ t,s\in\mathbb{R}
\end{aligned}$}{}{}}
\End

\Data\ID{1003188903}
\Updated{2020-07-15 03:07:21}
\Level{A}
\SubArea{Analytic geometry in a space}
\LANGen{
\Question{
Determine the relative position of the plane \( \rho \) with general equation \( 2x-y+z-2=0 \) and the straight line \( p \) with parametric equations:
\[ \begin{aligned}
x&=2-t, \\
y&=5-2t, \\
z&=3;\ t\in\mathbb{R}.
\end{aligned} \]}
{\( p \subset \rho \)}{\( p\parallel\rho\text{, }p\not{\!\!\subset} \rho \)}{\( p \) is intersecting the plane \( \rho \)}{}{}{}}
\LANGpl{
\Question{
Określ wzajemne położenie płaszczyzny \( \rho \)  i prostej \( p \) określonej równaniem parametrycznym \( 2x-y+z-2=0 \)
\[ \begin{aligned}
x&=2-t, \\
y&=5-2t, \\
z&=3;\ t\in\mathbb{R}.
\end{aligned} \]}
{\( p \subset \rho \)}{\( p\parallel\rho\text{, }p\not{\!\!\subset} \rho \)}{\( p \) przecina płaszczyznę \( \rho \)}{}{}{}}
\LANGcs{
\Question{
Vyšetřete vzájemnou polohu roviny \( \rho \): \( 2x-y+z-2=0 \) a přímky \( p \) dané parametrickými rovnicemi:
\[ \begin{aligned}
x&=2-t, \\
y&=5-2t, \\
z&=3;\ t\in\mathbb{R}.
\end{aligned} \]}
{\( p \subset \rho \)}{\( p\parallel\rho\text{, }p\not{\!\!\subset} \rho \)}{\( p \) protíná rovinu \( \rho \)}{}{}{}}
\LANGsk{
\Question{
Zistite vzájomnú polohu roviny \( \rho \): \( 2x-y+z-2=0 \) a priamky \( p \) danej parametrickými rovnicami:
\[ \begin{aligned}
x&=2-t, \\
y&=5-2t, \\
z&=3;\ t\in\mathbb{R}.
\end{aligned} \]}
{\( p \subset \rho \)}{\( p\parallel\rho\text{, }p\not{\!\!\subset} \rho \)}{\( p \) pretína rovinu \( \rho \)}{}{}{}}
\LANGes{
\Question{
Determina la posición relativa del plano \( \rho \) y la recta \( p \) cuyas ecuaciones paramétricas son: \( 2x-y+z-2=0 \) 
\[ \begin{aligned}
x&=2-t, \\
y&=5-2t, \\
z&=3;\ t\in\mathbb{R}.
\end{aligned} \]}
{\( p \subset \rho \)}{\( p\parallel\rho\text{, }p\not{\!\!\subset} \rho \)}{\( p \) corta el plano \( \rho \)}{}{}{}}
\End

\Data\ID{1003188904}
\Updated{2020-07-15 03:07:21}
\Level{A}
\SubArea{Analytic geometry in a space}
\LANGen{
\Question{
Determine the relative position of the plane \( \rho \) with general equation \( 7x-2y+z-2=0 \) and the straight line \( p \) with parametric equations:
\[ \begin{aligned}
x&=3+t, \\
y&=-5-2t, \\
z&=3-11t;\ t\in\mathbb{R}.
\end{aligned} \]}
{\( p\parallel \rho\text{, }p\not{\!\!\subset}\rho \)}{\( p \subset \rho \)}{\( p \) is intersecting the plane \( \rho \)}{}{}{}}
\LANGpl{
\Question{
Określ wzajemne położenie płaszczyzny \( \rho \) o równaniu ogólnym \( 7x-2y+z-2=0 \) i prostej  \( p \) określonej równaniem parametrycznym:
\[ \begin{aligned}
x&=3+t, \\
y&=-5-2t, \\
z&=3-11t;\ t\in\mathbb{R}.
\end{aligned} \]}
{\( p\parallel \rho\text{, }p\not{\!\!\subset}\rho \)}{\( p \subset \rho \)}{\( p \) przecina płaszczyznę  \( \rho \)}{}{}{}}
\LANGcs{
\Question{
Vyšetřete vzájemnou polohu roviny \( \rho \): \( 7x-2y+z-2=0 \) a přímky \( p \) dané parametrickými rovnicemi:
\[ \begin{aligned}
x&=3+t, \\
y&=-5-2t, \\
z&=3-11t;\ t\in\mathbb{R}.
\end{aligned} \]}
{\( p\parallel \rho\text{, }p\not{\!\!\subset}\rho \)}{\( p \subset \rho \)}{\( p \) protíná rovinu \( \rho \)}{}{}{}}
\LANGsk{
\Question{
Vyšetrite vzájomnú polohu roviny \( \rho \): \( 7x-2y+z-2=0 \) a priamky \( p \) danej parametrickými rovnicami:
\[ \begin{aligned}
x&=3+t, \\
y&=-5-2t, \\
z&=3-11t;\ t\in\mathbb{R}.
\end{aligned} \]}
{\( p\parallel \rho\text{, }p\not{\!\!\subset}\rho \)}{\( p \subset \rho \)}{\( p \) pretína rovinu \( \rho \)}{}{}{}}
\LANGes{
\Question{
Determina la posición relativa del plano \( \rho \) cuyas ecuaciones son \( 7x-2y+z-2=0 \) y la recta \( p \) cuyas ecuaciones paramétricas son:
\[ \begin{aligned}
x&=3+t, \\
y&=-5-2t, \\
z&=3-11t;\ t\in\mathbb{R}.
\end{aligned} \]}
{\( p\parallel \rho\text{, }p\not{\!\!\subset}\rho \)}{\( p \subset \rho \)}{\( p \) corta el plano \( \rho \)}{}{}{}}
\End

\Data\ID{1003188905}
\Updated{2020-07-15 03:07:21}
\Level{A}
\SubArea{Analytic geometry in a space}
\LANGen{
\Question{
Determine the relative position of the plane \( \rho \) with general equation \( 5x-4y+z-4=0 \) and the straight line \( p \) with parametric equations:
\[ \begin{aligned}
x&=-1+t,\\
y&=2-2t,\\
z&=3+t;\ t\in\mathbb{R}.
\end{aligned} \]}
{\( p \) is intersecting \( \rho \)}{\( p\parallel \rho\text{, } p\not{\!\!\subset}\rho \)}{\( p \subset \rho \)}{}{}{}}
\LANGpl{
\Question{
Określ wzajemne położenie płaszczyzny \( \rho \) o równaniu ogólnym \( 5x-4y+z-4=0 \) i prostej  \( p \) określonej równaniem parametrycznym:
\[ \begin{aligned}
x&=-1+t,\\
y&=2-2t,\\
z&=3+t;\ t\in\mathbb{R}.
\end{aligned} \]}
{\( p \) przecina\( \rho \)}{\( p\parallel \rho\text{, } p\not{\!\!\subset}\rho \)}{\( p \subset \rho \)}{}{}{}}
\LANGcs{
\Question{
Vyšetřete vzájemnou polohu roviny \( \rho \): \( 5x-4y+z-4=0 \) a přímky \( p \) dané parametrickými rovnicemi:
\[ \begin{aligned}
x&=-1+t,\\
y&=2-2t,\\
z&=3+t;\ t\in\mathbb{R}.
\end{aligned} \]}
{\( p \) protíná rovinu \( \rho \)}{\( p\parallel \rho\text{, } p\not{\!\!\subset}\rho \)}{\( p \subset \rho \)}{}{}{}}
\LANGsk{
\Question{
Vyšetrite vzájomnú polohu roviny \( \rho \): \( 5x-4y+z-4=0 \) a priamky \( p \) danej parametrickými rovnicami:
\[ \begin{aligned}
x&=-1+t,\\
y&=2-2t,\\
z&=3+t;\ t\in\mathbb{R}.
\end{aligned} \]}
{\( p \) pretína rovinu \( \rho \)}{\( p\parallel \rho\text{, } p\not{\!\!\subset}\rho \)}{\( p \subset \rho \)}{}{}{}}
\LANGes{
\Question{
Determina la posición relativa del plano \( \rho \) con la ecuación general \( 5x-4y+z-4=0 \) y la recta \( p \) cuyas ecuaciones paraméticas son:
\[ \begin{aligned}
x&=-1+t,\\
y&=2-2t,\\
z&=3+t;\ t\in\mathbb{R}.
\end{aligned} \]}
{\( p \) corta \( \rho \)}{\( p\parallel \rho\text{, } p\not{\!\!\subset}\rho \)}{\( p \subset \rho \)}{}{}{}}
\End

\Data\ID{1003188906}
\Updated{2020-07-15 03:07:21}
\Level{A}
\SubArea{Analytic geometry in a space}
\LANGen{
\Question{
Let there be planes \( \alpha \), \( \beta \), \( \gamma \) and \( \delta \) defined by their general equations:  
\[ \begin{aligned}
&\alpha\colon \frac23x-4y+6z-\frac83=0; \\
&\beta\colon x-2y+3z-4=0; \\
&\gamma\colon 2x-12y+18z-4 =0; \\
&\delta\colon x-6y+9z-4 =0.
\end{aligned} \]
Out of the following statements, select the one that is not true.}
{\( \alpha \parallel\delta\text{, }\alpha\neq\delta \)}{Planes \( \beta \) and \( \delta \) are intersecting.}{\( \gamma\parallel\delta\text{, }\gamma\neq\delta \)}{Planes \( \alpha \) and \( \beta \) are intersecting.}{\( \alpha = \delta \)}{}}
\LANGpl{
\Question{
Płaszczyzny  \( \alpha \), \( \beta \), \( \gamma \) i \( \delta \) są określone równaniami ogólnymi:  
\[ \begin{aligned}
&\alpha\colon \frac23x-4y+6z-\frac83=0; \\
&\beta\colon x-2y+3z-4=0; \\
&\gamma\colon 2x-12y+18z-4 =0; \\
&\delta\colon x-6y+9z-4 =0.
\end{aligned} \]
Z poniższych stwierdzeń, wybierz t nieprawdziwe.}
{\( \alpha \parallel\delta\text{, }\alpha\neq\delta \)}{Płaszczyzny  \( \beta \) i \( \delta \) przecinają się.}{\( \gamma\parallel\delta\text{, }\gamma\neq\delta \)}{Płaszczyzny \( \alpha \) i \( \beta \) przecinają się.}{\( \alpha = \delta \)}{}}
\LANGcs{
\Question{
Roviny \( \alpha \), \( \beta \), \( \gamma \) a \( \delta \) jsou dány svými obecnými rovnicemi:  
\[ \begin{aligned}
&\alpha\colon \frac23x-4y+6z-\frac83=0; \\
&\beta\colon x-2y+3z-4=0; \\
&\gamma\colon 2x-12y+18z-4 =0; \\
&\delta\colon x-6y+9z-4 =0.
\end{aligned} \]
Určete, které z následujících tvrzení není pravdivé.}
{\( \alpha \parallel\delta\text{, }\alpha\neq\delta \)}{Roviny \( \beta \) a \( \delta \) jsou různoběžné.}{\( \gamma\parallel\delta\text{, }\gamma\neq\delta \)}{Roviny \( \alpha \) a \( \beta \) jsou různoběžné.}{\( \alpha = \delta \)}{}}
\LANGsk{
\Question{
Roviny \( \alpha \), \( \beta \), \( \gamma \) a \( \delta \) sú dané svojimi všeobecnými rovnicami:  
\[ \begin{aligned}
&\alpha\colon \frac23x-4y+6z-\frac83=0; \\
&\beta\colon x-2y+3z-4=0; \\
&\gamma\colon 2x-12y+18z-4 =0; \\
&\delta\colon x-6y+9z-4 =0.
\end{aligned} \]
Určte, ktoré z nasledujúcich tvrdení nie je pravdivé.}
{\( \alpha \parallel\delta\text{, }\alpha\neq\delta \)}{Roviny \( \beta \) a \( \delta \) sú rôznobežné.}{\( \gamma\parallel\delta\text{, }\gamma\neq\delta \)}{Roviny \( \alpha \) a \( \beta \) sú rôznobežné.}{\( \alpha = \delta \)}{}}
\LANGes{
\Question{
Dados los planos \( \alpha \), \( \beta \), \( \gamma \) y \( \delta \) cuyas ecuaciones generales son:   
\[ \begin{aligned}
&\alpha\colon \frac23x-4y+6z-\frac83=0; \\
&\beta\colon x-2y+3z-4=0; \\
&\gamma\colon 2x-12y+18z-4 =0; \\
&\delta\colon x-6y+9z-4 =0.
\end{aligned} \]
Determina cuál de las siguientes afirmaciones no es verdadera:}
{\( \alpha \parallel\delta\text{, }\alpha\neq\delta \)}{Los planos \( \beta \) y \( \delta \) no son paralelos.}{\( \gamma\parallel\delta\text{, }\gamma\neq\delta \)}{Los planos \( \alpha \) y \( \beta \) no son paralelos.}{\( \alpha = \delta \)}{}}
\End

\Data\ID{1003188907}
\Updated{2020-07-15 03:07:21}
\Level{A}
\SubArea{Analytic geometry in a space}
\LANGen{
\Question{
We are given two intersecting planes \( x-6y+9z-4=0 \) and \( x-2y+3z-4=0 \). Find the parametric equations of their line of intersection \( p \).}
{\( \begin{aligned}
p\colon x&=4, \\
y&=\phantom{4+}\ 3t, \\
z&=\phantom{4+}\ 2t;\ t\in\mathbb{R}
\end{aligned} \)}{\( \begin{aligned}
p\colon x&=4+t , \\
y&=\phantom{4+}\ 3t , \\
z&=\phantom{4+}\ 2t;\ t\in\mathbb{R}
\end{aligned} \)}{\( \begin{aligned}
p\colon x&=4, \\
y&=\frac32+3t, \\
z&=1+2t;\ t\in\mathbb{R}
\end{aligned} \)}{\( \begin{aligned}
p\colon x&=4+t, \\
y&=\frac32+3t, \\
z&=1+2t;\ t\in\mathbb{R}
\end{aligned} \)}{}{}}
\LANGpl{
\Question{
Dane są dwie przecinające się płaszczyzny \( x-6y+9z-4=0 \) and \( x-2y+3z-4=0 \). Wskaż równanie parametryczne prostej ich przecięcia \( p \).}
{\( \begin{aligned}
p\colon x&=4, \\
y&=\phantom{4+}\ 3t, \\
z&=\phantom{4+}\ 2t;\ t\in\mathbb{R}
\end{aligned} \)}{\( \begin{aligned}
p\colon x&=4+t, \\
y&=\phantom{4+}\ 3t, \\
z&=\phantom{4+}\ 2t;\ t\in\mathbb{R}
\end{aligned} \)}{\( \begin{aligned}
p\colon x&=4, \\
y&=\frac32+3t, \\
z&=1+2t;\ t\in\mathbb{R}
\end{aligned} \)}{\( \begin{aligned}
p\colon x&=4+t, \\
y&=\frac32+3t, \\
z&=1+2t;\ t\in\mathbb{R}
\end{aligned} \)}{}{}}
\LANGcs{
\Question{
Jsou dány různoběžné roviny  \( x-6y+9z-4=0 \)  a  \( x-2y+3z-4=0 \). Určete parametrické rovnice jejich průsečnice \( p \).}
{\( \begin{aligned}
p\colon x&=4, \\
y&=\phantom{4+}\ 3t, \\
z&=\phantom{4+}\ 2t;\ t\in\mathbb{R}
\end{aligned} \)}{\( \begin{aligned}
p\colon x&=4+t, \\
y&=\phantom{4+}\ 3t, \\
z&=\phantom{4+}\ 2t;\ t\in\mathbb{R}
\end{aligned} \)}{\( \begin{aligned}
p\colon x&=4, \\
y&=\frac32+3t , \\
z&=1+2t;\ t\in\mathbb{R}
\end{aligned} \)}{\( \begin{aligned}
p\colon x&=4+t, \\
y&=\frac32+3t, \\
z&=1+2t;\ t\in\mathbb{R}
\end{aligned} \)}{}{}}
\LANGsk{
\Question{
Sú dané rôznobežné roviny  \( x-6y+9z-4=0 \)  a  \( x-2y+3z-4=0 \). Určte parametrické rovnice ich priesečnice \( p \).}
{\( \begin{aligned}
p\colon x&=4, \\
y&=\phantom{4+}\ 3t, \\
z&=\phantom{4+}\ 2t;\ t\in\mathbb{R}
\end{aligned} \)}{\( \begin{aligned}
p\colon x&=4+t, \\
y&=\phantom{4+}\ 3t, \\
z&=\phantom{4+}\ 2t;\ t\in\mathbb{R}
\end{aligned} \)}{\( \begin{aligned}
p\colon x&=4, \\
y&=\frac32+3t , \\
z&=1+2t;\ t\in\mathbb{R}
\end{aligned} \)}{\( \begin{aligned}
p\colon x&=4+t , \\
y&=\frac32+3t , \\
z&=1+2t;\ t\in\mathbb{R}
\end{aligned} \)}{}{}}
\LANGes{
\Question{
Dados los planos paralelos \( x-6y+9z-4=0 \) y \( x-2y+3z-4=0 \). Determina las ecuaciones paramétricas de su línea de intersecciones \( p \).}
{\( \begin{aligned}
p\colon x&=4, \\
y&=\phantom{4+}\ 3t, \\
z&=\phantom{4+}\ 2t;\ t\in\mathbb{R}
\end{aligned} \)}{\( \begin{aligned}
p\colon x&=4+t , \\
y&=\phantom{4+}\ 3t , \\
z&=\phantom{4+}\ 2t;\ t\in\mathbb{R}
\end{aligned} \)}{\( \begin{aligned}
p\colon x&=4, \\
y&=\frac32+3t, \\
z&=1+2t;\ t\in\mathbb{R}
\end{aligned} \)}{\( \begin{aligned}
p\colon x&=4+t, \\
y&=\frac32+3t, \\
z&=1+2t;\ t\in\mathbb{R}
\end{aligned} \)}{}{}}
\End

\Data\ID{1003188908}
\Updated{2020-07-15 03:07:21}
\Level{A}
\SubArea{Analytic geometry in a space}
\LANGen{
\Question{
Find the real number \( k \) so that the line \( AB \), where \( A=[6;k;5] \) and \( B=[3;-1;2] \), is parallel to the plane \( 2x-3y+9z-4=0 \).}
{\( k=10 \)}{\( k=12 \)}{\( k=8 \)}{\( k=9 \)}{}{}}
\LANGpl{
\Question{
Wskaż liczbę rzeczywistą  \( k \) tak, aby prosta \( AB \), gdzie \( A=[6;k;5] \) i \( B=[3;-1;2] \), była równoległa do płaszczyzny \( 2x-3y+9z-4=0 \).}
{\( k=10 \)}{\( k=12 \)}{\( k=8 \)}{\( k=9 \)}{}{}}
\LANGcs{
\Question{
Určete reálné číslo \( k \) tak, aby přímka \( AB \), kde \( A=[6;k;5] \) a \( B=[3;-1;2] \), byla rovnoběžná s rovinou \( 2x-3y+9z-4=0 \).}
{\( k=10 \)}{\( k=12 \)}{\( k=8 \)}{\( k=9 \)}{}{}}
\LANGsk{
\Question{
Určte reálne číslo \( k \) tak, aby priamka \( AB \), kde \( A=[6;k;5] \) a \( B=[3;-1;2] \), bola rovnobežná s rovinou \( 2x-3y+9z-4=0 \).}
{\( k=10 \)}{\( k=12 \)}{\( k=8 \)}{\( k=9 \)}{}{}}
\LANGes{
\Question{
Determina el número real \( k \) para que la recta \( AB \), donde \( A=[6;k;5] \) y \( B=[3;-1;2] \), sea paralela con el plano \( 2x-3y+9z-4=0 \).}
{\( k=10 \)}{\( k=12 \)}{\( k=8 \)}{\( k=9 \)}{}{}}
\End

\Data\ID{1103188705}
\Updated{2020-07-15 03:07:21}
\Level{A}
\SubArea{Analytic geometry in a space}
\IMG{1103188705.pdf}
\LANGen{
\Question{
Find parametric equations of the line \( p \) that passes through the point \( K=[4;2;3] \), is parallel to the \( xy \)-coordinate plane, and is intersecting the \( z \)-axis.}
{$\begin{aligned}
p\colon x&=4+2t, \\
y&=2+t, \\
z&=3;\ t\in\mathbb{R}
\end{aligned}$}{$\begin{aligned}
p\colon x&=4+2t, \\
y&=2+t, \\
z&=3+t;\ t\in\mathbb{R}
\end{aligned}$}{$\begin{aligned}
p\colon x&=4, \\
y&=2, \\
z&=3+3t;\ t\in\mathbb{R}
\end{aligned}$}{$\begin{aligned}
p\colon x&=4-2t, \\
y&=2-4t, \\
z&=3t;\ t\in\mathbb{R}
\end{aligned}$}{}{}}
\LANGpl{
\Question{
Wskaż równanie parametryczne prostej \( p \) przechodzącej przez punkt \( K=[4;2;3] \), równoległej do płaszczyzny w układzie współrzędnych  \( xy \), oraz przecinającą oś \( z \).}
{$\begin{aligned}
p\colon x&=4+2t, \\
y&=2+t, \\
z&=3;\ t\in\mathbb{R}
\end{aligned}$}{$\begin{aligned}
p\colon x&=4+2t, \\
y&=2+t, \\
z&=3+t;\ t\in\mathbb{R}
\end{aligned}$}{$\begin{aligned}
p\colon x&=4, \\
y&=2, \\
z&=3+3t;\ t\in\mathbb{R}
\end{aligned}$}{$\begin{aligned}
p\colon x&=4-2t, \\
y&=2-4t, \\
z&=3t;\ t\in\mathbb{R}
\end{aligned}$}{}{}}
\LANGcs{
\Question{
Určete parametrické rovnice přímky \( p \), která prochází bodem \( K=[4;2;3] \), je rovnoběžná se souřadnou rovinou \( (xy) \) a je různoběžná s osou \( z \).}
{$\begin{aligned}
p\colon x&=4+2t, \\
y&=2+t, \\
z&=3;\ t\in\mathbb{R}
\end{aligned}$}{$\begin{aligned}
p\colon x&=4+2t, \\
y&=2+t, \\
z&=3+t;\ t\in\mathbb{R}
\end{aligned}$}{$\begin{aligned}
p\colon x&=4, \\
y&=2, \\
z&=3+3t;\ t\in\mathbb{R}
\end{aligned}$}{$\begin{aligned}
p\colon x&=4-2t, \\
y&=2-4t, \\
z&=3t;\ t\in\mathbb{R}
\end{aligned}$}{}{}}
\LANGsk{
\Question{
Určte parametrické rovnice priamky \( p \), ktorá prechádza bodom \( K=[4;2;3] \), je rovnobežná so súradnicovou rovinou \( (xy) \) a je rôznobežná s osou \( z \).}
{$\begin{aligned}
p\colon x&=4+2t, \\
y&=2+t, \\
z&=3;\ t\in\mathbb{R}
\end{aligned}$}{$\begin{aligned}
p\colon x&=4+2t, \\
y&=2+t, \\
z&=3+t;\ t\in\mathbb{R}
\end{aligned}$}{$\begin{aligned}
p\colon x&=4, \\
y&=2, \\
z&=3+3t;\ t\in\mathbb{R}
\end{aligned}$}{$\begin{aligned}
p\colon x&=4-2t, \\
y&=2-4t, \\
z&=3t;\ t\in\mathbb{R}
\end{aligned}$}{}{}}
\LANGes{
\Question{
Determina las ecuaciones paramátricas de la recta \( p \) la cual pasa por el punto \( K=[4;2;3] \), es paralela con la coordenada \( xy \) y no es paralela con el eje \( z \).}
{$\begin{aligned}
p\colon x&=4+2t, \\
y&=2+t, \\
z&=3;\ t\in\mathbb{R}
\end{aligned}$}{$\begin{aligned}
p\colon x&=4+2t, \\
y&=2+t, \\
z&=3+t;\ t\in\mathbb{R}
\end{aligned}$}{$\begin{aligned}
p\colon x&=4, \\
y&=2, \\
z&=3+3t;\ t\in\mathbb{R}
\end{aligned}$}{$\begin{aligned}
p\colon x&=4-2t, \\
y&=2-4t, \\
z&=3t;\ t\in\mathbb{R}
\end{aligned}$}{}{}}
\End

\Data\ID{1103188706}
\Updated{2020-07-15 03:07:21}
\Level{A}
\SubArea{Analytic geometry in a space}
\IMG{1103188706.pdf}
\LANGen{
\Question{
We are given points \( A=[2;4;0] \) and \( B=[4;7;6] \). Find parametric equations of a line \( q \), which is the orthogonal projection of the line \( AB \) into the coordinate plane \( xy \).}
{$\begin{aligned}
p\colon x&=4+2t, \\
y&=7+3t, \\
z&=0;\ t\in\mathbb{R}
\end{aligned}$}{$\begin{aligned}
p\colon x&=2+4t, \\
y&=4+7t, \\
z&=6t;\ t\in\mathbb{R}
\end{aligned}$}{$\begin{aligned}
p\colon x&=4+2t, \\
y&=7+3t, \\
z&=6;\ t\in\mathbb{R}
\end{aligned}$}{$\begin{aligned}
p\colon x&=2-2t, \\
y&=4-3t, \\
z&=-6t;\ t\in\mathbb{R}
\end{aligned}$}{}{}}
\LANGpl{
\Question{
Dane są punkty \( A=[2;4;0] \) and \( B=[4;7;6] \). Wskaż równanie parametryczne prostej \( q \), która jest rzutem prostokątnym prostej \( AB \) na płaszczyźnie w układzie współrzędnych \( xy \).}
{$\begin{aligned}
p\colon x&=4+2t, \\
y&=7+3t, \\
z&=0;\ t\in\mathbb{R}
\end{aligned}$}{$\begin{aligned}
p\colon x&=2+4t, \\
y&=4+7t, \\
z&=6t;\ t\in\mathbb{R}
\end{aligned}$}{$\begin{aligned}
p\colon x&=4+2t, \\
y&=7+3t, \\
z&=6;\ t\in\mathbb{R}
\end{aligned}$}{$\begin{aligned}
p\colon x&=2-2t, \\
y&=4-3t, \\
z&=-6t;\ t\in\mathbb{R}
\end{aligned}$}{}{}}
\LANGcs{
\Question{
Jsou dány body \( A=[2;4;0] \) a \( B=[4;7;6] \). Určete parametrické rovnice přímky \( q \), která je pravoúhlým průmětem přímky \( AB \) do souřadné roviny \( (xy) \).}
{$\begin{aligned}
p\colon x&=4+2t, \\
y&=7+3t, \\
z&=0;\ t\in\mathbb{R}
\end{aligned}$}{$\begin{aligned}
p\colon x&=2+4t, \\
y&=4+7t, \\
z&=6t;\ t\in\mathbb{R}
\end{aligned}$}{$\begin{aligned}
p\colon x&=4+2t, \\
y&=7+3t, \\
z&=6;\ t\in\mathbb{R}
\end{aligned}$}{$\begin{aligned}
p\colon x&=2-2t, \\
y&=4-3t, \\
z&=-6t;\ t\in\mathbb{R}
\end{aligned}$}{}{}}
\LANGsk{
\Question{
Sú dané body \( A=[2;4;0] \) a \( B=[4;7;6] \). Určte parametrické rovnice priamky \( q \), ktorá je pravouhlým priemetom priamky \( AB \) do súradnicovej roviny \( (xy) \).}
{$\begin{aligned}
p\colon x&=4+2t, \\
y&=7+3t, \\
z&=0;\ t\in\mathbb{R}
\end{aligned}$}{$\begin{aligned}
p\colon x&=2+4t, \\
y&=4+7t, \\
z&=6t;\ t\in\mathbb{R}
\end{aligned}$}{$\begin{aligned}
p\colon x&=4+2t, \\
y&=7+3t, \\
z&=6;\ t\in\mathbb{R}
\end{aligned}$}{$\begin{aligned}
p\colon x&=2-2t, \\
y&=4-3t, \\
z&=-6t;\ t\in\mathbb{R}
\end{aligned}$}{}{}}
\LANGes{
\Question{
Dados los puntos \( A=[2;4;0] \) y \( B=[4;7;6] \). Determina las ecuaciones paramétricas de la recta \( q \), que es la proyección del segmento \( AB \) en las coordenadas \( xy \).}
{$\begin{aligned}
p\colon x&=4+2t, \\
y&=7+3t, \\
z&=0;\ t\in\mathbb{R}
\end{aligned}$}{$\begin{aligned}
p\colon x&=2+4t, \\
y&=4+7t, \\
z&=6t;\ t\in\mathbb{R}
\end{aligned}$}{$\begin{aligned}
p\colon x&=4+2t, \\
y&=7+3t, \\
z&=6;\ t\in\mathbb{R}
\end{aligned}$}{$\begin{aligned}
p\colon x&=2-2t, \\
y&=4-3t, \\
z&=-6t;\ t\in\mathbb{R}
\end{aligned}$}{}{}}
\End

\Data\ID{1103188804}
\Updated{2020-07-15 03:07:21}
\Level{A}
\SubArea{Analytic geometry in a space}
\IMG{1103188804.pdf}
\LANGen{
\Question{
Find the standard equation of the plane shown in the picture:}
{\( y+2z-4=0 \)}{\( x+y+2z-4=0 \)}{\( 2y+4z=-8 \)}{\( x+2y+4z-8=0 \)}{}{}}
\LANGpl{
\Question{
Wskaż równanie płaszczyzny znajdującej się na rysunku:}
{\( y+2z-4=0 \)}{\( x+y+2z-4=0 \)}{\( 2y+4z=-8 \)}{\( x+2y+4z-8=0 \)}{}{}}
\LANGcs{
\Question{
Určete obecnou rovnici roviny znázorněné na obrázku:}
{\( y+2z-4=0 \)}{\( x+y+2z-4=0 \)}{\( 2y+4z=-8 \)}{\( x+2y+4z-8=0 \)}{}{}}
\LANGsk{
\Question{
Určte všeobecnú rovnicu roviny znázornenej na obrázku:}
{\( y+2z-4=0 \)}{\( x+y+2z-4=0 \)}{\( 2y+4z=-8 \)}{\( x+2y+4z-8=0 \)}{}{}}
\LANGes{
\Question{
Busca la ecuación general del plano presentado en la imagen:}
{\( y+2z-4=0 \)}{\( x+y+2z-4=0 \)}{\( 2y+4z=-8 \)}{\( x+2y+4z-8=0 \)}{}{}}
\End

\Data\ID{1103188805}
\Updated{2020-07-15 03:07:21}
\Level{A}
\SubArea{Analytic geometry in a space}
\IMG{1103188805.pdf}
\LANGen{
\Question{
Find the standard equation of the plane shown in the picture:}
{\( 3x+5z-15=0 \)}{\( 3x-5z+15=0 \)}{\( 3x+15y+5z-15=0 \)}{\( 3x-15y+5z+15=0 \)}{}{}}
\LANGpl{
\Question{
Wskaż równanie płaszczyzny znajdującej się na rysunku:}
{\( 3x+5z-15=0 \)}{\( 3x-5z+15=0 \)}{\( 3x+15y+5z-15=0 \)}{\( 3x-15y+5z+15=0 \)}{}{}}
\LANGcs{
\Question{
Určete obecnou rovnici roviny znázorněné na obrázku:}
{\( 3x+5z-15=0 \)}{\( 3x-5z+15=0 \)}{\( 3x+15y+5z-15=0 \)}{\( 3x-15y+5z+15=0 \)}{}{}}
\LANGsk{
\Question{
Určte všeobecnú rovnicu roviny znázornenej na obrázku:}
{\( 3x+5z-15=0 \)}{\( 3x-5z+15=0 \)}{\( 3x+15y+5z-15=0 \)}{\( 3x-15y+5z+15=0 \)}{}{}}
\LANGes{
\Question{
Halla la ecuación general del plano presentado en la imagen:}
{\( 3x+5z-15=0 \)}{\( 3x-5z+15=0 \)}{\( 3x+15y+5z-15=0 \)}{\( 3x-15y+5z+15=0 \)}{}{}}
\End

\Data\ID{1103188806}
\Updated{2020-07-15 03:07:21}
\Level{A}
\SubArea{Analytic geometry in a space}
\IMG{1103188806.pdf}
\LANGen{
\Question{
Find the standard equation of the plane shown in the picture:}
{\( 3x+2y-12=0 \)}{\( 2x+3y-26=0 \)}{\( 2x+3y-z-12=0 \)}{\( 3x+2y+z-12=0 \)}{}{}}
\LANGpl{
\Question{
Wskaż równanie płaszczyzny znajdującej się na rysunku:}
{\( 3x+2y-12=0 \)}{\( 2x+3y-26=0 \)}{\( 2x+3y-z-12=0 \)}{\( 3x+2y+z-12=0 \)}{}{}}
\LANGcs{
\Question{
Určete obecnou rovnici roviny znázorněné na obrázku:}
{\( 3x+2y-12=0 \)}{\( 2x+3y-26=0 \)}{\( 2x+3y-z-12=0 \)}{\( 3x+2y+z-12=0 \)}{}{}}
\LANGsk{
\Question{
Určte všeobecnú rovnicu roviny znázornenej na obrázku:}
{\( 3x+2y-12=0 \)}{\( 2x+3y-26=0 \)}{\( 2x+3y-z-12=0 \)}{\( 3x+2y+z-12=0 \)}{}{}}
\LANGes{
\Question{
Halla la ecuación general del plano presentado en la imagen:}
{\( 3x+2y-12=0 \)}{\( 2x+3y-26=0 \)}{\( 2x+3y-z-12=0 \)}{\( 3x+2y+z-12=0 \)}{}{}}
\End

\Data\ID{1103188901}
\Updated{2020-07-15 03:07:21}
\Level{A}
\SubArea{Analytic geometry in a space}
\IMG{1103188901.pdf}
\LANGen{
\Question{
Find the general equation of the plane shown in the picture.}
{\( 15x+10y+12z-60=0 \)}{\( 4x+6y+5z-60=0 \)}{\( 10x+12y+15z-60=0 \)}{\( 30x+20y+24z-60=0 \)}{}{}}
\LANGpl{
\Question{
Wskaż równanie płaszczyzny znajdującej się na rysunku.}
{\( 15x+10y+12z-60=0 \)}{\( 4x+6y+5z-60=0 \)}{\( 10x+12y+15z-60=0 \)}{\( 30x+20y+24z-60=0 \)}{}{}}
\LANGcs{
\Question{
Určete obecnou rovnici roviny znázorněné na obrázku.}
{\( 15x+10y+12z-60=0 \)}{\( 4x+6y+5z-60=0 \)}{\( 10x+12y+15z-60=0 \)}{\( 30x+20y+24z-60=0 \)}{}{}}
\LANGsk{
\Question{
Určte všeobecnú rovnicu roviny znázornenej na obrázku:}
{\( 15x+10y+12z-60=0 \)}{\( 4x+6y+5z-60=0 \)}{\( 10x+12y+15z-60=0 \)}{\( 30x+20y+24z-60=0 \)}{}{}}
\LANGes{
\Question{
Halla la ecuación general del plano presentado en la imagen:}
{\( 15x+10y+12z-60=0 \)}{\( 4x+6y+5z-60=0 \)}{\( 10x+12y+15z-60=0 \)}{\( 30x+20y+24z-60=0 \)}{}{}}
\End

\Data\ID{1103188902}
\Updated{2020-07-15 03:07:21}
\Level{A}
\SubArea{Analytic geometry in a space}
\IMG{1103188902_0.pdf}
\LANGen{
\Question{
Assign the planes shown in the picture to the corresponding general equations.}
{\( \alpha\colon y-2=0;\ \beta\colon z-2=0;\ \gamma\colon x-2=0 \)}{\( \alpha\colon y+2=0;\ \beta\colon z+2=0;\ \gamma\colon x+2=0 \)}{\( \alpha\colon x+z-2=0;\ \beta\colon x+y-2=0;\ \gamma\colon y+z-2=0 \)}{\( \alpha\colon x-y+z-2=0;\ \beta\colon x+y-z-2=0;\ \gamma\colon -x+y+z-2=0 \)}{}{}}
\LANGpl{
\Question{
Przyporządkuj płaszczyzny przedstawione na rysunku do odpowiednich równań ogólnych.}
{\( \alpha\colon y-2=0;\ \beta\colon z-2=0;\ \gamma\colon x-2=0 \)}{\( \alpha\colon y+2=0;\ \beta\colon z+2=0;\ \gamma\colon x+2=0 \)}{\( \alpha\colon x+z-2=0;\ \beta\colon x+y-2=0;\ \gamma\colon y+z-2=0 \)}{\( \alpha\colon x-y+z-2=0;\ \beta\colon x+y-z-2=0;\ \gamma\colon -x+y+z-2=0 \)}{}{}}
\LANGcs{
\Question{
Rovinám znázorněným na obrázcích přiřaďte jejich obecné rovnice.}
{\( \alpha\colon y-2=0;\ \beta\colon z-2=0;\ \gamma\colon x-2=0 \)}{\( \alpha\colon y+2=0;\ \beta\colon z+2=0;\ \gamma\colon x+2=0 \)}{\( \alpha\colon x+z-2=0;\ \beta\colon x+y-2=0;\ \gamma\colon y+z-2=0 \)}{\( \alpha\colon x-y+z-2=0;\ \beta\colon x+y-z-2=0;\ \gamma\colon -x+y+z-2=0 \)}{}{}}
\LANGsk{
\Question{
Rovinám znázorneným na obrázkoch priraďte ich všeobecné rovnice.}
{\( \alpha\colon y-2=0;\ \beta\colon z-2=0;\ \gamma\colon x-2=0 \)}{\( \alpha\colon y+2=0;\ \beta\colon z+2=0;\ \gamma\colon x+2=0 \)}{\( \alpha\colon x+z-2=0;\ \beta\colon x+y-2=0;\ \gamma\colon y+z-2=0 \)}{\( \alpha\colon x-y+z-2=0;\ \beta\colon x+y-z-2=0;\ \gamma\colon -x+y+z-2=0 \)}{}{}}
\LANGes{
\Question{
Asigna las ecuaciones generales a sus planos presentados en las imágenes.}
{\( \alpha\colon y-2=0;\ \beta\colon z-2=0;\ \gamma\colon x-2=0 \)}{\( \alpha\colon y+2=0;\ \beta\colon z+2=0;\ \gamma\colon x+2=0 \)}{\( \alpha\colon x+z-2=0;\ \beta\colon x+y-2=0;\ \gamma\colon y+z-2=0 \)}{\( \alpha\colon x-y+z-2=0;\ \beta\colon x+y-z-2=0;\ \gamma\colon -x+y+z-2=0 \)}{}{}}
\End

\Data\ID{9000101001}
\Updated{2020-07-15 03:07:37}
\Level{A}
\SubArea{Analytic geometry in a space}
\LANGen{
\Question{
Determine whether two lines $p$ and $q$ are identical, parallel, intersecting or skew.
\[\begin{aligned}
 p\colon x & = 1 + t, & &
 \\y & = 2 - t, & &
 \\z & = 1 - t;\ t\in \mathbb{R} & &
\end{aligned}\] \[\begin{aligned}
 q\colon x & = 2s, & &
 \\y & = -1, & &
 \\z & = 2 - 2s;\ s\in \mathbb{R} & &
\end{aligned}\]}
{intersecting lines}{skew lines}{identical lines}{parallel lines, not identical}{}{}}
\LANGpl{
\Question{
Określ wzajemne położenie prostych w przestrzeni. 
\[\begin{aligned}
 p\colon x & = 1 + t, & &
 \\y & = 2 - t, & &
 \\z & = 1 - t;\ t\in \mathbb{R} & &
\end{aligned}\] \[\begin{aligned}
 q\colon x & = 2s, & &
 \\y & = -1, & &
 \\z & = 2 - 2s;\ s\in \mathbb{R} & &
\end{aligned}\]}
{Proste przecinające się.}{Proste skośne.}{Proste pokrywające się.}{Proste równoległe, nie pokrywające się.}{}{}}
\LANGcs{
\Question{
Určete vzájemnou polohu přímek
\(p\),
\(q\),
kde:
\[\begin{aligned}
 p\colon x & = 1 + t, & &
 \\y & = 2 - t, & &
 \\z & = 1 - t;\ t\in \mathbb{R} & &
\end{aligned}\] \[\begin{aligned}
 q\colon x & = 2s, & &
 \\y & = -1, & &
 \\z & = 2 - 2s;\ s\in \mathbb{R} & &
\end{aligned}\]}
{Dané přímky jsou různoběžné.}{Dané přímky jsou mimoběžné.}{Dané přímky jsou totožné.}{Dané přímky jsou rovnoběžné různé.}{}{}}
\LANGsk{
\Question{
Určte vzájomnú polohu priamok
\(p\),
\(q\),
kde:
\[\begin{aligned}
 p\colon x & = 1 + t, & &
 \\y & = 2 - t, & &
 \\z & = 1 - t;\ t\in \mathbb{R} & &
\end{aligned}\] \[\begin{aligned}
 q\colon x & = 2s, & &
 \\y & = -1, & &
 \\z & = 2 - 2s;\ s\in \mathbb{R} & &
\end{aligned}\]}
{Dané priamky sú rôznobežné.}{Dané priamky sú mimobežné.}{Dané priamky sú totožné.}{Dané priamky sú rovnobežné rôzne.}{}{}}
\LANGes{
\Question{
Determina la posición de dos rectas $p$ y $q$ si se sabe que:
\[\begin{aligned}
 p\colon x & = 1 + t, & &
 \\y & = 2 - t, & &
 \\z & = 1 - t;\ t\in \mathbb{R} & &
\end{aligned}\] \[\begin{aligned}
 q\colon x & = 2s, & &
 \\y & = -1, & &
 \\z & = 2 - 2s;\ s\in \mathbb{R} & &
\end{aligned}\]}
{Las rectas dadas son rectas no paralelas.}{Las rectas dadas son rectas sesgadas.}{Las rectas dadas son idénticas.}{Las rectas dadas son paralelas diferentes (no idénticas).}{}{}}
\End

\Data\ID{9000101002}
\Updated{2020-07-15 03:07:37}
\Level{A}
\SubArea{Analytic geometry in a space}
\LANGen{
\Question{
Find the intersection of the line \(AB\) 
and the line \(p\),
where \(A = [0;1;2]\),
\(B = [4;1;-2]\) and
\[ 
\begin{aligned}p\colon x& = 1 + t, &
\\y & = 2 - t,
\\z & = 1 - t;\ t\in \mathbb{R}.
\\ \end{aligned}
\]}
{\([2;1;0]\)}{\([1;2;1]\)}{\([3;0;-1]\)}{There is no intersection.}{}{}}
\LANGpl{
\Question{
Wyznacz punkt wspólny prostej  \(AB\) 
i prostej \(p\),
 \(A = [0;1;2]\),
\(B = [4;1;-2]\) i
\[ 
\begin{aligned}p\colon x& = 1 + t, &
\\y & = 2 - t,
\\z & = 1 - t;\ t\in \mathbb{R}.
\\ \end{aligned}
\]}
{\([2;1;0]\)}{\([1;2;1]\)}{\([3;0;-1]\)}{Brak punktów wspólnych.}{}{}}
\LANGcs{
\Question{
Jsou dány body \(A = [0;1;2]\),
\(B = [4;1;-2]\) a přímka
\(p\colon x = 1 + t;\: y = 2 - t;\: z = 1 - t,\: t\in \mathbb{R}\). Určete průsečík
přímky \(AB\) a
přímky \(p\),
případně zaškrtněte, že neexistuje.}
{\([2;1;0]\)}{\([1;2;1]\)}{\([3;0;-1]\)}{Průsečík daných přímek neexistuje.}{}{}}
\LANGsk{
\Question{
Sú dané body \(A = [0;1;2]\),
\(B = [4;1;-2]\) a priamka
\(p\colon x = 1 + t;\: y = 2 - t;\: z = 1 - t,\: t\in \mathbb{R}\). Určte priesečník
priamky \(AB\) a
priamky \(p\),
prípadne zaškrtnite, že neexistuje.}
{\([2;1;0]\)}{\([1;2;1]\)}{\([3;0;-1]\)}{Priesečník daných priamok neexistuje.}{}{}}
\LANGes{
\Question{
Dados los puntos \(A = [0;1;2]\) y
\(B = [4;1;-2]\) y la recta \(p\). 
\[ 
\begin{aligned}p\colon x& = 1 + t, &
\\y & = 2 - t,
\\z & = 1 - t;\ t\in \mathbb{R}.
\\ \end{aligned}
\]
Determina la intersección de la recta \(AB\) y la recta \(p\).
En el caso de que no exista marca la opción de que no existe.}
{\([2;1;0]\)}{\([1;2;1]\)}{\([3;0;-1]\)}{No hay intersección de las rectas.}{}{}}
\End

\Data\ID{9000101003}
\Updated{2020-07-15 03:07:37}
\Level{A}
\SubArea{Analytic geometry in a space}
\LANGen{
\Question{
Find the value of the real parameter \(m\in \mathbb{R}\) 
which ensures that the lines \(p\) 
and \(q\) 
are parallel and not identical. 
\[ 
\begin{aligned}p\colon x& = 1 + t, &
\\y & = 2 - t,
\\z & = 1 - t;\ t\in \mathbb{R}
\\ \end{aligned}\qquad \qquad \begin{aligned}q\colon x& = s, &
\\y & = -s,
\\z & = 3 + ms;\  s\in \mathbb{R}.
\\ \end{aligned}
\]}
{\(m = -1\)}{\(m = -2\)}{\(m = 0\)}{\(m = 1\)}{}{}}
\LANGpl{
\Question{
Wyznacz wartość rzeczywistą parametru \(m\in \mathbb{R}\) 
tak, aby proste \(p\) 
i \(q\) 
były równoległe i nie pokrywające się.
\[ 
\begin{aligned}p\colon x& = 1 + t, &
\\y & = 2 - t,
\\z & = 1 - t;\ t\in \mathbb{R}
\\ \end{aligned}\qquad \qquad \begin{aligned}q\colon x& = s, &
\\y & = -s,
\\z & = 3 + ms;\  s\in \mathbb{R}.
\\ \end{aligned}
\]}
{\(m = -1\)}{\(m = -2\)}{\(m = 0\)}{\(m = 1\)}{}{}}
\LANGcs{
\Question{
Určete hodnotu reálného parametru
\(m\) tak, aby
přímky \(p\colon x = 1 + t;\: y = 2 - t;\: z = 1 - t,\: t\in \mathbb{R}\) 
a \(q\colon x = s;\: y = -s;\: z = 3 + ms,\: s\in \mathbb{R}\) 
byly rovnoběžné různé.}
{\(m = -1\)}{\(m = -2\)}{\(m = 0\)}{\(m = 1\)}{}{}}
\LANGsk{
\Question{
Určte hodnotu reálneho parametra
\(m\) tak, aby
priamky \(p\colon x = 1 + t;\: y = 2 - t;\: z = 1 - t,\: t\in \mathbb{R}\) 
a \(q\colon x = s;\: y = -s;\: z = 3 + ms,\: s\in \mathbb{R}\) 
boli rovnobežné rôzne.}
{\(m = -1\)}{\(m = -2\)}{\(m = 0\)}{\(m = 1\)}{}{}}
\LANGes{
\Question{
Halla el valor del parámetro real \(m\in \mathbb{R}\) para que las rectas \(p\) y \(q\) sean paralelas y no indénticas.
\[ 
\begin{aligned}p\colon x& = 1 + t, &
\\y & = 2 - t,
\\z & = 1 - t;\ t\in \mathbb{R}
\\ \end{aligned}\qquad \qquad \begin{aligned}q\colon x& = s, &
\\y & = -s,
\\z & = 3 + ms;\  s\in \mathbb{R}.
\\ \end{aligned}
\]}
{\(m = -1\)}{\(m = -2\)}{\(m = 0\)}{\(m = 1\)}{}{}}
\End

\Data\ID{9000101004}
\Updated{2020-09-01 09:09:23}
\Level{A}
\SubArea{Analytic geometry in a space}
\LANGen{
\Question{
Find all the values of the real parameter  \(m\) 
so that  the lines \(p\) 
and \(q\) 
are skew lines. 
\[ 
\begin{aligned}p\colon x& = 1 + t, &
\\y & = 2 - t,
\\z & = 1 - t;\ t\in \mathbb{R}
\\ \end{aligned}\qquad \qquad \begin{aligned}q\colon x& = s, &
\\y & = 1 + s,
\\z & = 3 + ms;\ s\in \mathbb{R}
\\ \end{aligned}
\]}
{\(m\in\mathbb{R}\setminus\{-2\}\)}{No solution exists.}{The lines are skew for every real \(m\).}{\(m = -2\)}{}{}}
\LANGpl{
\Question{
Wyznacz rzeczywistą wartość parametru \(m\) 
tak, aby proste \(p\) 
i \(q\) 
były prostymi skośnymi. 
\[ 
\begin{aligned}p\colon x& = 1 + t, &
\\y & = 2 - t,
\\z & = 1 - t;\ t\in \mathbb{R}
\\ \end{aligned}\qquad \qquad \begin{aligned}q\colon x& = s, &
\\y & = 1 + s,
\\z & = 3 + ms;\ s\in \mathbb{R}
\\ \end{aligned}
\]}
{\(m\in\mathbb{R}\setminus\{-2\}\)}{Brak rozwiązania.}{Proste są skośne dla każdej  rzeczywistej wartości  \(m\).}{\(m = -2\)}{}{}}
\LANGcs{
\Question{
Určete parametr
\(m\) tak, aby přímky \(p\colon x = 1 + t;\: y = 2 - t;\: z = 1 - t,\: t\in \mathbb{R}\) 
a \(q\colon x = s;\: y = 1 + s;\: z = 3 + ms,\: s\in \mathbb{R}\) 
byly mimoběžné.}
{\(m\in\mathbb{R}\setminus\{-2\}\)}{Pro žádné reálné \(m\) 
nejsou dané přímky mimoběžné.}{Pro každé reálné \(m\) 
jsou dané přímky mimoběžné.}{\(m = -2\)}{}{}}
\LANGsk{
\Question{
Určte hodnotu reálneho parametra
\(m\) tak, aby priamky \(p\colon x = 1 + t;\: y = 2 - t;\: z = 1 - t,\: t\in \mathbb{R}\) 
a \(q\colon x = s;\: y = 1 + s;\: z = 3 + ms,\: s\in \mathbb{R}\) 
byly mimobežné.}
{\(m\in\mathbb{R}\setminus\{-2\}\)}{Pre žiadne reálne \(m\) 
nie sú dané priamky mimobežné.}{Pre každé reálne \(m\) 
sú dané priamky mimobežné.}{\(m = -2\)}{}{}}
\LANGes{
\Question{
Elige de las respuestas dadas el valor del parámetro real \(m\) para el que las rectas \(p\) y \(q\) sean rectas sesgadas.
\[ 
\begin{aligned}p\colon x& = 1 + t, &
\\y & = 2 - t,
\\z & = 1 - t;\ t\in \mathbb{R}
\\ \end{aligned}\qquad \qquad \begin{aligned}q\colon x& = s, &
\\y & = 1 + s,
\\z & = 3 + ms;\ s\in \mathbb{R}
\\ \end{aligned}
\]}
{\(m\in\mathbb{R}\setminus\{-2\}\)}{No hay solución.}{Las rectas son rectas sesgadas para cualquier \(m\) real.}{\(m = -2\)}{}{}}
\End

\Data\ID{9000101005}
\Updated{2020-07-15 03:07:37}
\Level{A}
\SubArea{Analytic geometry in a space}
\LANGen{
\Question{
Find the value of the real parameter \(m\) 
which ensures that the lines \(p\) 
and \(q\) 
are intersecting lines (with a unique common point). 
\[ 
\begin{aligned}p\colon x& = 1 + t, &
\\y & = 2 - t,
\\z & = 1 - t;\ t\in \mathbb{R}
\\ \end{aligned}\qquad \qquad \begin{aligned}q\colon x& = s, &
\\y & = 1 + s,
\\z & = 3 + ms;\ s\in \mathbb{R}
\\ \end{aligned}
\]}
{\(m = -2\)}{No solution exists.}{The lines are intersecting for every real
\(m\).}{\(m = 2\)}{}{}}
\LANGpl{
\Question{
Wyznacz rzeczywistą wartość parametru \(m\) 
tak, aby proste \(p\) 
i \(q\) 
były prostymi przecinającymi się. 
\[ 
\begin{aligned}p\colon x& = 1 + t, &
\\y & = 2 - t,
\\z & = 1 - t;\ t\in \mathbb{R}
\\ \end{aligned}\qquad \qquad \begin{aligned}q\colon x& = s, &
\\y & = 1 + s,
\\z & = 3 + ms;\ s\in \mathbb{R}
\\ \end{aligned}
\]}
{\(m = -2\)}{Brak rozwiązań.}{Proste są przecinające się  dla każdej rzeczywistej wartości 
\(m\).}{\(m = 2\)}{}{}}
\LANGcs{
\Question{
Určete hodnotu reálného parametru
\(m\) tak, aby
přímky \(p\colon x = 1 + t;\: y = 2 - t;\: z = 1 - t,\: t\in \mathbb{R}\) 
a \(q\colon x = s;\: y = 1 + s;\: z = 3 + ms,\: s\in \mathbb{R}\) 
byly různoběžné.}
{\(m = -2\)}{Pro žádné reálné \(m\) 
nejsou dané přímky různoběžné.}{Pro každé reálné \(m\) 
jsou dané přímky různoběžné.}{\(m = 2\)}{}{}}
\LANGsk{
\Question{
Určte hodnotu reálneho parametra
\(m\) tak, aby
priamky \(p\colon x = 1 + t;\: y = 2 - t;\: z = 1 - t,\: t\in \mathbb{R}\) 
a \(q\colon x = s;\: y = 1 + s;\: z = 3 + ms,\: s\in \mathbb{R}\) 
boli rôznobežné.}
{\(m = -2\)}{Pre žiadne reálne \(m\) 
nie sú dané priamky rôznobežné.}{Pre každé reálne \(m\) 
sú dané priamky rôznobežné.}{\(m = 2\)}{}{}}
\LANGes{
\Question{
Determina el valor del parámetro real \(m\) 
para que las rectas \(p\) y \(q\) 
sean rectas no paralelas.
\[ 
\begin{aligned}p\colon x& = 1 + t, &
\\y & = 2 - t,
\\z & = 1 - t;\ t\in \mathbb{R}
\\ \end{aligned}\qquad \qquad \begin{aligned}q\colon x& = s, &
\\y & = 1 + s,
\\z & = 3 + ms;\ s\in \mathbb{R}
\\ \end{aligned}
\]}
{\(m = -2\)}{No hay solución.}{Las rectas son rectas no paralelas para cualquier \(m\) real.}{\(m = 2\)}{}{}}
\End

\Data\ID{9000101006}
\Updated{2020-07-15 03:07:37}
\Level{A}
\SubArea{Analytic geometry in a space}
\LANGen{
\Question{
Find the value of the real parameter \(m\) 
which ensures that the following lines are parallel and not identical lines.
\[ 
\begin{aligned}p\colon x& = 1 + t, &
\\y & = 2 - t,
\\z & = 1 - t;\  t\in \mathbb{R}
\\ \end{aligned}\qquad \qquad \begin{aligned}q\colon x& = s, &
\\y & = 1 + s,
\\z & = 3 + ms;\ s\in \mathbb{R}
\\ \end{aligned}
\]}
{No solution exists.}{The lines are parallel and not identical for every real
\(m\).}{\(m = -2\)}{\(m = 2\)}{}{}}
\LANGpl{
\Question{
Wyznacz rzeczywistą wartość parametru \(m\) 
tak, aby proste były równoległe i nie pokrywające się.
\[ 
\begin{aligned}p\colon x& = 1 + t, &
\\y & = 2 - t,
\\z & = 1 - t;\  t\in \mathbb{R}
\\ \end{aligned}\qquad \qquad \begin{aligned}q\colon x& = s, &
\\y & = 1 + s,
\\z & = 3 + ms;\ s\in \mathbb{R}
\\ \end{aligned}
\]}
{Brak rozwiązań.}{Proste są równoległe i nie pokrywające się dla każdej rzeczywistej wartości 
\(m\).}{\(m = -2\)}{\(m = 2\)}{}{}}
\LANGcs{
\Question{
Určete hodnotu reálného parametru
\(m\) tak, aby
přímky \(p\colon x = 1 + t;\: y = 2 - t;\: z = 1 - t,\: t\in \mathbb{R}\) 
a \(q\colon x = s;\: y = 1 + s;\: z = 3 + ms,\: s\in \mathbb{R}\) 
byly rovnoběžné různé.}
{Pro žádné reálné \(m\) 
nejsou dané přímky rovnoběžné různé.}{Pro každé reálné \(m\) 
jsou dané přímky rovnoběžné různé.}{\(m = -2\)}{\(m = 2\)}{}{}}
\LANGsk{
\Question{
Určte hodnotu reálneho parametra
\(m\) tak, aby
priamky \(p\colon x = 1 + t;\: y = 2 - t;\: z = 1 - t,\: t\in \mathbb{R}\) 
a \(q\colon x = s;\: y = 1 + s;\: z = 3 + ms,\: s\in \mathbb{R}\) 
boli rovnobežné rôzne.}
{Pre žiadne reálne \(m\) 
nie sú dané priamky rovnobežné rôzne.}{Pre každé reálne \(m\) 
sú dané priamky rovnobežné rôzne.}{\(m = -2\)}{\(m = 2\)}{}{}}
\LANGes{
\Question{
Halla el valor del parámetro real \(m\) para que las rectas siguientes sean paralelas no idénticas.
\[ 
\begin{aligned}p\colon x& = 1 + t, &
\\y & = 2 - t,
\\z & = 1 - t;\  t\in \mathbb{R}
\\ \end{aligned}\qquad \qquad \begin{aligned}q\colon x& = s, &
\\y & = 1 + s,
\\z & = 3 + ms;\ s\in \mathbb{R}
\\ \end{aligned}
\]}
{No hay solución.}{Las rectas son paralelas no idénticas para cada valor \(m\) real.}{\(m = -2\)}{\(m = 2\)}{}{}}
\End

\Data\ID{9000101007}
\Updated{2020-07-15 03:07:37}
\Level{A}
\SubArea{Analytic geometry in a space}
\LANGen{
\Question{
Find the value of the real parameter \(m\) 
which ensures that the following two lines are identical. 
\[ 
\begin{aligned}p\colon x& = 1 + t, &
\\y & = 2 - t,
\\z & = 1 - t;\ t\in \mathbb{R}
\\ \end{aligned}\qquad \qquad \begin{aligned}q\colon x& = s, &
\\y & = 1 + s,
\\z & = 3 + ms;\  s\in \mathbb{R}
\\ \end{aligned}
\]}
{No solution exists.}{The lines are identical for every real \(m\).}{\(m = -2\)}{\(m = 2\)}{}{}}
\LANGpl{
\Question{
Wyznacz rzeczywistą wartość parametru \(m\) 
tak, aby podane proste były prostymi pokrywającymi się. 
\[ 
\begin{aligned}p\colon x& = 1 + t, &
\\y & = 2 - t,
\\z & = 1 - t;\ t\in \mathbb{R}
\\ \end{aligned}\qquad \qquad \begin{aligned}q\colon x& = s, &
\\y & = 1 + s,
\\z & = 3 + ms;\  s\in \mathbb{R}
\\ \end{aligned}
\]}
{Brak rozwiązania.}{Proste są prostymi pokrywającymi się dla każdej rzeczywistej wartości  \(m\).}{\(m = -2\)}{\(m = 2\)}{}{}}
\LANGcs{
\Question{
Určete hodnotu reálného parametru
\(m\) tak, aby
přímky \(p\colon x = 1 + t;\: y = 2 - t;\: z = 1 - t,\: t\in \mathbb{R}\) 
a \(q\colon x = s;\: y = 1 + s;\: z = 3 + ms,\: s\in \mathbb{R}\) 
byly totožné.}
{Pro žádné reálné \(m\) 
nejsou dané přímky totožné.}{Pro každé reálné \(m\) 
jsou dané přímky totožné.}{\(m = -2\)}{\(m = 2\)}{}{}}
\LANGsk{
\Question{
Určte hodnotu reálneho parametra
\(m\) tak, aby
priamky \(p\colon x = 1 + t;\: y = 2 - t;\: z = 1 - t,\: t\in \mathbb{R}\) 
a \(q\colon x = s;\: y = 1 + s;\: z = 3 + ms,\: s\in \mathbb{R}\) 
boli totožné.}
{Pre žiadne reálne \(m\) 
nie sú dané priamky totožné.}{Pre každé reálne \(m\) 
sú dané priamky totožné.}{\(m = -2\)}{\(m = 2\)}{}{}}
\LANGes{
\Question{
Halla el valor del parámetro real \(m\) para que las rectas siguientes sean idénticas.
\[ 
\begin{aligned}p\colon x& = 1 + t, &
\\y & = 2 - t,
\\z & = 1 - t;\ t\in \mathbb{R}
\\ \end{aligned}\qquad \qquad \begin{aligned}q\colon x& = s, &
\\y & = 1 + s,
\\z & = 3 + ms;\  s\in \mathbb{R}
\\ \end{aligned}
\]}
{No hay solución.}{Las rectas son idénticas para cada \(m\) real.}{\(m = -2\)}{\(m = 2\)}{}{}}
\End

\Data\ID{9000101008}
\Updated{2020-07-15 03:07:37}
\Level{A}
\SubArea{Analytic geometry in a space}
\LANGen{
\Question{
Given points \(A = [0;2;-1]\),
\(B = [1;3;-3]\),
\(C = [-1;1;2]\),
\(D = [-2;0;3]\),
find the point which is the intersection of the lines
\(AB\) and
\(CD\)?}
{Point \(D\)}{Point \(A\)}{Point \(B\)}{Point \(C\)}{}{}}
\LANGpl{
\Question{
Dane są punkty \(A = [0;2;-1]\),
\(B = [1;3;-3]\),
\(C = [-1;1;2]\),
\(D = [-2;0;3]\),
wyznacz punkt przecięcia prostych
\(AB\) i
\(CD\)}
{Punkt \(D\)}{Punkt \(A\)}{Punkt \(B\)}{Punkt \(C\)}{}{}}
\LANGcs{
\Question{
Jsou dány body \(A = [0;2;-1]\),
\(B = [1;3;-3]\),
\(C = [-1;1;2]\),
\(D = [-2;0;3]\).
Který z těchto bodů je průsečíkem přímek
\(AB\) a
\(CD\)?}
{Bod \(D\).}{Bod \(A\).}{Bod \(B\).}{Bod \(C\).}{}{}}
\LANGsk{
\Question{
Sú dané body \(A = [0;2;-1]\),
\(B = [1;3;-3]\),
\(C = [-1;1;2]\),
\(D = [-2;0;3]\).
Ktorý z týchto bodov je priesečníkom priamok
\(AB\) a
\(CD\)?}
{Bod \(D\).}{Bod \(A\).}{Bod \(B\).}{Bod \(C\).}{}{}}
\LANGes{
\Question{
Dados los puntos \(A = [0;2;-1]\),
\(B = [1;3;-3]\),
\(C = [-1;1;2]\),
\(D = [-2;0;3]\),
encuentra cuál de estos puntos es la intersección de las rectas \(AB\) y
\(CD\)?}
{El punto \(D\)}{El punto \(A\)}{El punto \(B\)}{El punto \(C\)}{}{}}
\End

\Data\ID{9000101009}
\Updated{2020-07-15 03:07:37}
\Level{A}
\SubArea{Analytic geometry in a space}
\LANGen{
\Question{
Determine whether the following two lines are identical, parallel, intersecting or
skew.
\[\begin{aligned}
 a\colon x & = t, & &
 \\y & = -t, & &
 \\z & = 1 - t;\ t\in \mathbb{R} & &
\end{aligned}\]\[\begin{aligned}
 b\colon x & = -s, & &
 \\y & = s, & &
 \\z & = 1 + s;\ s\in \mathbb{R} & &
\end{aligned}\]}
{identical lines}{skew lines}{intersecting lines}{parallel, not identical lines}{}{}}
\LANGpl{
\Question{
Określ wzajemne położenie prostych a i b  w przestrzeni  jeśli

\[\begin{aligned}
 a\colon x & = t, & &
 \\y & = -t, & &
 \\z & = 1 - t;\ t\in \mathbb{R} & &
\end{aligned}\]\[\begin{aligned}
 b\colon x & = -s, & &
 \\y & = s, & &
 \\z & = 1 + s;\ s\in \mathbb{R} & &
\end{aligned}\]}
{proste pokrywające się}{proste skośne}{proste przecinające się}{proste równoległe, nie pokrywające się}{}{}}
\LANGcs{
\Question{
Určete vzájemnou polohu přímek
\(a\),
\(b\),
kde:

\[\begin{aligned}
 a\colon x & = t, & &
 \\y & = -t, & &
 \\z & = 1 - t;\ t\in \mathbb{R} & &
\end{aligned}\]\[\begin{aligned}
 b\colon x & = -s, & &
 \\y & = s, & &
 \\z & = 1 + s;\ s\in \mathbb{R} & &
\end{aligned}\]}
{Dané přímky jsou totožné.}{Dané přímky jsou mimoběžné.}{Dané přímky jsou různoběžné.}{Dané přímky jsou rovnoběžné různé.}{}{}}
\LANGsk{
\Question{
Určte vzájomnú polohu priamok
\(a\),
\(b\),
kde:

\[\begin{aligned}
 a\colon x & = t, & &
 \\y & = -t, & &
 \\z & = 1 - t;\ t\in \mathbb{R} & &
\end{aligned}\]\[\begin{aligned}
 b\colon x & = -s, & &
 \\y & = s, & &
 \\z & = 1 + s;\ s\in \mathbb{R} & &
\end{aligned}\]}
{Dané priamky sú totožné.}{Dané priamky sú mimobežné.}{Dané priamky sú rôznobežné.}{Dané priamky sú rovnobežné rôzne.}{}{}}
\LANGes{
\Question{
Determina la posición relativa de estas dos rectas.
\[\begin{aligned}
 a\colon x & = t, & &
 \\y & = -t, & &
 \\z & = 1 - t;\ t\in \mathbb{R} & &
\end{aligned}\]\[\begin{aligned}
 b\colon x & = -s, & &
 \\y & = s, & &
 \\z & = 1 + s;\ s\in \mathbb{R} & &
\end{aligned}\]}
{Son rectas idénticas.}{Son rectas sesgadas.}{Son rectas no paralelas.}{Son rectas paralelas no idénticas.}{}{}}
\End

\Data\ID{9000101010}
\Updated{2020-07-15 03:07:37}
\Level{A}
\SubArea{Analytic geometry in a space}
\LANGen{
\Question{
Determine whether the following two lines are identical, parallel, intersecting or
skew.

\[\begin{aligned}
 a\colon x & = t, & &
 \\y & = -t, & &
 \\z & = 1 - t;\ t\in \mathbb{R} & &
\end{aligned}\]\[\begin{aligned}
 b\colon x & = -s, & &
 \\y & = s, & &
 \\z & = -1 + s;\ s\in \mathbb{R} & &
\end{aligned}\]}
{parallel, not identical lines}{skew lines}{intersecting lines}{identical lines}{}{}}
\LANGpl{
\Question{
Określ wzajemne położenie prostych a i b w przestrzeni jeśli

\[\begin{aligned}
 a\colon x & = t, & &
 \\y & = -t, & &
 \\z & = 1 - t;\ t\in \mathbb{R} & &
\end{aligned}\]\[\begin{aligned}
 b\colon x & = -s, & &
 \\y & = s, & &
 \\z & = -1 + s;\ s\in \mathbb{R} & &
\end{aligned}\]}
{proste równoległe, nie pokrywające się}{proste skośne}{proste przecinające się}{proste pokrywające się}{}{}}
\LANGcs{
\Question{
Určete vzájemnou polohu přímek
\(a\),
\(b\),
kde:

\[\begin{aligned}
 a\colon x & = t, & &
 \\y & = -t, & &
 \\z & = 1 - t;\ t\in \mathbb{R} & &
\end{aligned}\]\[\begin{aligned}
 b\colon x & = -s, & &
 \\y & = s, & &
 \\z & = -1 + s;\ s\in \mathbb{R} & &
\end{aligned}\]}
{Dané přímky jsou rovnoběžné různé.}{Dané přímky jsou mimoběžné.}{Dané přímky jsou různoběžné.}{Dané přímky jsou totožné.}{}{}}
\LANGsk{
\Question{
Určte vzájomnú polohu priamok
\(a\),
\(b\),
kde:
\[\begin{aligned}
 a\colon x & = t, & &
 \\y & = -t, & &
 \\z & = 1 - t;\ t\in \mathbb{R} & &
\end{aligned}\]\[\begin{aligned}
 b\colon x & = -s, & &
 \\y & = s, & &
 \\z & = -1 + s;\ s\in \mathbb{R} & &
\end{aligned}\]}
{Dané priamky sú rovnobežné rôzne.}{Dané priamky sú mimobežné.}{Dané priamky sú rôznobežné.}{Dané priamky sú totožné.}{}{}}
\LANGes{
\Question{
Determina la posición relativa de estas dos rectas:

\[\begin{aligned}
 a\colon x & = t, & &
 \\y & = -t, & &
 \\z & = 1 - t;\ t\in \mathbb{R} & &
\end{aligned}\]\[\begin{aligned}
 b\colon x & = -s, & &
 \\y & = s, & &
 \\z & = -1 + s;\ s\in \mathbb{R} & &
\end{aligned}\]}
{Son rectas paralelas no idénticas.}{Son rectas sesgadas.}{Son rectas no paralelas.}{Son rectas idénticas.}{}{}}
\End

\Data\ID{9000106601}
\Updated{2020-07-15 03:07:37}
\Level{A}
\SubArea{Analytic geometry in a space}
\LANGen{
\Question{
Determine whether the following two lines are identical, parallel, intersecting or
skew. 
\[ 
\begin{aligned}[t] p\colon x& = -6 - t,&
\\y & = 7 + t,
\\z & = -2t;\ t\in \mathbb{R}
\\ \end{aligned}\qquad \qquad \begin{aligned}[t] q\colon x& = -1 - 2s, &
\\y & = 2 + 2s,
\\z & = 10 - 4s;\ s\in \mathbb{R}
\\ \end{aligned}
\]}
{identical lines}{parallel lines, not identical}{intersecting lines}{skew lines}{}{}}
\LANGpl{
\Question{
Określ wzajemne położenie prostych w przestrzeni. 

\[ 
\begin{aligned}[t] p\colon x& = -6 - t,&
\\y & = 7 + t,
\\z & = -2t;\ t\in \mathbb{R}
\\ \end{aligned}\qquad \qquad \begin{aligned}[t] q\colon x& = -1 - 2s, &
\\y & = 2 + 2s,
\\z & = 10 - 4s;\ s\in \mathbb{R}
\\ \end{aligned}
\]}
{proste pokrywające się}{proste równoległe, nie pokrywające się}{proste przecinające się}{proste skośne}{}{}}
\LANGcs{
\Question{
Určete vzájemnou polohu přímek
\(p\) a
\(q\) v
prostoru, je-li
\[\begin{aligned}
 p & = \{[-6 - t;\ 7 + t;\ -2t]\text{,}\ t\in \mathbb{R}\}\text{,} & &
 \\q & = \{[-1 - 2s;\ 2 + 2s;\ 10 - 4s]\text{,}\ s\in \mathbb{R}\}\text{.} & &
\end{aligned}\]}
{Dané přímky jsou totožné.}{Dané přímky jsou rovnoběžné různé.}{Dané přímky jsou různoběžné.}{Dané přímky jsou mimoběžné.}{}{}}
\LANGsk{
\Question{
Určte vzájomnú polohu priamok
\(p\) a
\(q\) v
priestore, ak je
\[\begin{aligned}
 p & = \{[-6 - t;\ 7 + t;\ -2t]\text{,}\ t\in \mathbb{R}\}\text{,} & &
 \\q & = \{[-1 - 2s;\ 2 + 2s;\ 10 - 4s]\text{,}\ s\in \mathbb{R}\}\text{.} & &
\end{aligned}\]}
{Dané priamky sú totožné.}{Dané priamky sú rovnobežné rôzne.}{Dané priamky sú rôznobežné.}{Dané priamky sú mimobežné.}{}{}}
\LANGes{
\Question{
Determina la posición relativa entre las siguientes rectas. 
\[ 
\begin{aligned}[t] p\colon x& = -6 - t,&
\\y & = 7 + t,
\\z & = -2t;\ t\in \mathbb{R}
\\ \end{aligned}\qquad \qquad \begin{aligned}[t] q\colon x& = -1 - 2s, &
\\y & = 2 + 2s,
\\z & = 10 - 4s;\ s\in \mathbb{R}
\\ \end{aligned}
\]}
{rectas idénticas}{rectas paralelas no idénticas}{rectas no paralelas}{rectas sesgadas}{}{}}
\End

\Data\ID{9000106602}
\Updated{2020-07-15 03:07:37}
\Level{A}
\SubArea{Analytic geometry in a space}
\LANGen{
\Question{
Determine whether the following two lines are identical, parallel, intersecting or
skew. 
\[ 
\begin{aligned}[t] p\colon x& = -3 + 2t,&
\\y & = 1 - t,
\\z & = 3 - 2t;\ t\in \mathbb{R}
\\ \end{aligned}\qquad \qquad \begin{aligned}[t] q\colon x& = 2 - 4s, &
\\y & = -3 + 2s,
\\z & = 6 + 4s;\ s\in \mathbb{R}
\\ \end{aligned}
\]}
{parallel lines, not identical}{identical lines}{intersecting lines}{skew lines}{}{}}
\LANGpl{
\Question{
Określ wzajemne położenie prostych w przestrzeni.  
\[ 
\begin{aligned}[t] p\colon x& = -3 + 2t,&
\\y & = 1 - t,
\\z & = 3 - 2t;\ t\in \mathbb{R}
\\ \end{aligned}\qquad \qquad \begin{aligned}[t] q\colon x& = 2 - 4s, &
\\y & = -3 + 2s,
\\z & = 6 + 4s;\ s\in \mathbb{R}
\\ \end{aligned}
\]}
{proste równoległe, nie pokrywające się}{proste pokrywające się}{proste przecinające się}{proste skośne}{}{}}
\LANGcs{
\Question{
Určete vzájemnou polohu přímek
\(p\) a
\(q\) v
prostoru, je-li
\[\begin{aligned}
 p & = \{[-3 + 2t;\ 1 - t;\ 3 - 2t]\text{,}\ t\in \mathbb{R}\}\text{,} & &
 \\q & = \{[2 - 4s;\ -3 + 2s;\ 6 + 4s]\text{,}\ s\in \mathbb{R}\}\text{.} & &
\end{aligned}\]}
{Dané přímky jsou rovnoběžné různé.}{Dané přímky jsou totožné.}{Dané přímky jsou různoběžné.}{Dané přímky jsou mimoběžné.}{}{}}
\LANGsk{
\Question{
Určte vzájomnú polohu priamok
\(p\) a
\(q\) v
priestore, ak
\[\begin{aligned}
 p & = \{[-3 + 2t;\ 1 - t;\ 3 - 2t]\text{,}\ t\in \mathbb{R}\}\text{,} & &
 \\q & = \{[2 - 4s;\ -3 + 2s;\ 6 + 4s]\text{,}\ s\in \mathbb{R}\}\text{.} & &
\end{aligned}\]}
{Dané priamky sú rovnobežné rôzne.}{Dané priamky sú totožné.}{Dané priamky sú rôznobežné.}{Dané priamky sú mimobežné.}{}{}}
\LANGes{
\Question{
Determina la posición relativa de las siguientes rectas.  
\[ 
\begin{aligned}[t] p\colon x& = -3 + 2t,&
\\y & = 1 - t,
\\z & = 3 - 2t;\ t\in \mathbb{R}
\\ \end{aligned}\qquad \qquad \begin{aligned}[t] q\colon x& = 2 - 4s, &
\\y & = -3 + 2s,
\\z & = 6 + 4s;\ s\in \mathbb{R}
\\ \end{aligned}
\]}
{rectas paralelas no idénticas}{rectas idénticas}{rectas no paralelas}{rectas sesgadas}{}{}}
\End

\Data\ID{9000106603}
\Updated{2020-07-15 03:07:37}
\Level{A}
\SubArea{Analytic geometry in a space}
\LANGen{
\Question{
Determine whether the following two lines are identical, parallel, intersecting or
skew. 
\[ 
\begin{aligned}[t] p\colon x& = -1 - t, &
\\y & = 11 - 2t,
\\z & = 1 + t;\ t\in \mathbb{R}
\\ \end{aligned}\qquad \qquad \begin{aligned}[t] q\colon x& = -3 + s, &
\\y & = 4 - s,
\\z & = 6 + 2s;\ s\in \mathbb{R}
\\ \end{aligned}
\]}
{intersecting lines}{parallel lines, not identical}{identical lines}{skew lines}{}{}}
\LANGpl{
\Question{
Określ wzajemne położenie prostych w przestrzeni.

\[ 
\begin{aligned}[t] p\colon x& = -1 - t, &
\\y & = 11 - 2t,
\\z & = 1 + t;\ t\in \mathbb{R}
\\ \end{aligned}\qquad \qquad \begin{aligned}[t] q\colon x& = -3 + s, &
\\y & = 4 - s,
\\z & = 6 + 2s;\ s\in \mathbb{R}
\\ \end{aligned}
\]}
{proste przecinające się}{proste równoległe, nie pokrywające się}{proste pokrywające się}{proste skośne}{}{}}
\LANGcs{
\Question{
Určete vzájemnou polohu přímek
\(p\) a
\(q\) v
prostoru, je-li
\[\begin{aligned}
 p & = \{[-1 - t;\ 11 - 2t;\ 1 + t]\text{,}\ t\in \mathbb{R}\}\text{,} & &
 \\q & = \{[-3 + s;\ 4 - s;\ 6 + 2s]\text{,}\ s\in \mathbb{R}\}\text{.} & &
\end{aligned}\]}
{Dané přímky jsou různoběžné.}{Dané přímky jsou rovnoběžné různé.}{Dané přímky jsou totožné.}{Dané přímky jsou mimoběžné.}{}{}}
\LANGsk{
\Question{
Určte vzájomnú polohu priamok
\(p\) a
\(q\) v
priestore, ak
\[\begin{aligned}
 p & = \{[-1 - t;\ 11 - 2t;\ 1 + t]\text{,}\ t\in \mathbb{R}\}\text{,} & &
 \\q & = \{[-3 + s;\ 4 - s;\ 6 + 2s]\text{,}\ s\in \mathbb{R}\}\text{.} & &
\end{aligned}\]}
{Dané priamky sú rôznobežné.}{Dané priamky sú rovnobežné rôzne.}{Dané priamky sú totožné.}{Dané priamky sú mimobežné.}{}{}}
\LANGes{
\Question{
Determina la posición relativa de las siguientes rectas. 
\[ 
\begin{aligned}[t] p\colon x& = -1 - t, &
\\y & = 11 - 2t,
\\z & = 1 + t;\ t\in \mathbb{R}
\\ \end{aligned}\qquad \qquad \begin{aligned}[t] q\colon x& = -3 + s, &
\\y & = 4 - s,
\\z & = 6 + 2s;\ s\in \mathbb{R}
\\ \end{aligned}
\]}
{rectas no paralelas}{rectas paralelas, no idénticas}{rectas idénticas}{rectas sesgadas}{}{}}
\End

\Data\ID{9000106604}
\Updated{2020-07-15 03:07:37}
\Level{A}
\SubArea{Analytic geometry in a space}
\LANGen{
\Question{
Determine whether the following two lines are identical, parallel, intersecting or
skew. 
\[ 
\begin{aligned}[t] p\colon x& = 1 + 3t&
\\y & = 2 - 6t
\\z & = 3t;\ t\in \mathbb{R}
\\ \end{aligned}\qquad \qquad \begin{aligned}[t] q\colon x& = 4 - 2s &
\\y & = 1 + 4s
\\z & = 3 - 2s;\ s\in \mathbb{R}
\\ \end{aligned}
\]}
{parallel lines, not identical}{identical lines}{intersecting lines}{skew lines}{}{}}
\LANGpl{
\Question{
Określ wzajemne położenie prostych w przestrzeni
\[ 
\begin{aligned}[t] p\colon x& = 1 + 3t&
\\y & = 2 - 6t
\\z & = 3t;\ t\in \mathbb{R}
\\ \end{aligned}\qquad \qquad \begin{aligned}[t] q\colon x& = 4 - 2s &
\\y & = 1 + 4s
\\z & = 3 - 2s;\ s\in \mathbb{R}
\\ \end{aligned}
\]}
{proste równoległe, nie pokrywające się}{proste pokrywające się}{proste przecinające się}{proste skośne}{}{}}
\LANGcs{
\Question{
Určete vzájemnou polohu přímek
\(p\) a
\(q\) v
prostoru, je-li
\[\begin{aligned}
 p = \{[1 + 3t;\ 2 - 6t;\ 3t]\text{,}\ t\in \mathbb{R}\}, & &q\colon &x = 4 - 2s, & & & &
 \\ & & &y = 1 + 4s, & & & &
 \\ & & &z = 3 - 2s;\ s\in \mathbb{R}\text{.} & & & &
\end{aligned}\]}
{Dané přímky jsou rovnoběžné různé.}{Dané přímky jsou totožné.}{Dané přímky jsou různoběžné.}{Dané přímky jsou mimoběžné.}{}{}}
\LANGsk{
\Question{
Určte vzájomnú polohu priamok
\(p\) a
\(q\) v
priestore, ak
\[\begin{aligned}
 p = \{[1 + 3t;\ 2 - 6t;\ 3t], t\in \mathbb{R}\}\text{,} & &q\colon &x = 4 - 2s, & & & &
 \\ & & &y = 1 + 4s, & & & &
 \\ & & &z = 3 - 2s;\ s\in \mathbb{R}. & & & &
\end{aligned}\]}
{Dané priamky sú rovnobežné rôzne.}{Dané priamky sú totožné.}{Dané priamky sú rôznobežné.}{Dané priamky sú mimobežné.}{}{}}
\LANGes{
\Question{
Determina la posición relativa de las siguientes rectas. 
\[ 
\begin{aligned}[t] p\colon x& = 1 + 3t&
\\y & = 2 - 6t
\\z & = 3t;\ t\in \mathbb{R}
\\ \end{aligned}\qquad \qquad \begin{aligned}[t] q\colon x& = 4 - 2s &
\\y & = 1 + 4s
\\z & = 3 - 2s;\ s\in \mathbb{R}
\\ \end{aligned}
\]}
{rectas paralelas, no idénticas}{rectas idénticas}{rectas no paralelas}{rectas sesgadas}{}{}}
\End

\Data\ID{9000106605}
\Updated{2020-07-15 03:07:37}
\Level{A}
\SubArea{Analytic geometry in a space}
\LANGen{
\Question{
Determine whether the following two lines are identical, parallel, intersecting or
skew. 
\[ 
\begin{aligned}[t] p\colon x& = 5 - 3t, &
\\y & = t,
\\z & = 5 - t;\ t\in \mathbb{R}
\\ \end{aligned}\qquad \qquad \begin{aligned}[t] q\colon x& = -4 + 3s,&
\\y & = 3 - s,
\\z & = 2 + s;\ s\in \mathbb{R}
\\ \end{aligned}
\]}
{identical lines}{parallel lines, not identical}{intersecting lines}{skew lines}{}{}}
\LANGpl{
\Question{
Określ wzajemne położenie prostych w przestrzeni. 
\[ 
\begin{aligned}[t] p\colon x& = 5 - 3t, &
\\y & = t,
\\z & = 5 - t;\ t\in \mathbb{R}
\\ \end{aligned}\qquad \qquad \begin{aligned}[t] q\colon x& = -4 + 3s,&
\\y & = 3 - s,
\\z & = 2 + s;\ s\in \mathbb{R}
\\ \end{aligned}
\]}
{proste pokrywające się}{proste równoległe, nie pokrywające się}{proste przecinające się}{proste skośne}{}{}}
\LANGcs{
\Question{
Určete vzájemnou polohu přímek
\(p\) a
\(q\) v
prostoru, je-li
\[\begin{aligned}
 p = \{[5 - 3t;\ t;\ 5 - t]\text{,}\ t\in \mathbb{R}\}, & &q\colon &x = -4 + 3s, & & & &
 \\ & & &y =\phantom{ -}3 -\phantom{ 3}s, & & & &
 \\ & & &z =\phantom{ -}2 +\phantom{ 3}s;\ s\in \mathbb{R}. & & & &
\end{aligned}\]}
{Dané přímky jsou totožné.}{Dané přímky jsou rovnoběžné různé.}{Dané přímky jsou různoběžné.}{Dané přímky jsou mimoběžné.}{}{}}
\LANGsk{
\Question{
Určte vzájomnú polohu priamok
\(p\) a
\(q\) v
priestore, ak
\[\begin{aligned}
 p = \{[5 - 3t;\ t;\ 5 - t]\text{,}\ t\in \mathbb{R}\}, & &q\colon &x = -4 + 3s, & & & &
 \\ & & &y =\phantom{ -}3 -\phantom{ 3}s, & & & &
 \\ & & &z =\phantom{ -}2 +\phantom{ 3}s;\ s\in \mathbb{R}. & & & &
\end{aligned}\]}
{Dané priamky sú totožné.}{Dané priamky sú rovnobežné rôzne.}{Dané priamky sú rôznobežné.}{Dané priamky sú mimobežné.}{}{}}
\LANGes{
\Question{
Determina la posición relativa de las siguientes rectas. 
\[ 
\begin{aligned}[t] p\colon x& = 5 - 3t, &
\\y & = t,
\\z & = 5 - t;\ t\in \mathbb{R}
\\ \end{aligned}\qquad \qquad \begin{aligned}[t] q\colon x& = -4 + 3s,&
\\y & = 3 - s,
\\z & = 2 + s;\ s\in \mathbb{R}
\\ \end{aligned}
\]}
{rectas idénticas}{rectas paralelas, no idénticas}{rectas no paralelas}{rectas sesgadas}{}{}}
\End

\Data\ID{9000106606}
\Updated{2020-07-15 03:07:37}
\Level{A}
\SubArea{Analytic geometry in a space}
\LANGen{
\Question{
Determine whether the following two lines are identical, parallel, intersecting or
skew. 
\[ 
\begin{aligned}[t] p\colon x& = 2t, &
\\y & = 3 - t,
\\z & = 4 - t;\ t\in \mathbb{R}
\\ \end{aligned}\qquad \qquad \begin{aligned}[t] q\colon x& = 2 - 2s, &
\\y & = -1 + s,
\\z & = 6 + 3s;\ s\in \mathbb{R}
\\ \end{aligned}
\]}
{skew lines}{parallel lines, not identical}{intersecting lines}{identical lines}{}{}}
\LANGpl{
\Question{
Określ wzajemne położenie prostych w przestrzeni. 
\[ 
\begin{aligned}[t] p\colon x& = 2t, &
\\y & = 3 - t,
\\z & = 4 - t;\ t\in \mathbb{R}
\\ \end{aligned}\qquad \qquad \begin{aligned}[t] q\colon x& = 2 - 2s, &
\\y & = -1 + s,
\\z & = 6 + 3s;\ s\in \mathbb{R}
\\ \end{aligned}
\]}
{proste skośne}{proste równoległe, nie pokrywające się}{proste przecinające się}{proste pokrywające się}{}{}}
\LANGcs{
\Question{
Určete vzájemnou polohu přímek
\(p\) a
\(q\) v
prostoru, je-li
\[\begin{aligned}
 p = \{[2t;\ 3 - t;\ 4 - t]\text{,}\ t\in \mathbb{R}\}, & &q\colon &x =\phantom{ -}2 - 2s, & & & &
 \\ & & &y = -1 +\phantom{ 4}s, & & & &
 \\ & & &z =\phantom{ -}6 + 3s;\ s\in \mathbb{R}. & & & &
\end{aligned}\]}
{Dané přímky jsou mimoběžné.}{Dané přímky jsou rovnoběžné různé.}{Dané přímky jsou různoběžné.}{Dané přímky jsou totožné.}{}{}}
\LANGsk{
\Question{
Určte vzájomnú polohu priamok
\(p\) a
\(q\) v
priestore, ak
\[\begin{aligned}
 p = \{[2t;\ 3 - t;\ 4 - t]\text{,}\ t\in \mathbb{R}\}, & &q\colon &x =\phantom{ -}2 - 2s, & & & &
 \\ & & &y = -1 +\phantom{ 4}s, & & & &
 \\ & & &z =\phantom{ -}6 + 3s;\ s\in \mathbb{R}. & & & &
\end{aligned}\]}
{Dané priamky sú mimobežné.}{Dané priamky sú rovnobežné rôzne.}{Dané priamky sú rôznobežné.}{Dané priamky sú totožné.}{}{}}
\LANGes{
\Question{
Determina la posición relativa de las siguientes rectas.
\[ 
\begin{aligned}[t] p\colon x& = 2t, &
\\y & = 3 - t,
\\z & = 4 - t;\ t\in \mathbb{R}
\\ \end{aligned}\qquad \qquad \begin{aligned}[t] q\colon x& = 2 - 2s, &
\\y & = -1 + s,
\\z & = 6 + 3s;\ s\in \mathbb{R}
\\ \end{aligned}
\]}
{rectas sesgadas}{rectas paralelas, no idénticas}{rectas no paralelas}{rectas idénticas}{}{}}
\End

\Data\ID{9000106607}
\Updated{2020-07-15 03:07:37}
\Level{A}
\SubArea{Analytic geometry in a space}
\LANGen{
\Question{
Determine whether the following two lines are identical, parallel, intersecting or
skew.
\[\begin{aligned}
 p\colon &x = 2, &q\colon &x =\phantom{ -}1 -\phantom{ 3}s, & & & &
 \\ &y = 3 -\phantom{ 2}t, & &y =\phantom{ -}2 + 3s, & & & &
 \\ &z = 3 + 2t;\ t\in \mathbb{R}, & &z = -1 - 2s;\ s\in \mathbb{R} & & & &
\end{aligned}\]}
{skew lines}{parallel lines, not identical}{intersecting lines}{identical lines}{}{}}
\LANGpl{
\Question{
Określ wzajemne położenie prostych w przestrzeni.
\[\begin{aligned}
 p\colon &x = 2, &q\colon &x =\phantom{ -}1 -\phantom{ 3}s, & & & &
 \\ &y = 3 -\phantom{ 2}t, & &y =\phantom{ -}2 + 3s, & & & &
 \\ &z = 3 + 2t;\ t\in \mathbb{R}, & &z = -1 - 2s;\ s\in \mathbb{R} & & & &
\end{aligned}\]}
{proste skośne}{proste równoległe, nie pokrywające się}{proste przecinające się}{proste pokrywające się}{}{}}
\LANGcs{
\Question{
Určete vzájemnou polohu přímek
\(p\) a
\(q\) v
prostoru, je-li
\[\begin{aligned}
 p\colon &x = 2, &q\colon &x =\phantom{ -}1 -\phantom{ 3}s, & & & &
 \\ &y = 3 -\phantom{ 2}t, & &y =\phantom{ -}2 + 3s, & & & &
 \\ &z = 3 + 2t;\ t\in \mathbb{R}, & &z = -1 - 2s;\ s\in \mathbb{R}. & & & &
\end{aligned}\]}
{Dané přímky jsou mimoběžné.}{Dané přímky jsou rovnoběžné různé.}{Dané přímky jsou různoběžné.}{Dané přímky jsou totožné.}{}{}}
\LANGsk{
\Question{
Určte vzájomnú polohu priamok
\(p\) a
\(q\) v
priestore, ak
\[\begin{aligned}
 p\colon &x = 2, &q\colon &x =\phantom{ -}1 -\phantom{ 3}s, & & & &
 \\ &y = 3 -\phantom{ 2}t, & &y =\phantom{ -}2 + 3s, & & & &
 \\ &z = 3 + 2t;\ t\in \mathbb{R}, & &z = -1 - 2s;\ s\in \mathbb{R}. & & & &
\end{aligned}\]}
{Dané priamky sú mimobežné.}{Dané priamky sú rovnobežné rôzne.}{Dané priamky sú rôznobežné.}{Dané priamky sú totožné.}{}{}}
\LANGes{
\Question{
Determina la posición relativa de las siguientes rectas.
\[\begin{aligned}
 p\colon &x = 2, &q\colon &x =\phantom{ -}1 -\phantom{ 3}s, & & & &
 \\ &y = 3 -\phantom{ 2}t, & &y =\phantom{ -}2 + 3s, & & & &
 \\ &z = 3 + 2t;\ t\in \mathbb{R}, & &z = -1 - 2s;\ s\in \mathbb{R} & & & &
\end{aligned}\]}
{rectas sesgadas}{rectas paralelas, no idénticas}{rectas no paralelas}{rectas idénticas}{}{}}
\End

\Data\ID{9000106608}
\Updated{2020-07-15 03:07:37}
\Level{A}
\SubArea{Analytic geometry in a space}
\LANGen{
\Question{
Determine whether the following two lines are identical, parallel, intersecting or
skew.
\[\begin{aligned}
 p\colon\, &x = 2, &q\colon\, &x =\phantom{ 1} - s, & & & &
 \\ &y = 2 + t, & &y = 4, & & & &
 \\ &z = 3;\ t\in \mathbb{R}, & &z = 1 - s;\ s\in \mathbb{R} & & & &
\end{aligned}\]}
{intersecting lines}{parallel lines, not identical}{skew lines}{identical lines}{}{}}
\LANGpl{
\Question{
Określ wzajemne położenie prostych w przestrzeni.

\[\begin{aligned}
 p\colon\, &x = 2, &q\colon\, &x =\phantom{ 1} - s, & & & &
 \\ &y = 2 + t, & &y = 4, & & & &
 \\ &z = 3;\ t\in \mathbb{R}, & &z = 1 - s;\ s\in \mathbb{R} & & & &
\end{aligned}\]}
{proste przecinające się}{proste równoległe, nie pokrywające się}{proste skośne}{proste pokrywające się}{}{}}
\LANGcs{
\Question{
Určete vzájemnou polohu přímek
\(p\) a
\(q\) v
prostoru, je-li

\[\begin{aligned}
 p\colon\, &x = 2, &q\colon\, &x =\phantom{ 1} - s, & & & &
 \\ &y = 2 + t, & &y = 4, & & & &
 \\ &z = 3;\ t\in \mathbb{R}, & &z = 1 - s;\ s\in \mathbb{R}. & & & &
\end{aligned}\]}
{Dané přímky jsou různoběžné.}{Dané přímky jsou rovnoběžné různé.}{Dané přímky jsou mimoběžné.}{Dané přímky jsou totožné.}{}{}}
\LANGsk{
\Question{
Určte vzájomnú polohu priamok
\(p\) a
\(q\) v
priestore, ak

\[\begin{aligned}
 p\colon\, &x = 2, &q\colon\, &x =\phantom{ 1} - s, & & & &
 \\ &y = 2 + t, & &y = 4, & & & &
 \\ &z = 3;\ t\in \mathbb{R}, & &z = 1 - s;\ s\in \mathbb{R}. & & & &
\end{aligned}\]}
{Dané priamky sú rôznobežné.}{Dané priamky sú rovnobežné rôzne.}{Dané priamky sú mimobežné.}{Dané priamky sú totožné.}{}{}}
\LANGes{
\Question{
Determina la posición relativa de las siguientes rectas.
\[\begin{aligned}
 p\colon\, &x = 2, &q\colon\, &x =\phantom{ 1} - s, & & & &
 \\ &y = 2 + t, & &y = 4, & & & &
 \\ &z = 3;\ t\in \mathbb{R}, & &z = 1 - s;\ s\in \mathbb{R} & & & &
\end{aligned}\]}
{rectas no paralelas}{rectas paralelas, no idénticas}{rectas sesgadas}{rectas idénticas}{}{}}
\End

\Data\ID{9000106609}
\Updated{2020-07-15 03:07:37}
\Level{A}
\SubArea{Analytic geometry in a space}
\LANGen{
\Question{
Determine whether two lines are identical, parallel, intersecting or skew. The first line is the line passes
through the points \(A = [3;-2;1]\) 
and \(B = [0;7;7]\) 
and the second line is the line passes through the points
\(C = [5;-8;-3]\) and
\(D = [6;-11;-5]\).}
{identical lines}{parallel lines, not identical}{intersecting lines}{skew lines}{}{}}
\LANGpl{
\Question{
Określ wzajemne położenie prostych w przestrzeni. Pierwsza prosta przechodzi 
przez punkty \(A = [3;-2;1]\) 
i \(B = [0;7;7]\) 
druga prosta przechodzi przez punkty
\(C = [5;-8;-3]\) i
\(D = [6;-11;-5]\).}
{proste pokrywające się}{proste równoległe, nie pokrywające się}{proste przecinające się}{proste skośne}{}{}}
\LANGcs{
\Question{
Určete vzájemnou polohu přímek
\(p\) a
\(q\) v prostoru,
je-li přímka \(p\) 
dána body \(A = [3;-2;1]\),
\(B = [0;7;7]\) a
přímka \(q\) 
body \(C = [5;-8;-3]\),
\(D = [6;-11;-5]\).}
{Dané přímky jsou totožné.}{Dané přímky jsou rovnoběžné různé.}{Dané přímky jsou různoběžné.}{Dané přímky jsou mimoběžné.}{}{}}
\LANGsk{
\Question{
Určte vzájomnú polohu priamok
\(p\) a
\(q\) v priestore,
ak je priamka \(p\) 
daná bodmi \(A = [3;-2;1]\),
\(B = [0;7;7]\) a
priamka \(q\) 
bodmi \(C = [5;-8;-3]\),
\(D = [6;-11;-5]\).}
{Dané priamky sú totožné.}{Dané priamky sú rovnobežné rôzne.}{Dané priamky sú rôznobežné.}{Dané priamky sú mimobežné.}{}{}}
\LANGes{
\Question{
Determina la posición relativa de las rectas
\(p\) y
\(q\). Los puntos \(A = [3;-2;1]\),
\(B = [0;7;7]\) pasan por la recta \(p\)
 y los puntos \(C = [5;-8;-3]\),
\(D = [6;-11;-5]\) pasan por la recta \(q\).}
{Las rectas son idénticas.}{Las rectas son paralelas, no idénticas.}{Las rectas no son paralelas.}{Son las rectas sesgadas.}{}{}}
\End

\Data\ID{9000106610}
\Updated{2020-07-15 03:07:37}
\Level{A}
\SubArea{Analytic geometry in a space}
\LANGen{
\Question{
Determine whether two lines are identical, parallel, intersecting or skew. The first line is the line passes
through the points \(A = [1;-4;2]\) 
and \(B = [3;0;0]\) 
and the second line is the line passes through the points
\(C = [3;-5;5]\) and
\(D = [-1;-3;-1]\).}
{intersecting lines}{parallel lines, not identical}{identical lines}{skew lines}{}{}}
\LANGpl{
\Question{
Określ wzajemne położenie prostych w przestrzeni. Pierwsza prosta 
przechodzi przez punkty \(A = [1;-4;2]\) 
i \(B = [3;0;0]\), druga prosta przechodzi przez punkty
\(C = [3;-5;5]\) i
\(D = [-1;-3;-1]\).}
{proste przecinające się}{proste równoległe, nie pokrywające się}{proste pokrywające się}{proste skośne}{}{}}
\LANGcs{
\Question{
Určete vzájemnou polohu přímek
\(p\) a
\(q\) v prostoru,
je-li přímka \(p\) 
dána body \(A = [1;-4;2]\),
\(B = [3;0;0]\) a
přímka \(q\) 
body \(C = [3;-5;5]\),
\(D = [-1;-3;-1]\).}
{Dané přímky jsou různoběžné.}{Dané přímky jsou rovnoběžné různé.}{Dané přímky jsou totožné.}{Dané přímky jsou mimoběžné.}{}{}}
\LANGsk{
\Question{
Určte vzájomnú polohu priamok
\(p\) a
\(q\) v priestore,
ak priamka \(p\) 
daná bodmi \(A = [1;-4;2]\),
\(B = [3;0;0]\) a
priamka \(q\) 
bodmi \(C = [3;-5;5]\),
\(D = [-1;-3;-1]\).}
{Dané priamky sú rôznobežné.}{Dané priamky sú rovnobežné rôzne.}{Dané priamky sú totožné.}{Dané priamky sú mimobežné.}{}{}}
\LANGes{
\Question{
Determina la posición relativa de la recta
\(p\) y
\(q\). Los puntos 
\(A = [1;-4;2]\),
\(B = [3;0;0]\) pasan por la recta \(p\) y
los puntos \(C = [3;-5;5]\),
\(D = [-1;-3;-1]\) pasan por la recta \(q\).}
{Las rectas no son paralelas.}{Las rectas son paralelas, no idénticas.}{Las rectas son idénticas.}{Son rectas sesgadas.}{}{}}
\End

\Data\ID{9000117402}
\Updated{2020-07-15 03:07:37}
\Level{A}
\SubArea{Analytic geometry in a space}
\LANGen{
\Question{
Determine whether the following planes
\(\rho \) and
\(\sigma \) are
parallel, identical or intersecting. 
\[ 
\begin{aligned}[t] \rho \colon &x = 2 + u - v, &
\\&y = 1 + 2u + 4v,
\\&z = -1 + 3u + 3v;\ u,v\in \mathbb{R},
\\ \end{aligned}\qquad \begin{aligned}[t] \sigma \colon &x = 2 + r - s, &
\\&y = 7 + 2r + 4s,
\\&z = 5 + 3r + 3s;\ s,t\in \mathbb{R}.
\\ \end{aligned}
\]}
{identical}{parallel, not identical}{intersecting}{}{}{}}
\LANGpl{
\Question{
Określ wzajemne położenie płaszczyzn
\(\rho \) i
\(\sigma \). 
\[ 
\begin{aligned}[t] \rho \colon &x = 2 + u - v, &
\\&y = 1 + 2u + 4v,
\\&z = -1 + 3u + 3v;\ u,v\in \mathbb{R},
\\ \end{aligned}\qquad \begin{aligned}[t] \sigma \colon &x = 2 + r - s, &
\\&y = 7 + 2r + 4s,
\\&z = 5 + 3r + 3s;\ s,t\in \mathbb{R}.
\\ \end{aligned}
\]}
{pokrywające się}{równoległe, nie pokrywające się}{przecinające się}{}{}{}}
\LANGcs{
\Question{
Jsou dány roviny \(\rho \) 
a \(\sigma \).
Určete jejich vzájemnou polohu. 
\[ 
\begin{aligned}[t] \rho \colon &x = 2 + u - v, &
\\&y = 1 + 2u + 4v,
\\&z = -1 + 3u + 3v;\ u,v\in \mathbb{R},
\\ \end{aligned}\qquad \begin{aligned}[t] \sigma \colon &x = 2 + r - s, &
\\&y = 7 + 2r + 4s,
\\&z = 5 + 3r + 3s;\ s,t\in \mathbb{R}.
\\ \end{aligned}
\]}
{Dané roviny jsou totožné.}{Dané roviny jsou rovnoběžné různé.}{Dané roviny jsou různoběžné.}{}{}{}}
\LANGsk{
\Question{
Sú dané roviny \(\rho \) 
a \(\sigma \).
Určte ich vzájomnú polohu. 
\[ 
\begin{aligned}[t] \rho \colon &x = 2 + u - v, &
\\&y = 1 + 2u + 4v,
\\&z = -1 + 3u + 3v;\ u,v\in \mathbb{R},
\\ \end{aligned}\qquad \begin{aligned}[t] \sigma \colon &x = 2 + r - s, &
\\&y = 7 + 2r + 4s,
\\&z = 5 + 3r + 3s;\ s,t\in \mathbb{R}.
\\ \end{aligned}
\]}
{Dané roviny sú totožné.}{Dané roviny sú rovnobežné rôzne.}{Dané roviny sú rôznobežné.}{}{}{}}
\LANGes{
\Question{
Determina la posición relativa entre los planos \(\rho \) y
\(\sigma \).
\[ 
\begin{aligned}[t] \rho \colon &x = 2 + u - v, &
\\&y = 1 + 2u + 4v,
\\&z = -1 + 3u + 3v;\ u,v\in \mathbb{R},
\\ \end{aligned}\qquad \begin{aligned}[t] \sigma \colon &x = 2 + r - s, &
\\&y = 7 + 2r + 4s,
\\&z = 5 + 3r + 3s;\ s,t\in \mathbb{R}.
\\ \end{aligned}
\]}
{Los planos son idénticos.}{Los planos son paralelos, no idénticos.}{Los planos no son paralelos.}{}{}{}}
\End

\Data\ID{9000117403}
\Updated{2020-07-15 03:07:37}
\Level{A}
\SubArea{Analytic geometry in a space}
\LANGen{
\Question{
Determine whether the following planes
\(\rho \) and
\(\sigma \) are
parallel, identical or intersecting. 
\[ 
\begin{aligned}[t] \rho \colon &x = -u + v, &
\\&y = u + 2v,
\\&z = -u - v;\ u,v\in \mathbb{R},
\\ \end{aligned}\qquad \sigma \colon x-2y-3z+1 = 0
\]}
{parallel, not identical}{identical}{intersecting}{}{}{}}
\LANGpl{
\Question{
Określ wzajemne położenie płaszczyzn
\(\rho \) i
\(\sigma \).
\[ 
\begin{aligned}[t] \rho \colon &x = -u + v, &
\\&y = u + 2v,
\\&z = -u - v;\ u,v\in \mathbb{R},
\\ \end{aligned}\qquad \sigma \colon x-2y-3z+1 = 0
\]}
{równoległe, nie pokrywające się}{pokrywające się}{przecinające się}{}{}{}}
\LANGcs{
\Question{
Jsou dány roviny \(\rho \) 
a \(\sigma \).
Určete jejich vzájemnou polohu. 
\[ 
\begin{aligned}[t] \rho \colon &x = -u + v, &
\\&y = u + 2v,
\\&z = -u - v;\ u,v\in \mathbb{R},
\\ \end{aligned}\qquad \sigma \colon x-2y-3z+1 = 0
\]}
{Dané roviny jsou rovnoběžné různé.}{Dané roviny jsou totožné.}{Dané roviny jsou různoběžné.}{}{}{}}
\LANGsk{
\Question{
Sú dané roviny \(\rho \) 
a \(\sigma \).
Určte ich vzájomnú polohu. 
\[ 
\begin{aligned}[t] \rho \colon &x = -u + v, &
\\&y = u + 2v,
\\&z = -u - v;\ u,v\in \mathbb{R},
\\ \end{aligned}\qquad \sigma \colon x-2y-3z+1 = 0
\]}
{Dané roviny sú rovnobežné rôzne.}{Dané roviny sú totožné.}{Dané roviny sú rôznobežné.}{}{}{}}
\LANGes{
\Question{
Determina la posición relativa entre los planos
\(\rho \) y 
\(\sigma \).
\[ 
\begin{aligned}[t] \rho \colon &x = -u + v, &
\\&y = u + 2v,
\\&z = -u - v;\ u,v\in \mathbb{R},
\\ \end{aligned}\qquad \sigma \colon x-2y-3z+1 = 0
\]}
{Los planos son paralelos, no idénticos.}{Los planos son idénticos.}{Los planos no son paralelos.}{}{}{}}
\End

\Data\ID{9000117404}
\Updated{2020-07-15 03:07:37}
\Level{A}
\SubArea{Analytic geometry in a space}
\LANGen{
\Question{
Determine whether the following planes are parallel, identical or intersecting.
\[\begin{aligned}
 \rho \colon \frac{3}
{8}x + \frac{1} 
{2}y -\frac{2} 
{3}z - 1 = 0,\qquad \sigma \colon \frac{3}
{4}x + y -\frac{4} 
{3}z - 2 = 0 & &
\end{aligned}\]}
{identical}{parallel, not identical}{intersecting}{}{}{}}
\LANGpl{
\Question{
Określ wzajemne położenie płaszczyzn 
\[\begin{aligned}
 \rho \colon \frac{3}
{8}x + \frac{1} 
{2}y -\frac{2} 
{3}z - 1 = 0,\qquad \sigma \colon \frac{3}
{4}x + y -\frac{4} 
{3}z - 2 = 0 & &
\end{aligned}\]}
{pokrywające się}{równoległe, nie pokrywające się}{przecinające się}{}{}{}}
\LANGcs{
\Question{
Jsou dány roviny
\[\begin{aligned}
 \rho \colon \frac{3}
{8}x + \frac{1} 
{2}y -\frac{2} 
{3}z - 1 = 0,\quad \sigma \colon \frac{3}
{4}x + y -\frac{4} 
{3}z - 2 = 0. & &
\end{aligned}\]
Určete jejich vzájemnou polohu.}
{Dané roviny jsou totožné.}{Dané roviny jsou rovnoběžné různé.}{Dané roviny jsou různoběžné.}{}{}{}}
\LANGsk{
\Question{
Sú dané roviny
\[\begin{aligned}
 \rho \colon \frac{3}
{8}x + \frac{1} 
{2}y -\frac{2} 
{3}z - 1 = 0,\quad \sigma \colon \frac{3}
{4}x + y -\frac{4} 
{3}z - 2 = 0. & &
\end{aligned}\]
Určte ich vzájomnú polohu.}
{Dané roviny sú totožné.}{Dané roviny sú rovnobežné rôzne.}{Dané roviny sú rôznobežné.}{}{}{}}
\LANGes{
\Question{
Determina la posición relativa entre los planos.
\[\begin{aligned}
 \rho \colon \frac{3}
{8}x + \frac{1} 
{2}y -\frac{2} 
{3}z - 1 = 0,\qquad \sigma \colon \frac{3}
{4}x + y -\frac{4} 
{3}z - 2 = 0 & &
\end{aligned}\]}
{Los planos son idénticos.}{Los planos son paralelos, no idénticos.}{Los planos no son paralelos.}{}{}{}}
\End

\Data\ID{9000117405}
\Updated{2020-07-15 03:07:37}
\Level{A}
\SubArea{Analytic geometry in a space}
\LANGen{
\Question{
Determine whether the following planes are parallel, identical or intersecting.
\[\begin{aligned}
 \rho \colon 3x - y - 4z + 2 = 0,\qquad \sigma \colon 6x - 2y - 8z + 5 = 0 & &
\end{aligned}\]}
{parallel, not identical}{identical}{intersecting}{}{}{}}
\LANGpl{
\Question{
Określ wzajemne położenie płaszczyzn.
\[\begin{aligned}
 \rho \colon 3x - y - 4z + 2 = 0,\qquad \sigma \colon 6x - 2y - 8z + 5 = 0 & &
\end{aligned}\]}
{równoległe, nie pokrywające się}{pokrywające się}{przecinające się}{}{}{}}
\LANGcs{
\Question{
Jsou dány roviny
\[\begin{aligned}
 \rho \colon 3x - y - 4z + 2 = 0,\quad \sigma \colon 6x - 2y - 8z + 5 = 0. & &
\end{aligned}\]
Určete jejich vzájemnou polohu.}
{Dané roviny jsou rovnoběžné různé.}{Dané roviny jsou totožné.}{Dané roviny jsou různoběžné.}{}{}{}}
\LANGsk{
\Question{
Sú dané roviny
\[\begin{aligned}
 \rho \colon 3x - y - 4z + 2 = 0,\quad \sigma \colon 6x - 2y - 8z + 5 = 0. & &
\end{aligned}\]
Určte ich vzájomnú polohu.}
{Dané roviny sú rovnobežné rôzne.}{Dané roviny sú totožné.}{Dané roviny sú rôznobežné.}{}{}{}}
\LANGes{
\Question{
Determina la posición relativa entre los planos.
\[\begin{aligned}
 \rho \colon 3x - y - 4z + 2 = 0,\qquad \sigma \colon 6x - 2y - 8z + 5 = 0 & &
\end{aligned}\]}
{Los planos son paralelos, no idénticos.}{Los planos son idénticos.}{Los planos no son paralelos.}{}{}{}}
\End

\Data\ID{9000117406}
\Updated{2020-07-15 03:07:37}
\Level{A}
\SubArea{Analytic geometry in a space}
\LANGen{
\Question{
Determine whether the following planes are parallel, identical or intersecting.
\[\begin{aligned}
 \rho \colon \frac{3}
{2}x -\frac{1} 
{4}y + \frac{2} 
{3}z -\frac{2} 
{5} = 0,\qquad \sigma \colon \frac{2}
{3}x - 4y + \frac{3} 
{2}z -\frac{5} 
{2} = 0 & &
\end{aligned}\]}
{intersecting}{identical}{parallel, not identical}{}{}{}}
\LANGpl{
\Question{
Określ wzajemne położenie płaszczyzn.
\[\begin{aligned}
 \rho \colon \frac{3}
{2}x -\frac{1} 
{4}y + \frac{2} 
{3}z -\frac{2} 
{5} = 0,\qquad \sigma \colon \frac{2}
{3}x - 4y + \frac{3} 
{2}z -\frac{5} 
{2} = 0 & &
\end{aligned}\]}
{przecinające się}{pokrywające się}{równoległe, nie pokrywające się}{}{}{}}
\LANGcs{
\Question{
Jsou dány roviny
\[\begin{aligned}
 \rho \colon \frac{3}
{2}x -\frac{1} 
{4}y + \frac{2} 
{3}z -\frac{2} 
{5} = 0,\quad \sigma \colon \frac{2}
{3}x - 4y + \frac{3} 
{2}z -\frac{5} 
{2} = 0. & &
\end{aligned}\]
Určete jejich vzájemnou polohu.}
{Dané roviny jsou různoběžné.}{Dané roviny jsou totožné.}{Dané roviny jsou rovnoběžné různé.}{}{}{}}
\LANGsk{
\Question{
Sú dané roviny
\[\begin{aligned}
 \rho \colon \frac{3}
{2}x -\frac{1} 
{4}y + \frac{2} 
{3}z -\frac{2} 
{5} = 0,\quad \sigma \colon \frac{2}
{3}x - 4y + \frac{3} 
{2}z -\frac{5} 
{2} = 0. & &
\end{aligned}\]
Určte ich vzájomnú polohu.}
{Dané roviny sú rôznobežné.}{Dané roviny sú totožné.}{Dané roviny sú rovnobežné rôzne.}{}{}{}}
\LANGes{
\Question{
Determina la posición relativa entre los planos.
\[\begin{aligned}
 \rho \colon \frac{3}
{2}x -\frac{1} 
{4}y + \frac{2} 
{3}z -\frac{2} 
{5} = 0,\qquad \sigma \colon \frac{2}
{3}x - 4y + \frac{3} 
{2}z -\frac{5} 
{2} = 0 & &
\end{aligned}\]}
{Los planos no son paralelos.}{Los planos son idénticos.}{Los planos son paralelos, no idénticos.}{}{}{}}
\End

\Data\ID{1003189005}
\Updated{2020-07-15 04:07:40}
\Level{B}
\SubArea{Analytic geometry in a space}
\LANGen{
\Question{
We are given a straight line \( p \) by parametric equations
\begin{align*}
x&=1+t, \\
y&= 1+2t, \\ 
z&= 4-t;\ t\in\mathbb{R}.
\end{align*}
Find the parametric equations of the line \( p' \) that is an orthogonal projection of the line \( p \) into the coordinate \(xy\)-plane .}
{$\begin{aligned}
p'\colon x&=5+s, \\
      y&= 9+2s, \\
      z&= 0;\ s\in\mathbb{R}
\end{aligned}$}{$\begin{aligned}
p'\colon x&=5+s, \\
y&= 9-2s, \\ 
z&=0;\ s\in\mathbb{R}
\end{aligned}$}{$\begin{aligned}
p'\colon x&=1+s, \\
y&=1+2s, \\
z&= 4;\ s\in\mathbb{R} 
\end{aligned}$}{$\begin{aligned}
p'\colon x&=5+2s, \\
y&=9+s, \\ 
z&= 0;\ s\in\mathbb{R} 
\end{aligned}$}{}{}}
\LANGpl{
\Question{
Dana jest prosta \( p \) określona równaniem parametrycznym
\begin{align*}
x&=1+t, \\
y&= 1+2t, \\ 
z&= 4-t;\ t\in\mathbb{R}.
\end{align*}
Wskaż równanie parametryczne prostej \( p' \), która jest rzutem prostopadłym prostej \( p \) na płaszczyźnie w układzie współrzędnych  \(xy\) .}
{$\begin{aligned}
p'\colon x&=5+s, \\
      y&= 9+2s, \\
      z&= 0;\ s\in\mathbb{R}
\end{aligned}$}{$\begin{aligned}
p'\colon x&=5+s, \\
y&= 9-2s, \\ 
z&=0;\ s\in\mathbb{R}
\end{aligned}$}{$\begin{aligned}
p'\colon x&=1+s, \\
y&=1+2s, \\
z&= 4;\ s\in\mathbb{R} 
\end{aligned}$}{$\begin{aligned}
p'\colon x&=5+2s, \\
y&=9+s, \\ 
z&= 0;\ s\in\mathbb{R} 
\end{aligned}$}{}{}}
\LANGcs{
\Question{
Přímka \( p \) je dána parametrickými rovnicemi:
\begin{align*}
x&=1+t, \\
y&= 1+2t, \\ 
z&= 4-t;\ t\in\mathbb{R}.
\end{align*}
Určete parametrické rovnice přímky \( p' \), která je kolmým průmětem přímky \( p \) do souřadné roviny \((xy)\) .}
{$\begin{aligned}
p'\colon x&=5+s, \\
      y&= 9+2s, \\
      z&= 0;\ s\in\mathbb{R}
\end{aligned}$}{$\begin{aligned}
p'\colon x&=5+s, \\
y&= 9-2s, \\ 
z&=0;\ s\in\mathbb{R}
\end{aligned}$}{$\begin{aligned}
p'\colon x&=1+s, \\
y&=1+2s, \\
z&= 4;\ s\in\mathbb{R} 
\end{aligned}$}{$\begin{aligned}
p'\colon x&=5+2s, \\
y&=9+s, \\ 
z&= 0;\ s\in\mathbb{R} 
\end{aligned}$}{}{}}
\LANGsk{
\Question{
Priamka \( p \) ja daná parametrickými rovnicami:
\begin{align*}
x&=1+t, \\
y&= 1+2t, \\ 
z&= 4-t;\ t\in\mathbb{R}.
\end{align*}
Určte parametrické rovnice priamky \( p' \), ktorá je kolmým priemetom priamky \( p \) do súradnicovej roviny \((xy)\) .}
{$\begin{aligned}
p'\colon x&=5+s, \\
      y&= 9+2s, \\
      z&= 0;\ s\in\mathbb{R}
\end{aligned}$}{$\begin{aligned}
p'\colon x&=5+s, \\
y&= 9-2s, \\ 
z&=0;\ s\in\mathbb{R}
\end{aligned}$}{$\begin{aligned}
p'\colon x&=1+s, \\
y&=1+2s, \\
z&= 4;\ s\in\mathbb{R} 
\end{aligned}$}{$\begin{aligned}
p'\colon x&=5+2s, \\
y&=9+s, \\ 
z&= 0;\ s\in\mathbb{R} 
\end{aligned}$}{}{}}
\LANGes{
\Question{
Dada la recta \( p \) cuyas ecuaciones paramétricas son:
\begin{align*}
x&=1+t, \\
y&= 1+2t, \\ 
z&= 4-t;\ t\in\mathbb{R}.
\end{align*}
Determina las ecuaciones paramétricas de la recta \( p' \) que es proyección perpendicular de la recta \( p \) en las coordenadas \(xy\).}
{$\begin{aligned}
p'\colon x&=5+s, \\
      y&= 9+2s, \\
      z&= 0;\ s\in\mathbb{R}
\end{aligned}$}{$\begin{aligned}
p'\colon x&=5+s, \\
y&= 9-2s, \\ 
z&=0;\ s\in\mathbb{R}
\end{aligned}$}{$\begin{aligned}
p'\colon x&=1+s, \\
y&=1+2s, \\
z&= 4;\ s\in\mathbb{R} 
\end{aligned}$}{$\begin{aligned}
p'\colon x&=5+2s, \\
y&=9+s, \\ 
z&= 0;\ s\in\mathbb{R} 
\end{aligned}$}{}{}}
\End

\Data\ID{1103189001}
\Updated{2020-07-15 03:07:21}
\Level{B}
\SubArea{Analytic geometry in a space}
\IMG{1103189001_0.pdf}
\LANGen{
\Question{
Find the general form of the equation of the plane \( \alpha \) that is perpendicular to the straight line \( p \) given by:
 \begin{align*}
x&=7+t, \\
y&=2t, \\ 
z&=4-t;\ t\in\mathbb{R},
\end{align*} 
and passes through the point \( A=[1;0;4] \). Consequently, find the coordinates of the point \( B \) which is the point of intersection of \( p \) and \( \alpha \) (see the picture).}
{\( \alpha\colon x+2y-z+3=0;\ B=[6;-2;5] \)}{\( \alpha\colon x+2y-z-3;\ B=[6;-2;5] \)}{\( \alpha\colon x+2y-z-3=0;\ B=[8;2;3] \)}{\( \alpha\colon x+2y-z+3=0;\ B=[8;2;3] \)}{}{}}
\LANGpl{
\Question{
Wskaż równanie ogólne płaszczyzny  \( \alpha \) tak, aby była prostopadła do prostej \( p \) określonej przez:
\begin{align*}
x&=7+t, \\
y&=2t, \\ 
z&=4-t;\ t\in\mathbb{R},
\end{align*}
oraz przechodziła przez punkt \( A=[1;0;4] \). Wskaż również współrzędne punktu  \( B \), który jest punktem przecięcia \( p \) oraz \( \alpha \) (spójrz na rysunek).}
{\( \alpha\colon x+2y-z+3=0;\ B=[6;-2;5] \)}{\( \alpha\colon x+2y-z-3;\ B=[6;-2;5] \)}{\( \alpha\colon x+2y-z-3=0;\ B=[8;2;3] \)}{\( \alpha\colon x+2y-z+3=0;\ B=[8;2;3] \)}{}{}}
\LANGcs{
\Question{
Určete obecnou rovnici roviny \( \alpha \), která je kolmá k přímce \( p \):
\begin{align*}
x&=7+t, \\
y&=2t, \\ 
z&=4-t;\ t\in\mathbb{R},
\end{align*}
a prochází bodem \( A=[1;0;4] \). Dále vypočtěte souřadnice bodu \( B \), ve kterém přímka \( p \) protíná rovinu \( \alpha \) (viz obrázek).}
{\( \alpha\colon x+2y-z+3=0;\ B=[6;-2;5] \)}{\( \alpha\colon x+2y-z-3;\ B=[6;-2;5] \)}{\( \alpha\colon x+2y-z-3=0;\ B=[8;2;3] \)}{\( \alpha\colon x+2y-z+3=0;\ B=[8;2;3] \)}{}{}}
\LANGsk{
\Question{
Určte všeobecnú rovnicu roviny \( \alpha \), ktorá je kolmá na priamku \( p \):
\begin{align*}
x&=7+t, \\
y&=2t, \\ 
z&=4-t;\ t\in\mathbb{R},
\end{align*}
a prechádza bodom \( A=[1;0;4] \). Ďalej vypočítajte súradnice bodu \( B \), v ktorom priamka \( p \) pretína rovinu \( \alpha \) (viď obrázok).}
{\( \alpha\colon x+2y-z+3=0;\ B=[6;-2;5] \)}{\( \alpha\colon x+2y-z-3;\ B=[6;-2;5] \)}{\( \alpha\colon x+2y-z-3=0;\ B=[8;2;3] \)}{\( \alpha\colon x+2y-z+3=0;\ B=[8;2;3] \)}{}{}}
\LANGes{
\Question{
Determina la ecuación general del plano \( \alpha \) que es perpendicular a la recta \( p \):
 \begin{align*}
x&=7+t, \\
y&=2t, \\ 
z&=4-t;\ t\in\mathbb{R},
\end{align*} 
y pasa por el punto \( A=[1;0;4] \). Luego cuenta las coordenadas del punto \( B \) en el que la recta \( p \) corta el plano \( \alpha \) (mira la imagen).}
{\( \alpha\colon x+2y-z+3=0;\ B=[6;-2;5] \)}{\( \alpha\colon x+2y-z-3;\ B=[6;-2;5] \)}{\( \alpha\colon x+2y-z-3=0;\ B=[8;2;3] \)}{\( \alpha\colon x+2y-z+3=0;\ B=[8;2;3] \)}{}{}}
\End

\Data\ID{1103189002}
\Updated{2020-07-15 03:07:21}
\Level{B}
\SubArea{Analytic geometry in a space}
\IMG{1103189002.pdf}
\LANGen{
\Question{
Find the general form of the equation of the plane \( \beta \) that passes through the points \( M=[-1;1;-3] \) and \( N=[0;2;-1] \) and is perpendicular to the plane \( \alpha \) given by \( 3x-y+2=0 \) (see the picture).}
{\( \beta\colon x+3y-2z-8=0 \)}{\( \beta\colon x+3z+10=0 \)}{\( \beta\colon x+3z+3=0 \)}{\( \beta\colon x+3y-2z+8=0 \)}{}{}}
\LANGpl{
\Question{
Wskaż równanie ogólne płaszczyzny  \( \beta \) przechodzącej prze punkty \( M=[-1;1;-3] \) i \( N=[0;2;-1] \) oraz prostopadłej do płaszczyzny\( \alpha \) określonej przez \( 3x-y+2=0 \) (spójrz na rysunek).}
{\( \beta\colon x+3y-2z-8=0 \)}{\( \beta\colon x+3z+10=0 \)}{\( \beta\colon x+3z+3=0 \)}{\( \beta\colon x+3y-2z+8=0 \)}{}{}}
\LANGcs{
\Question{
Určete obecnou rovnici roviny \( \beta \), která prochází body \( M=[-1;1;-3] \) a \( N=[0;2;-1] \) a je kolmá k rovině \( \alpha \): \( 3x-y+2=0 \) (viz obrázek).}
{\( \beta\colon x+3y-2z-8=0 \)}{\( \beta\colon x+3z+10=0 \)}{\( \beta\colon x+3z+3=0 \)}{\( \beta\colon x+3y-2z+8=0 \)}{}{}}
\LANGsk{
\Question{
Určte všeobecnú rovnicu roviny \( \beta \), ktorá prechádza bodmi \( M=[-1;1;-3] \) a \( N=[0;2;-1] \) a je kolmá na rovinu \( \alpha \): \( 3x-y+2=0 \) (viď obrázok).}
{\( \beta\colon x+3y-2z-8=0 \)}{\( \beta\colon x+3z+10=0 \)}{\( \beta\colon x+3z+3=0 \)}{\( \beta\colon x+3y-2z+8=0 \)}{}{}}
\LANGes{
\Question{
Determina la ecuación general del plano \( \beta \) el cual pasa por los puntos \( M=[-1;1;-3] \) y \( N=[0;2;-1] \) y es perpendicular al plano  \( \alpha \): \( 3x-y+2=0 \) (mira la imagen).}
{\( \beta\colon x+3y-2z-8=0 \)}{\( \beta\colon x+3z+10=0 \)}{\( \beta\colon x+3z+3=0 \)}{\( \beta\colon x+3y-2z+8=0 \)}{}{}}
\End

\Data\ID{1103189003}
\Updated{2020-07-15 03:07:21}
\Level{B}
\SubArea{Analytic geometry in a space}
\IMG{1103189003.pdf}
\LANGen{
\Question{
Find the general form of the equation of the plane \( \beta \) that passes through the straight line \( p \) given by parametric equations
\begin{align*}
x&=1+2t, \\
y&=-2t, \\ 
z&=1+t;\ t\in\mathbb{R},
\end{align*}
and is perpendicular to the plane \( \alpha \) given by \( x+3y-z-7=0 \) (see the picture).}
{\( \beta\colon x-3y-8z+7=0 \)}{\( \beta\colon 2x-2y+z-3=0 \)}{\( \beta\colon x-3y-8z-7=0 \)}{\( \beta\colon 2x-2y+z+3=0 \)}{}{}}
\LANGpl{
\Question{
Wskaż równanie ogólne płaszczyzny  \( \beta \) przechodzącej przez prostą \( p \) określoną równaniem parametrycznym
\begin{align*}
x&=1+2t, \\
y&=-2t, \\ 
z&=1+t;\ t\in\mathbb{R},
\end{align*}
oraz prostopadłą do płaszczyzny  \( \alpha \) określoną przez \( x+3y-z-7=0 \) (spójrz na rysunek).}
{\( \beta\colon x-3y-8z+7=0 \)}{\( \beta\colon 2x-2y+z-3=0 \)}{\( \beta\colon x-3y-8z-7=0 \)}{\( \beta\colon 2x-2y+z+3=0 \)}{}{}}
\LANGcs{
\Question{
Určete obecnou rovnici roviny \( \beta \), která prochází přímkou \( p \) danou parametrickými rovnicemi
\begin{align*}
x&=1+2t, \\
y&=-2t, \\ 
z&=1+t;\ t\in\mathbb{R},
\end{align*}
a je kolmá k rovině \( \alpha \): \( x+3y-z-7=0 \) (viz obrázek).}
{\( \beta\colon x-3y-8z+7=0 \)}{\( \beta\colon 2x-2y+z-3=0 \)}{\( \beta\colon x-3y-8z-7=0 \)}{\( \beta\colon 2x-2y+z+3=0 \)}{}{}}
\LANGsk{
\Question{
Určte všeobecnú rovnicu roviny \( \beta \), ktorá prechádza priamkou \( p \) danou parametrickými rovnicami
\begin{align*}
x&=1+2t, \\
y&=-2t, \\ 
z&=1+t;\ t\in\mathbb{R},
\end{align*} 
a je kolmá k rovine \( \alpha \): \( x+3y-z-7=0 \) (viď obrázok).}
{\( \beta\colon x-3y-8z+7=0 \)}{\( \beta\colon 2x-2y+z-3=0 \)}{\( \beta\colon x-3y-8z-7=0 \)}{\( \beta\colon 2x-2y+z+3=0 \)}{}{}}
\LANGes{
\Question{
Determina la ecuación general del plano \( \beta \) que pasa por la recta \( p \) dada por las ecuaciones paramétricas:
\begin{align*}
x&=1+2t, \\
y&=-2t, \\ 
z&=1+t;\ t\in\mathbb{R},
\end{align*}
y es perpendicular al plano \( \alpha \): \( x+3y-z-7=0 \) (mira la imagen).}
{\( \beta\colon x-3y-8z+7=0 \)}{\( \beta\colon 2x-2y+z-3=0 \)}{\( \beta\colon x-3y-8z-7=0 \)}{\( \beta\colon 2x-2y+z+3=0 \)}{}{}}
\End

\Data\ID{1103189004}
\Updated{2020-09-21 12:09:17}
\Level{B}
\SubArea{Analytic geometry in a space}
\IMG{1103189004_0.pdf}
\LANGen{
\Question{
We are given the point \( A=[2;-1;-4] \) and planes \( \rho \) by \( x-y+3z-5=0 \) and \( \sigma \) by \( 2x-y-z-8=0 \). Find the general form of the equation of the plane \( \alpha \) which passes through the point \( A \) and is perpendicular to both planes (see the picture).}
{\( \alpha\colon 4x+7y+z+3=0 \)}{\( \alpha\colon -2x+5y-3z-3=0 \)}{\( \alpha\colon 4x-7y+z+3=0 \)}{\( \alpha\colon 2x-5y+3z+3=0 \)}{}{}}
\LANGpl{
\Question{
Dany jest punkt \( A=[2;-1;-4] \) oraz płaszczyzny \( \rho \) określona przez \( x-y+3z-5=0 \) i \( \sigma \) określona przez \( 2x-y-z-8=0 \). Wskaż równanie ogólne płaszczyzny  \( \alpha \) przechodzącej przez punkt \( A \) oraz prostopadłej do obu płaszczyzn  (spójrz na rysunek).}
{\( \alpha\colon 4x+7y+z+3=0 \)}{\( \alpha\colon -2x+5y-3z-3=0 \)}{\( \alpha\colon 4x-7y+z+3=0 \)}{\( \alpha\colon 2x-5y+3z+3=0 \)}{}{}}
\LANGcs{
\Question{
Je dán bod \( A=[2;-1;-4] \) a roviny \( \rho \): \( x-y+3z-5=0 \) a \( \sigma \): \( 2x-y-z-8=0 \). Určete obecnou rovnici roviny \( \alpha \), která prochází bodem \( A \) a je kolmá k oběma daným rovinám (viz obrázek).}
{\( \alpha\colon 4x+7y+z+3=0 \)}{\( \alpha\colon -2x+5y-3z-3=0 \)}{\( \alpha\colon 4x-7y+z+3=0 \)}{\( \alpha\colon 2x-5y+3z+3=0 \)}{}{}}
\LANGsk{
\Question{
Je daný bod \( A=[2;-1;-4] \) a roviny \( \rho \): \( x-y+3z-5=0 \) a \( \sigma \): \( 2x-y-z-8=0 \). Určte všeobecnú rovnicu roviny \( \alpha \), ktorá prechádza bodom \( A \) a je kolmá k obom daným rovinám (viď obrázok).}
{\( \alpha\colon 4x+7y+z+3=0 \)}{\( \alpha\colon -2x+5y-3z-3=0 \)}{\( \alpha\colon 4x-7y+z+3=0 \)}{\( \alpha\colon 2x-5y+3z+3=0 \)}{}{}}
\LANGes{
\Question{
Dado el punto \( A=[2;-1;-4] \) y los planos \( \rho \): \( x-y+3z-5=0 \) y \( \sigma \): \( 2x-y-z-8=0 \). Determina la ecuación general del plano \( \alpha \) que pasa por el punto \( A \) y es perpendicular a los dos planos dados (mira la imagen).}
{\( \alpha\colon 4x+7y+z+3=0 \)}{\( \alpha\colon -2x+5y-3z-3=0 \)}{\( \alpha\colon 4x-7y+z+3=0 \)}{\( \alpha\colon 2x-5y+3z+3=0 \)}{}{}}
\End

\Data\ID{9000101101}
\Updated{2020-07-15 03:07:37}
\Level{B}
\SubArea{Analytic geometry in a space}
\LANGen{
\Question{
Find the distance from the point \(A = [0;1;-3]\) 
to the point \(B = [-1;3;0]\).}
{\(\sqrt{14}\)}{\(3\)}{\(4\)}{\(\sqrt{26}\)}{}{}}
\LANGpl{
\Question{
Wyznacz odległość z punktu \(A = [0;1;-3]\) 
do  punktu \(B = [-1;3;0]\).}
{\(\sqrt{14}\)}{\(3\)}{\(4\)}{\(\sqrt{26}\)}{}{}}
\LANGcs{
\Question{
Jsou dány body \(A = [0;1;-3]\) 
a \(B = [-1;3;0]\).
Vypočítejte jejich vzdálenost.}
{\(\sqrt{14}\)}{\(3\)}{\(4\)}{\(\sqrt{26}\)}{}{}}
\LANGsk{
\Question{
Sú dané body \(A = [0;1;-3]\) 
a \(B = [-1;3;0]\).
Vypočítajte ich vzdialenosť.}
{\(\sqrt{14}\)}{\(3\)}{\(4\)}{\(\sqrt{26}\)}{}{}}
\LANGes{
\Question{
Calcula la distancia relativa entre el punto \(A = [0;1;-3]\) 
y el punto \(B = [-1;3;0]\).}
{\(\sqrt{14}\)}{\(3\)}{\(4\)}{\(\sqrt{26}\)}{}{}}
\End

\Data\ID{9000101102}
\Updated{2020-07-15 03:07:37}
\Level{B}
\SubArea{Analytic geometry in a space}
\LANGen{
\Question{
Find the distance between the point \(A = [1;0;1]\) 
and the line \(p\).
\[ 
\begin{aligned}p\colon x& = 2, &
\\y & = 3t,
\\z & = 1 - t;\ t\in \mathbb{R}
\\ \end{aligned}
\]}
{\(1\)}{\(0\)}{\(2\)}{\(3\)}{}{}}
\LANGpl{
\Question{
Wyznacz odległość pomiędzy punktem \(A = [1;0;1]\) 
a prostą \(p\).
\[ 
\begin{aligned}p\colon x& = 2, &
\\y & = 3t,
\\z & = 1 - t;\ t\in \mathbb{R}
\\ \end{aligned}
\]}
{\(1\)}{\(0\)}{\(2\)}{\(3\)}{}{}}
\LANGcs{
\Question{
Je dán bod \(A = [1;0;1]\) 
a přímka \(p\colon x = 2;y = 3t;z = 1 - t\),
\(t\in \mathbb{R}\). Vypočítejte
vzdálenost bodu \(A\) 
od přímky \(p\).}
{\(1\)}{\(0\)}{\(2\)}{\(3\)}{}{}}
\LANGsk{
\Question{
Je daný bod \(A = [1;0;1]\) 
a priamka \(p\colon x = 2;y = 3t;z = 1 - t\),
\(t\in \mathbb{R}\). Vypočítajte
vzdialenosť bodu \(A\) 
od priamky \(p\).}
{\(1\)}{\(0\)}{\(2\)}{\(3\)}{}{}}
\LANGes{
\Question{
Encuentra la distancia relativa entre el punto \(A = [1;0;1]\) 
y la recta \(p\).
\[ 
\begin{aligned}p\colon x& = 2, &
\\y & = 3t,
\\z & = 1 - t;\ t\in \mathbb{R}
\\ \end{aligned}
\]}
{\(1\)}{\(0\)}{\(2\)}{\(3\)}{}{}}
\End

\Data\ID{9000101103}
\Updated{2020-07-15 03:07:37}
\Level{B}
\SubArea{Analytic geometry in a space}
\LANGen{
\Question{
Find the distance between two parallel lines
\(p\) and
\(q\).
\[ 
\begin{aligned}p\colon x& = 2, &
\\y & = 3t,
\\z & = 1 - t;\  t\in \mathbb{R}
\\ \end{aligned}\qquad \qquad \begin{aligned}q\colon x& = 3, &
\\y & = 6s,
\\z & = 1 - 2s;\ s\in \mathbb{R}
\\ \end{aligned}
\]}
{\(1\)}{\(2\)}{\(3\)}{\(4\)}{}{}}
\LANGpl{
\Question{
Wyznacz odległość pomiędzy dwoma prostymi równoległymi
\(p\) i
\(q\).
\[ 
\begin{aligned}p\colon x& = 2, &
\\y & = 3t,
\\z & = 1 - t;\  t\in \mathbb{R}
\\ \end{aligned}\qquad \qquad \begin{aligned}q\colon x& = 3, &
\\y & = 6s,
\\z & = 1 - 2s;\ s\in \mathbb{R}
\\ \end{aligned}
\]}
{\(1\)}{\(2\)}{\(3\)}{\(4\)}{}{}}
\LANGcs{
\Question{
Jsou dány dvě rovnoběžné přímky
\(p\colon x = 2;y = 3t;z = 1 - t\),
\(t\in \mathbb{R}\),
\(q\colon x = 3;y = 6s;z = 1 - 2s\),
\(s\in \mathbb{R}\).
Vypočítejte jejich vzdálenost.}
{\(1\)}{\(2\)}{\(3\)}{\(4\)}{}{}}
\LANGsk{
\Question{
Sú dané dve rovnobežné priamky
\(p\colon x = 2;y = 3t;z = 1 - t\),
\(t\in \mathbb{R}\),
\(q\colon x = 3;y = 6s;z = 1 - 2s\),
\(s\in \mathbb{R}\).
Vypočítajte ich vzdialenosť.}
{\(1\)}{\(2\)}{\(3\)}{\(4\)}{}{}}
\LANGes{
\Question{
Halla la distancia entre dos rectas paralelas \(p\) y
\(q\).
\[ 
\begin{aligned}p\colon x& = 2, &
\\y & = 3t,
\\z & = 1 - t;\  t\in \mathbb{R}
\\ \end{aligned}\qquad \qquad \begin{aligned}q\colon x& = 3, &
\\y & = 6s,
\\z & = 1 - 2s;\ s\in \mathbb{R}
\\ \end{aligned}
\]}
{\(1\)}{\(2\)}{\(3\)}{\(4\)}{}{}}
\End

\Data\ID{9000101104}
\Updated{2020-07-15 03:07:37}
\Level{B}
\SubArea{Analytic geometry in a space}
\LANGen{
\Question{
Find the distance between the point \(A = [-1;1;0]\) 
and the plane \(\alpha \).
\[ 
 \alpha \colon 12y + 5z - 2 = 0
\]}
{\(\frac{10}
{13}\)}{\(\frac{15}
{13}\)}{\(\frac{17}
{13}\)}{\(\frac{14}
{13}\)}{}{}}
\LANGpl{
\Question{
Określ odległość pomiędzy punktem \(A = [-1;1;0]\), 
a płaszczyzną  \(\alpha \).
\[ 
 \alpha \colon 12y + 5z - 2 = 0
\]}
{\(\frac{10}
{13}\)}{\(\frac{15}
{13}\)}{\(\frac{17}
{13}\)}{\(\frac{14}
{13}\)}{}{}}
\LANGcs{
\Question{
Je dán bod \(A = [-1;1;0]\) a
rovina \(\alpha \colon 12y + 5z - 2 = 0\). Určete
vzdálenost bodu \(A\) 
od roviny \(\alpha \).}
{\(\frac{10}
{13}\)}{\(\frac{15}
{13}\)}{\(\frac{17}
{13}\)}{\(\frac{14}
{13}\)}{}{}}
\LANGsk{
\Question{
Je daný bod \(A = [-1;1;0]\) a
rovina \(\alpha \colon 12y + 5z - 2 = 0\). Určte
vzdialenosť bodu \(A\) 
od roviny \(\alpha \).}
{\(\frac{10}
{13}\)}{\(\frac{15}
{13}\)}{\(\frac{17}
{13}\)}{\(\frac{14}
{13}\)}{}{}}
\LANGes{
\Question{
Halla la distancia entre el punto \(A = [-1;1;0]\) y el plano \(\alpha \).
\[ 
 \alpha \colon 12y + 5z - 2 = 0
\]}
{\(\frac{10}
{13}\)}{\(\frac{15}
{13}\)}{\(\frac{17}
{13}\)}{\(\frac{14}
{13}\)}{}{}}
\End

\Data\ID{9000101105}
\Updated{2020-07-15 03:07:37}
\Level{B}
\SubArea{Analytic geometry in a space}
\LANGen{
\Question{
Find the distance between two parallel planes
\(\alpha \) and
\(\beta \).
\[ \begin{aligned}
 \alpha& \colon x - 1.5y - 3z - 1 = 0,\\ \beta &\colon 2x - 3y - 6z + 3 = 0
\end{aligned}\]}
{\(\frac{5}
{7}\)}{\(\frac{1}
{7}\)}{\(\frac{2}
{49}\)}{\(\frac{2}
{41}\)}{}{}}
\LANGpl{
\Question{
Określ odległość pomiędzy dwoma równoległymi płaszczyznami 
\(\alpha \) i
\(\beta \).
\[ \begin{aligned}
 \alpha& \colon x - 1{,}5y - 3z - 1 = 0,\\ \beta &\colon 2x - 3y - 6z + 3 = 0
\end{aligned}\]}
{\(\frac{5}
{7}\)}{\(\frac{1}
{7}\)}{\(\frac{2}
{49}\)}{\(\frac{2}
{41}\)}{}{}}
\LANGcs{
\Question{
Jsou dány dvě rovnoběžné roviny
\(\alpha \colon x - 1{,}5y - 3z - 1 = 0\),
\(\beta \colon 2x - 3y - 6z + 3 = 0\).
Určete vzdálenost těchto rovin.}
{\(\frac{5}
{7}\)}{\(\frac{1}
{7}\)}{\(\frac{2}
{49}\)}{\(\frac{2}
{41}\)}{}{}}
\LANGsk{
\Question{
Sú dané dve rovnobežné roviny
\(\alpha \colon x - 1{,}5y - 3z - 1 = 0\),
\(\beta \colon 2x - 3y - 6z + 3 = 0\).
Určte vzdialenosť týchto rovín.}
{\(\frac{5}
{7}\)}{\(\frac{1}
{7}\)}{\(\frac{2}
{49}\)}{\(\frac{2}
{41}\)}{}{}}
\LANGes{
\Question{
Halla la distancia entre dos planos paralelos 
\(\alpha \) y
\(\beta \).
\[ \begin{aligned}
 \alpha& \colon x - 1.5y - 3z - 1 = 0,\\ \beta &\colon 2x - 3y - 6z + 3 = 0
\end{aligned}\]}
{\(\frac{5}
{7}\)}{\(\frac{1}
{7}\)}{\(\frac{2}
{49}\)}{\(\frac{2}
{41}\)}{}{}}
\End

\Data\ID{9000101106}
\Updated{2020-07-15 03:07:37}
\Level{B}
\SubArea{Analytic geometry in a space}
\LANGen{
\Question{
In the following list identify a point which does not have zero distance from the line
\(m\).
\[ 
\begin{aligned}m\colon x& = s, &
\\y & = 8 - s,
\\z & = 1 + 3s;\  s\in \mathbb{R}
\\ \end{aligned}
\]}
{\([2;6;10]\)}{\([0;8;1]\)}{\([1;7;4]\)}{\([8;0;25]\)}{}{}}
\LANGpl{
\Question{
Wyznacz punkt, tak, aby  odległość punktu od prostej 
\(m\) nie miała wartości zerowej.
\[ 
\begin{aligned}m\colon x& = s, &
\\y & = 8 - s,
\\z & = 1 + 3s;\  s\in \mathbb{R}
\\ \end{aligned}
\]}
{\([2;6;10]\)}{\([0;8;1]\)}{\([1;7;4]\)}{\([8;0;25]\)}{}{}}
\LANGcs{
\Question{
Je dána přímka \(m\colon x = s;y = 8 - s;z = 1 + 3s\),
\(s\in \mathbb{R}\). Vyberte bod, který
nemá od přímky \(m\) 
vzdálenost \(v = 0\).}
{\([2;6;10]\)}{\([0;8;1]\)}{\([1;7;4]\)}{\([8;0;25]\)}{}{}}
\LANGsk{
\Question{
Je daná priamka \(m\colon x = s;y = 8 - s;z = 1 + 3s\),
\(s\in \mathbb{R}\). Vyberte bod, ktorý
nemá od priamky \(m\) 
vzdialenosť \(v = 0\).}
{\([2;6;10]\)}{\([0;8;1]\)}{\([1;7;4]\)}{\([8;0;25]\)}{}{}}
\LANGes{
\Question{
Dada la recta \(m\colon x = s;y = 8 - s;z = 1 + 3s\),
\(s\in \mathbb{R}\). Elige el punto cuya distancia de la recta \(m\) 
no es \(v = 0\).}
{\([2;6;10]\)}{\([0;8;1]\)}{\([1;7;4]\)}{\([8;0;25]\)}{}{}}
\End

\Data\ID{9000101107}
\Updated{2020-07-15 03:07:37}
\Level{B}
\SubArea{Analytic geometry in a space}
\LANGen{
\Question{
Find the distance between the line \(p\) 
and the plane \(\alpha \).
\[ 
\alpha \colon x-3y+2z-4 = 0,\qquad \qquad \begin{aligned}[t] p\colon x& = 1 + t, &
\\y & = -3t,
\\z & = 2;\ t\in \mathbb{R}
\\ \end{aligned}
\]}
{\(0\)}{\(\frac{5}
{\sqrt{17}}\)}{\(2\)}{\(1\)}{}{}}
\LANGpl{
\Question{
Określ odległość pomiędzy prostą \(p\), 
a płaszczyzną  \(\alpha \).
\[ 
\alpha \colon x-3y+2z-4 = 0,\qquad \qquad \begin{aligned}[t] p\colon x& = 1 + t, &
\\y & = -3t,
\\z & = 2;\ t\in \mathbb{R}
\\ \end{aligned}
\]}
{\(0\)}{\(\frac{5}
{\sqrt{17}}\)}{\(2\)}{\(1\)}{}{}}
\LANGcs{
\Question{
Vypočítejte vzdálenost přímky
\(p\) a roviny
\(\alpha \).
\[ 
\alpha \colon x-3y+2z-4 = 0,\qquad \qquad \begin{aligned}[t] p\colon x& = 1 + t, &
\\y & = -3t,
\\z & = 2;\ t\in \mathbb{R}
\\ \end{aligned}
\]}
{\(0\)}{\(\frac{5}
{\sqrt{17}}\)}{\(2\)}{\(1\)}{}{}}
\LANGsk{
\Question{
Vypočítajte vzdialenosť priamky
\(p\) a roviny
\(\alpha \).
\[ 
\alpha \colon x-3y+2z-4 = 0,\qquad \qquad \begin{aligned}[t] p\colon x& = 1 + t, &
\\y & = -3t,
\\z & = 2;\ t\in \mathbb{R}
\\ \end{aligned}
\]}
{\(0\)}{\(\frac{5}
{\sqrt{17}}\)}{\(2\)}{\(1\)}{}{}}
\LANGes{
\Question{
Halla la distancia entre la línea \(p\) y el plano \(\alpha \).
\[ 
\alpha \colon x-3y+2z-4 = 0,\qquad \qquad \begin{aligned}[t] p\colon x& = 1 + t, &
\\y & = -3t,
\\z & = 2;\ t\in \mathbb{R}
\\ \end{aligned}
\]}
{\(0\)}{\(\frac{5}
{\sqrt{17}}\)}{\(2\)}{\(1\)}{}{}}
\End

\Data\ID{9000101108}
\Updated{2020-07-15 04:07:14}
\Level{B}
\SubArea{Analytic geometry in a space}
\LANGen{
\Question{
Find the distance between the line \(q\) 
and the plane \(\beta \).
\[ 
\beta \colon x+4y+2z-4 = 0,\qquad \qquad \begin{aligned}[t] q\colon x& = 4, &
\\y & = -2t,
\\z & = 1 + 4t;\ t\in \mathbb{R}
\\ \end{aligned}
\]}
{\(\frac{2}
{\sqrt{21}}\)}{\(\frac{4}
{\sqrt{21}}\)}{\(0\)}{\(1\)}{}{}}
\LANGpl{
\Question{
Określ odległość pomiędzy prostą  \(q\), 
a płaszczyzną  \(\beta \).
\[ 
\beta \colon x+4y+2z-4 = 0,\qquad \qquad \begin{aligned}[t] q\colon x& = 4, &
\\y & = -2t,
\\z & = 1 + 4t;\ t\in \mathbb{R}
\\ \end{aligned}
\]}
{\(\frac{2}
{\sqrt{21}}\)}{\(\frac{4}
{\sqrt{21}}\)}{\(0\)}{\(1\)}{}{}}
\LANGcs{
\Question{
Vypočítejte vzdálenost přímky
\(q\) a roviny
\(\beta \). \[ 
\beta \colon x+4y+2z-4 = 0,\qquad \qquad \begin{aligned}[t] q\colon x& = 4, &
\\y & = -2t,
\\z & = 1 + 4t;\ t\in \mathbb{R}
\\ \end{aligned}
\]}
{\(\frac{2}
{\sqrt{21}}\)}{\(\frac{4}
{\sqrt{21}}\)}{\(0\)}{\(1\)}{}{}}
\LANGsk{
\Question{
Vypočítajte vzdialenosť priamky
\(q\) a roviny
\(\beta \).\[ 
\beta \colon x+4y+2z-4 = 0,\qquad \qquad \begin{aligned}[t] q\colon x& = 4, &
\\y & = -2t,
\\z & = 1 + 4t;\ t\in \mathbb{R}
\\ \end{aligned}
\]}
{\(\frac{2}
{\sqrt{21}}\)}{\(\frac{4}
{\sqrt{21}}\)}{\(0\)}{\(1\)}{}{}}
\LANGes{
\Question{
Halla la distancia entre la recta \(q\) y el plano \(\beta \).
\[ 
\beta \colon x+4y+2z-4 = 0,\qquad \qquad \begin{aligned}[t] q\colon x& = 4, &
\\y & = -2t,
\\z & = 1 + 4t;\ t\in \mathbb{R}
\\ \end{aligned}
\]}
{\(\frac{2}
{\sqrt{21}}\)}{\(\frac{4}
{\sqrt{21}}\)}{\(0\)}{\(1\)}{}{}}
\End

\Data\ID{9000101109}
\Updated{2020-07-15 03:07:37}
\Level{B}
\SubArea{Analytic geometry in a space}
\IMG{001011_09-figure0.pdf}
\LANGen{
\Question{
The points \(A = [0;5;0]\),
\(B = [5;5;0]\),
\(C = [5;0;0]\),
\(D = [0;0;0]\) define a cube
\(ABCDEFGH\). Find the distance
between the line \(AB\) 
and the plane \(EFG\).}
{\(5\)}{\(3\)}{\(4\)}{\(6\)}{}{}}
\LANGpl{
\Question{
Punkty \(A = [0;5;0]\),
\(B = [5;5;0]\),
\(C = [5;0;0]\),
\(D = [0;0;0]\) tworzą sześcian
\(ABCDEFGH\). Wyznacz odległość
pomiędzy prostą  \(AB\), 
a płaszczyzną  \(EFG\).}
{\(5\)}{\(3\)}{\(4\)}{\(6\)}{}{}}
\LANGcs{
\Question{
Jsou dány body \(A = [0;5;0]\),
\(B = [5;5;0]\),
\(C = [5;0;0]\),
\(D = [0;0;0]\), které tvoří vrcholy
krychle \(ABCDEFGH\). Vypočítejte
vzdálenost přímky \(AB\) 
od roviny \(EFG\).}
{\(5\)}{\(3\)}{\(4\)}{\(6\)}{}{}}
\LANGsk{
\Question{
Sú dané body \(A = [0;5;0]\),
\(B = [5;5;0]\),
\(C = [5;0;0]\),
\(D = [0;0;0]\), ktoré tvoria vrcholy
kocky \(ABCDEFGH\). Vypočítajte
vzdialenosť priamky \(AB\) 
od roviny \(EFG\).}
{\(5\)}{\(3\)}{\(4\)}{\(6\)}{}{}}
\LANGes{
\Question{
Dados los puntos \(A = [0;5;0]\),
\(B = [5;5;0]\),
\(C = [5;0;0]\),
\(D = [0;0;0]\) que definen el cubo
\(ABCDEFGH\). Halla la distancia 
entre la recta \(AB\) 
y el plano \(EFG\).}
{\(5\)}{\(3\)}{\(4\)}{\(6\)}{}{}}
\End

\Data\ID{9000101110}
\Updated{2020-07-15 03:07:37}
\Level{B}
\SubArea{Analytic geometry in a space}
\IMG{001011_10-figure0.pdf}
\LANGen{
\Question{
The points \(A = [0;5;0]\),
\(B = [5;5;0]\),
\(C = [5;0;0]\),
\(D = [0;0;0]\) define a cube
\(ABCDEFGH\). Find the distance
between the point \(A\) 
and the point \(F\).}
{\(\sqrt{50}\)}{\(\sqrt{5}\)}{\(5\)}{\(10\)}{}{}}
\LANGpl{
\Question{
Punkty \(A = [0;5;0]\),
\(B = [5;5;0]\),
\(C = [5;0;0]\),
\(D = [0;0;0]\) tworzą sześcian
\(ABCDEFGH\). Określ odległość pomiędzy punktem \(A\), 
a punktem  \(F\).}
{\(\sqrt{50}\)}{\(\sqrt{5}\)}{\(5\)}{\(10\)}{}{}}
\LANGcs{
\Question{
Jsou dány body \(A = [0;5;0]\),
\(B = [5;5;0]\),
\(C = [5;0;0]\),
\(D = [0;0;0]\), které tvoří vrcholy
krychle \(ABCDEFGH\). Vypočítejte
vzdálenost bodu \(A\) 
od bodu \(F\).}
{\(\sqrt{50}\)}{\(\sqrt{5}\)}{\(5\)}{\(10\)}{}{}}
\LANGsk{
\Question{
Sú dané body \(A = [0;5;0]\),
\(B = [5;5;0]\),
\(C = [5;0;0]\),
\(D = [0;0;0]\), ktoré tvoria vrcholy
kocky \(ABCDEFGH\). Vypočítajte
vzdialenosť bodu \(A\) 
od bodu \(F\).}
{\(\sqrt{50}\)}{\(\sqrt{5}\)}{\(5\)}{\(10\)}{}{}}
\LANGes{
\Question{
Dados los puntos \(A = [0;5;0]\),
\(B = [5;5;0]\),
\(C = [5;0;0]\),
\(D = [0;0;0]\) que definen el cubo
\(ABCDEFGH\). Halla la distancia
entre el punto \(A\) 
y el punto \(F\).}
{\(\sqrt{50}\)}{\(\sqrt{5}\)}{\(5\)}{\(10\)}{}{}}
\End

\Data\ID{9000101901}
\Updated{2020-07-15 03:07:37}
\Level{B}
\SubArea{Analytic geometry in a space}
\LANGen{
\Question{
Find the angle between two lines and round your answer to the nearest minute. 
\[ 
\begin{aligned}p\colon x& = 2 - t , & 
\\y & = 3t ,
\\z & = 1 ;\ t\in \mathbb{R}
\\ \end{aligned}\qquad \qquad \begin{aligned}q\colon x& = 2s, & 
\\y & = 4s ,
\\z & = 1 - s;\ s\in \mathbb{R}
\\ \end{aligned}
\]}
{\(46^{\circ }22'\)}{\(0^{\circ }\)}{\(67^{\circ }18'\)}{\(90^{\circ }\)}{}{}}
\LANGpl{
\Question{
Wyznacz kąt pomiędzy dwoma prostymi w przestrzeni i zaokrągli do pełnych stopni i minut. 
\[ 
\begin{aligned}p\colon x& = 2 - t,&
\\y & = 3t,
\\z & = 1;\ t\in \mathbb{R}
\\ \end{aligned}\qquad \qquad \begin{aligned}q\colon x& = 2s, &
\\y & = 4s,
\\z & = 1 - s;\ s\in \mathbb{R}
\\ \end{aligned}
\]}
{\(46^{\circ }22'\)}{\(0^{\circ }\)}{\(67^{\circ }18'\)}{\(90^{\circ }\)}{}{}}
\LANGcs{
\Question{
Určete odchylku přímek \(p\),
\(q\),
kde
\[ \begin{alignedat}{80}
 p\colon &x = 2 - t,\ y = 3t,\ z = 1,\ t\in \mathbb{R}, & &
 \\q\colon &x = 2s,\ y = 4s,\ z = 1 - s,\ s\in \mathbb{R}. & &
\\\end{alignedat}\]
Výsledek zaokrouhlete na minuty.}
{\(46^{\circ }22'\)}{\(0^{\circ }\)}{\(67^{\circ }18'\)}{\(90^{\circ }\)}{}{}}
\LANGsk{
\Question{
Určte odchýlku priamok \(p\),
\(q\),
kde
\[ \begin{alignedat}{80}
 p\colon &x = 2 - t,\ y = 3t,\ z = 1,\ t\in \mathbb{R}, & &
 \\q\colon &x = 2s,\ y = 4s,\ z = 1 - s,\ s\in \mathbb{R}. & &
\\\end{alignedat}\]
Výsledok zaokrúhlite na minúty.}
{\(46^{\circ }22'\)}{\(0^{\circ }\)}{\(67^{\circ }18'\)}{\(90^{\circ }\)}{}{}}
\LANGes{
\Question{
Halla el ángulo entre dos rectas y aproxima tu respuesta a minutos.  
\[ 
\begin{aligned}p\colon x& = 2 - t , & 
\\y & = 3t ,
\\z & = 1 ;\ t\in \mathbb{R}
\\ \end{aligned}\qquad \qquad \begin{aligned}q\colon x& = 2s, & 
\\y & = 4s ,
\\z & = 1 - s;\ s\in \mathbb{R}
\\ \end{aligned}
\]}
{\(46^{\circ }22'\)}{\(0^{\circ }\)}{\(67^{\circ }18'\)}{\(90^{\circ }\)}{}{}}
\End

\Data\ID{9000101902}
\Updated{2020-07-15 03:07:37}
\Level{B}
\SubArea{Analytic geometry in a space}
\LANGen{
\Question{
Given points \(A = [0;1;2]\),
\(B = [1;2;0]\),
\(C = [1;2;3]\), find the angle
between the lines \(AB\) 
and \(AC\).
Round your answer to the nearest degree.}
{\(90^{\circ }\)}{\(0^{\circ }\)}{\(30^{\circ }\)}{\(60^{\circ }\)}{}{}}
\LANGpl{
\Question{
Dane są punkty \(A = [0;1;2]\),
\(B = [1;2;0]\),
\(C = [1;2;3]\), wyznacz kąt
pomiędzy prostymi \(AB\) 
i \(AC\).
Zaokrągli  odpowiedź do pełnych stopni.}
{\(90^{\circ }\)}{\(0^{\circ }\)}{\(30^{\circ }\)}{\(60^{\circ }\)}{}{}}
\LANGcs{
\Question{
Jsou dány body \(A = [0;1;2]\),
\(B = [1;2;0]\),
\(C = [1;2;3]\). Určete odchylku
přímek \(AB\) 
a \(AC\).
Výsledek zaokrouhlete na stupně.}
{\(90^{\circ }\)}{\(0^{\circ }\)}{\(30^{\circ }\)}{\(60^{\circ }\)}{}{}}
\LANGsk{
\Question{
Sú dané body \(A = [0;1;2]\),
\(B = [1;2;0]\),
\(C = [1;2;3]\). Určte odchýlku
priamok \(AB\) 
a \(AC\).
Výsledok zaokrúhlite na stupňa.}
{\(90^{\circ }\)}{\(0^{\circ }\)}{\(30^{\circ }\)}{\(60^{\circ }\)}{}{}}
\LANGes{
\Question{
Dados los puntos \(A = [0;1;2]\),
\(B = [1;2;0]\),
\(C = [1;2;3]\), halla el ángulo entre la recta \(AB\) 
y la recta \(AC\).
Aproxima el resultado a grados.}
{\(90^{\circ }\)}{\(0^{\circ }\)}{\(30^{\circ }\)}{\(60^{\circ }\)}{}{}}
\End

\Data\ID{9000101903}
\Updated{2020-07-15 03:07:37}
\Level{B}
\SubArea{Analytic geometry in a space}
\LANGen{
\Question{
Given points \(A = [-1;0;3]\),
\(B = [0;2;0]\), find the angle
between the line \(AB\) 
and the line \(m\).
\[ 
\begin{aligned}m\colon x& = 1 + 2t, &
\\y & = -3t,
\\z & = 1;\  t\in \mathbb{R}
\\ \end{aligned}
\]
Round your answer to the nearest minute.}
{\(72^{\circ }45'\)}{\(0^{\circ }\)}{\(48^{\circ }15'\)}{\(90^{\circ }\)}{}{}}
\LANGpl{
\Question{
Dane są punkty \(A = [-1;0;3]\),
\(B = [0;2;0]\), wyznacz kąt
pomiędzy prostą \(AB\), 
a prostą \(m\).
\[ 
\begin{aligned}m\colon x& = 1 + 2t, &
\\y & = -3t,
\\z & = 1;\ t\in \mathbb{R}
\\ \end{aligned}\]
Zaokrągli odpowiedź do pełnych stopni i minut.}
{\(72^{\circ }45'\)}{\(0^{\circ }\)}{\(48^{\circ }15'\)}{\(90^{\circ }\)}{}{}}
\LANGcs{
\Question{
Jsou dány body \(A = [-1;0;3]\),
\(B = [0;2;0]\) a přímka
\(m\colon x = 1 + 2t,\, y = -3t,\, z = 1,\, t\in \mathbb{R}\). Určete odchylku
přímek \(AB\) 
a \(m\).
Výsledek zaokrouhlete na minuty.}
{\(72^{\circ }45'\)}{\(0^{\circ }\)}{\(48^{\circ }15'\)}{\(90^{\circ }\)}{}{}}
\LANGsk{
\Question{
Sú dané body \(A = [-1;0;3]\),
\(B = [0;2;0]\) a priamka
\(m\colon x = 1 + 2t,\, y = -3t,\, z = 1,\, t\in \mathbb{R}\). Určte odchýlku
priamok \(AB\) 
a \(m\).
Výsledok zaokrúhlite na minúty.}
{\(72^{\circ }45'\)}{\(0^{\circ }\)}{\(48^{\circ }15'\)}{\(90^{\circ }\)}{}{}}
\LANGes{
\Question{
Dados los puntos \(A = [-1;0;3]\),
\(B = [0;2;0]\), halla el ángulo
entre la recta \(AB\) 
y la recta \(m\).
\[ 
\begin{aligned}m\colon x& = 1 + 2t, &
\\y & = -3t,
\\z & = 1;\  t\in \mathbb{R}
\\ \end{aligned}
\]
Aproxima el resultado a los minutos.}
{\(72^{\circ }45'\)}{\(0^{\circ }\)}{\(48^{\circ }15'\)}{\(90^{\circ }\)}{}{}}
\End

\Data\ID{9000101904}
\Updated{2020-07-15 03:07:37}
\Level{B}
\SubArea{Analytic geometry in a space}
\LANGen{
\Question{
Find the angle between the \(x\)-axis
and the line \(p\).
\[ 
\begin{aligned}p\colon x& = 2 - t, &
\\y & = 3t,
\\z & = 1;\ t\in \mathbb{R}
\\ \end{aligned}
\]
Round your answer to the nearest minute.}
{\(71^{\circ }34'\)}{\(0^{\circ }\)}{\(69^{\circ }17'\)}{\(90^{\circ }\)}{}{}}
\LANGpl{
\Question{
Wyznacz kąt pomiędzy osią \(x\),
a prostą \(p\).
\[ 
\begin{aligned}p\colon x& = 2 - t,&
\\y & = 3t,
\\z & = 1;\ t\in \mathbb{R}
\\ \end{aligned}
\]
Zaokrągli odpowiedź do pełnych stopni i minut.}
{\(71^{\circ }34'\)}{\(0^{\circ }\)}{\(69^{\circ }17'\)}{\(90^{\circ }\)}{}{}}
\LANGcs{
\Question{
Určete odchylku přímky \(p\colon x = 2 - t,\, y = 3t,\, z = 1,\, t\in \mathbb{R}\) 
od osy \(x\) 
soustavy souřadnic. Výsledek zaokrouhlete na minuty.}
{\(71^{\circ }34'\)}{\(0^{\circ }\)}{\(69^{\circ }17'\)}{\(90^{\circ }\)}{}{}}
\LANGsk{
\Question{
Určte odchýlku priamky \(p\colon x = 2 - t,\, y = 3t,\, z = 1,\, t\in \mathbb{R}\) 
od osy \(x\) 
sústavy súradníc. Výsledok zaokrúhlite na minúty.}
{\(71^{\circ }34'\)}{\(0^{\circ }\)}{\(69^{\circ }17'\)}{\(90^{\circ }\)}{}{}}
\LANGes{
\Question{
Halla el ángulo entre el eje \(x\)
y la recta \(p\).
\[ 
\begin{aligned}p\colon x& = 2 - t, &
\\y & = 3t,
\\z & = 1;\ t\in \mathbb{R}
\\ \end{aligned}
\]
Aproxima la respuesta a los minutos.}
{\(71^{\circ }34'\)}{\(0^{\circ }\)}{\(69^{\circ }17'\)}{\(90^{\circ }\)}{}{}}
\End

\Data\ID{9000101905}
\Updated{2020-07-15 03:07:37}
\Level{B}
\SubArea{Analytic geometry in a space}
\LANGen{
\Question{
The points \(A = [0;5;0]\),
\(B = [5;5;0]\),
\(C = [5;0;0]\),
\(D = [0;0;0]\) define the cube
\(ABCDEFGH\). Find the angle
between the lines \(BF\) 
and \(AC\).
Round your answer to the nearest minute.}
{\(90^{\circ }\)}{\(0^{\circ }\)}{\(45^{\circ }\)}{\(73^{\circ }47'\)}{}{}}
\LANGpl{
\Question{
Punkty \(A = [0;5;0]\),
\(B = [5;5;0]\),
\(C = [5;0;0]\),
\(D = [0;0;0]\) tworzą sześcian
\(ABCDEFGH\). Wyznacz  kąt pomiędzy
prostymi \(BF\) 
i \(AC\).
Zaokrągli odpowiedź do pełnych stopni i minut.}
{\(90^{\circ }\)}{\(0^{\circ }\)}{\(45^{\circ }\)}{\(73^{\circ }47'\)}{}{}}
\LANGcs{
\Question{
Jsou dány body \(A = [0;5;0]\),
\(B = [5;5;0]\),
\(C = [5;0;0]\),
\(D = [0;0;0]\), které tvoří vrcholy
krychle \(ABCDEFGH\). Určete
odchylku přímek \(BF\) 
a \(AC\).
Výsledek zaokrouhlete na minuty.}
{\(90^{\circ }\)}{\(0^{\circ }\)}{\(45^{\circ }\)}{\(73^{\circ }47'\)}{}{}}
\LANGsk{
\Question{
Sú dané body \(A = [0;5;0]\),
\(B = [5;5;0]\),
\(C = [5;0;0]\),
\(D = [0;0;0]\), ktoré tvoria vrcholy
kocky \(ABCDEFGH\). Určte
odchýlku priamok \(BF\) 
a \(AC\).
Výsledok zaokrúhlite na minúty.}
{\(90^{\circ }\)}{\(0^{\circ }\)}{\(45^{\circ }\)}{\(73^{\circ }47'\)}{}{}}
\LANGes{
\Question{
Dados los puntos \(A = [0;5;0]\),
\(B = [5;5;0]\),
\(C = [5;0;0]\),
\(D = [0;0;0]\) que definen el cubo
\(ABCDEFGH\). Halla el ángulo 
entre la recta \(BF\) 
y \(AC\).
Aproxima el resultado a los minutos.}
{\(90^{\circ }\)}{\(0^{\circ }\)}{\(45^{\circ }\)}{\(73^{\circ }47'\)}{}{}}
\End

\Data\ID{9000101906}
\Updated{2020-07-15 03:07:37}
\Level{B}
\SubArea{Analytic geometry in a space}
\LANGen{
\Question{
Find the angle between the planes \(\alpha \) 
and \(\beta \).
\[ 
 \alpha \colon 2x - 5y + 3z - 4 = 0,\qquad \beta \colon x - 3 = 0
\]
Round your answer to the nearest minute.}
{\(71^{\circ }4'\)}{\(21^{\circ }42'\)}{\(82^{\circ }19'\)}{\(85^{\circ }2'\)}{}{}}
\LANGpl{
\Question{
Wyznacz kąt pomiędzy płaszczyznami \(\alpha \) 
i \(\beta \).
\[ 
 \alpha \colon 2x - 5y + 3z - 4 = 0,\qquad \beta \colon x - 3 = 0
\]
Zaokrągli odpowiedź do pełnych stopni i minut.}
{\(71^{\circ }4'\)}{\(21^{\circ }42'\)}{\(82^{\circ }19'\)}{\(85^{\circ }2'\)}{}{}}
\LANGcs{
\Question{
Určete odchylku rovin \(\alpha \colon 2x - 5y + 3z - 4 = 0\),
\(\beta \colon x - 3 = 0\).
Výsledek zaokrouhlete na minuty.}
{\(71^{\circ }4'\)}{\(21^{\circ }42'\)}{\(82^{\circ }19'\)}{\(85^{\circ }2'\)}{}{}}
\LANGsk{
\Question{
Určte odchýlku rovín \(\alpha \colon 2x - 5y + 3z - 4 = 0\),
\(\beta \colon x - 3 = 0\).
Výsledok zaokrúhlite na minúty.}
{\(71^{\circ }4'\)}{\(21^{\circ }42'\)}{\(82^{\circ }19'\)}{\(85^{\circ }2'\)}{}{}}
\LANGes{
\Question{
Halla el ángulo entre los planos \(\alpha \) 
y \(\beta \).
\[ 
 \alpha \colon 2x - 5y + 3z - 4 = 0,\qquad \beta \colon x - 3 = 0
\]
Aproxima el resultado a los minutos.}
{\(71^{\circ }4'\)}{\(21^{\circ }42'\)}{\(82^{\circ }19'\)}{\(85^{\circ }2'\)}{}{}}
\End

\Data\ID{9000101907}
\Updated{2020-07-15 03:07:37}
\Level{B}
\SubArea{Analytic geometry in a space}
\LANGen{
\Question{
The general plane \(\alpha \) 
has the equation 
\[ 
 \alpha \colon 3z - 4 = 0
\]
and the plane \(\beta \) has a
normal vector \(\vec{n} = (0;0;1)\). Find
the angle between \(\alpha \) 
and \(\beta \) 
and round your answer to the nearest degree.}
{\(0^{\circ }\)}{\(30^{\circ }\)}{\(45^{\circ }\)}{\(90^{\circ }\)}{}{}}
\LANGpl{
\Question{
Dana jest płaszczyzna wyrażona równaniem skalarnym \(\alpha \)  
\[ 
 \alpha \colon 3z - 4 = 0
\]
oraz płaszczyzna  \(\beta \) o wektorze 
prostopadłym \(\vec{n} = (0;0;1)\). Wyznacz 
kąt pomiędzy płaszczyznami \(\alpha \) 
i \(\beta \). 
Zaokrągli odpowiedź do pełnych stopni i minut.}
{\(0^{\circ }\)}{\(30^{\circ }\)}{\(45^{\circ }\)}{\(90^{\circ }\)}{}{}}
\LANGcs{
\Question{
Určete odchylku roviny \(\alpha \colon 3z - 4 = 0\) 
od roviny, která má normálový vektor
\(\vec{n} = (0;0;1)\).
Výsledek zaokrouhlete na minuty.}
{\(0^{\circ }\)}{\(30^{\circ }\)}{\(45^{\circ }\)}{\(90^{\circ }\)}{}{}}
\LANGsk{
\Question{
Určte odchýlku roviny \(\alpha \colon 3z - 4 = 0\) 
od roviny, ktorá má normálový vektor
\(\vec{n} = (0;0;1)\).
Výsledok zaokrúhlite na minúty.}
{\(0^{\circ }\)}{\(30^{\circ }\)}{\(45^{\circ }\)}{\(90^{\circ }\)}{}{}}
\LANGes{
\Question{
El plano general \(\alpha \) 
tiene la siguiente ecuación:
\[ 
 \alpha \colon 3z - 4 = 0
\]
y el plano \(\beta \) tiene el vector \(\vec{n} = (0;0;1)\). Halla el ángulo entre \(\alpha \) 
y \(\beta \) 
y aproxima el resultado a los minutos.}
{\(0^{\circ }\)}{\(30^{\circ }\)}{\(45^{\circ }\)}{\(90^{\circ }\)}{}{}}
\End

\Data\ID{9000101908}
\Updated{2020-07-15 03:07:37}
\Level{B}
\SubArea{Analytic geometry in a space}
\LANGen{
\Question{
Find the angle between the line \(p\) 
and the plane \(\alpha \).
\[ 
\alpha \colon x-3z+5 = 0;\qquad \qquad \begin{aligned}[t] p\colon x& = 3, &
\\y & = 3t,
\\z & = 1 - t;\ t\in \mathbb{R}
\\ \end{aligned}
\]
Round your answer to the nearest minute.}
{\(17^{\circ }27'\)}{\(0^{\circ }\)}{\(47^{\circ }33'\)}{\(90^{\circ }\)}{}{}}
\LANGpl{
\Question{
Wyznacz kąt pomiędzy prostą \(p\), 
a płaszczyzną \(\alpha \).
\[ 
\alpha \colon x-3z+5 = 0;\qquad \qquad \begin{aligned}[t] p\colon x& = 3, &
\\y & = 3t,
\\z & = 1 - t;\ t\in \mathbb{R}
\\ \end{aligned}
\]
Zaokrągli odpowiedź do pełnych stopni i minut.}
{\(17^{\circ }27'\)}{\(0^{\circ }\)}{\(47^{\circ }33'\)}{\(90^{\circ }\)}{}{}}
\LANGcs{
\Question{
Určete odchylku přímky \(p\colon x = 3,\, y = 3t,\, z = 1 - t,\, t\in \mathbb{R}\) 
a roviny \(\alpha \colon x - 3z + 5 = 0\).
Výsledek zaokrouhlete na minuty.}
{\(17^{\circ }27'\)}{\(0^{\circ }\)}{\(47^{\circ }33'\)}{\(90^{\circ }\)}{}{}}
\LANGsk{
\Question{
Určte odchýlku priamky \(p\colon x = 3,\, y = 3t,\, z = 1 - t,\, t\in \mathbb{R}\) 
a roviny \(\alpha \colon x - 3z + 5 = 0\).
Výsledok zaokrúhlite na minúty.}
{\(17^{\circ }27'\)}{\(0^{\circ }\)}{\(47^{\circ }33'\)}{\(90^{\circ }\)}{}{}}
\LANGes{
\Question{
Halla el ángulo entre la recta \(p\) 
y el plano \(\alpha \).
\[ 
\alpha \colon x-3z+5 = 0;\qquad \qquad \begin{aligned}[t] p\colon x& = 3, &
\\y & = 3t,
\\z & = 1 - t;\ t\in \mathbb{R}
\\ \end{aligned}
\]
Aproxima el resultado a los minutos.}
{\(17^{\circ }27'\)}{\(0^{\circ }\)}{\(47^{\circ }33'\)}{\(90^{\circ }\)}{}{}}
\End

\Data\ID{9000101909}
\Updated{2020-07-15 03:07:37}
\Level{B}
\SubArea{Analytic geometry in a space}
\LANGen{
\Question{
Given points \(A = [1;0;2]\),
\(B = [1;0;0]\) and the
plane \(\alpha \),
\[ 
 \alpha \colon 2x - 4y = 0,
\]
find the angle between the line \(AB\) 
and the plane \(\alpha \).
Round your answer to the nearest minute.}
{\(0^{\circ }\)}{\(22^{\circ }48'\)}{\(45^{\circ }19'\)}{\(90^{\circ }\)}{}{}}
\LANGpl{
\Question{
Dane są punkty \(A = [1;0;2]\),
\(B = [1;0;0]\) oraz
płaszczyzna \(\alpha \),
\[ 
 \alpha \colon 2x - 4y = 0,
\]
wyznacz kąt pomiędzy prostą \(AB\), 
a płaszczyzną  \(\alpha \).
Zaokrągli odpowiedź do pełnych stopni i minut.}
{\(0^{\circ }\)}{\(22^{\circ }48'\)}{\(45^{\circ }19'\)}{\(90^{\circ }\)}{}{}}
\LANGcs{
\Question{
Jsou dány body \(A = [1;0;2]\),
\(B = [1;0;0]\). Určete odchylku
přímky \(AB\) 
a roviny \(\alpha \colon 2x - 4y = 0\).
Výsledek zaokrouhlete na minuty.}
{\(0^{\circ }\)}{\(22^{\circ }48'\)}{\(45^{\circ }19'\)}{\(90^{\circ }\)}{}{}}
\LANGsk{
\Question{
Sú dané body \(A = [1;0;2]\),
\(B = [1;0;0]\). Určte odchýlku
priamky \(AB\) 
a roviny \(\alpha \colon 2x - 4y = 0\).
Výsledok zaokrúhlite na minúty.}
{\(0^{\circ }\)}{\(22^{\circ }48'\)}{\(45^{\circ }19'\)}{\(90^{\circ }\)}{}{}}
\LANGes{
\Question{
Dados los puntos \(A = [1;0;2]\),
\(B = [1;0;0]\) y el plano \(\alpha \),
\[ 
 \alpha \colon 2x - 4y = 0,
\] halla el ángulo entre la recta \(AB\) 
y el plano \(\alpha \).
Aproxima el resultado a los minutos.}
{\(0^{\circ }\)}{\(22^{\circ }48'\)}{\(45^{\circ }19'\)}{\(90^{\circ }\)}{}{}}
\End

\Data\ID{9000101910}
\Updated{2020-07-15 03:07:37}
\Level{B}
\SubArea{Analytic geometry in a space}
\IMG{001019_10-figure0.pdf}
\LANGen{
\Question{
The points \(A = [0;5;0]\),
\(B = [5;5;0]\),
\(C = [5;0;0]\) and
\(D = [0;0;0]\) define the cube
\(ABCDEFGH\). Find the angle
between the line \(BF\) 
and the plane \(AFE\).
Round your answer to the nearest minute.}
{\(0^{\circ }\)}{\(35^{\circ }16'\)}{\(45^{\circ }\)}{\(90^{\circ }\)}{}{}}
\LANGpl{
\Question{
Punkty \(A = [0;5;0]\),
\(B = [5;5;0]\),
\(C = [5;0;0]\) i
\(D = [0;0;0]\) tworzą sześcian
\(ABCDEFGH\). Wyznacz kąt
pomiędzy prostą \(BF\), 
a płaszczyzną  \(AFE\).
Zaokrągli odpowiedź do pełnych stopni i minut.}
{\(0^{\circ }\)}{\(35^{\circ }16'\)}{\(45^{\circ }\)}{\(90^{\circ }\)}{}{}}
\LANGcs{
\Question{
Jsou dány body \(A = [0;5;0]\),
\(B = [5;5;0]\),
\(C = [5;0;0]\),
\(D = [0;0;0]\), které tvoří vrcholy
krychle \(ABCDEFGH\). Určete
odchylku přímky \(BF\) 
a roviny \(AFE\).
Výsledek zaokrouhlete na minuty.}
{\(0^{\circ }\)}{\(35^{\circ }16'\)}{\(45^{\circ }\)}{\(90^{\circ }\)}{}{}}
\LANGsk{
\Question{
Sú dané body \(A = [0;5;0]\),
\(B = [5;5;0]\),
\(C = [5;0;0]\),
\(D = [0;0;0]\), ktoré tvoria vrcholy
kocky \(ABCDEFGH\). Určte
odchýlku priamky \(BF\) 
a roviny \(AFE\).
Výsledok zaokrúhlite na minúty.}
{\(0^{\circ }\)}{\(35^{\circ }16'\)}{\(45^{\circ }\)}{\(90^{\circ }\)}{}{}}
\LANGes{
\Question{
Dados los puntos \(A = [0;5;0]\),
\(B = [5;5;0]\),
\(C = [5;0;0]\) y
\(D = [0;0;0]\) que definen el cubo
\(ABCDEFGH\). Halla el ángulo
entre la recta \(BF\) 
y el plano \(AFE\).
Aproxima el resultado a los minutos.}
{\(0^{\circ }\)}{\(35^{\circ }16'\)}{\(45^{\circ }\)}{\(90^{\circ }\)}{}{}}
\End

\Data\ID{9000106301}
\Updated{2020-07-15 03:07:37}
\Level{B}
\SubArea{Analytic geometry in a space}
\IMG{001063_01-figure0.pdf}
\LANGen{
\Question{
Find the line $k$ which is perpendicular to the plane
\(\alpha \) 
\[ 
 \alpha \colon 2x + y - z - 5 = 0
\]
and passes through the point \(A = [0;0;1]\).}
{\(\begin{aligned}[t] x& =\phantom{ 1 -} 2t, &
\\y& =\phantom{ 1 -}\ t,
\\z& = 1 - t;\ t\in \mathbb{R}
\\ \end{aligned}\)}{\(\begin{aligned}[t] x& =\phantom{ -}2 + 2m, &
\\y& =\phantom{ -}1 +\phantom{ 2}m,
\\z& = -1 -\phantom{ 2}m;\ m\in \mathbb{R}
\\ \end{aligned}\)}{\(\begin{aligned}[t] x& =\phantom{ -}2k, &
\\y& =\phantom{ -2}k,
\\z& = -\phantom{2}k;\ k\in \mathbb{R}
\\ \end{aligned}\)}{\(\begin{aligned}[t] x& =\phantom{ -}2, &
\\y& =\phantom{ -}1,
\\z& = -1 + u;\ u\in \mathbb{R}
\\ \end{aligned}\)}{}{}}
\LANGpl{
\Question{
Wyznacz prostą $k$ prostopadłą do płaszczyzny
\(\alpha \) 
\[ 
 \alpha \colon 2x + y - z - 5 = 0
\]
i przechodzącą przez punkt \(A = [0;0;1]\).}
{\(\begin{aligned}[t] x& =\phantom{ 1 -} 2t, &
\\y& =\phantom{ 1 -}\ t,
\\z& = 1 - t;\ t\in \mathbb{R}
\\ \end{aligned}\)}{\(\begin{aligned}[t] x& =\phantom{ -}2 + 2m, &
\\y& =\phantom{ -}1 +\phantom{ 2}m,
\\z& = -1 -\phantom{ 2}m;\ m\in \mathbb{R}
\\ \end{aligned}\)}{\(\begin{aligned}[t] x& =\phantom{ -}2k, &
\\y& =\phantom{ -2}k,
\\z& = -\phantom{2}k;\ k\in \mathbb{R}
\\ \end{aligned}\)}{\(\begin{aligned}[t] x& =\phantom{ -}2, &
\\y& =\phantom{ -}1,
\\z& = -1 + u;\ u\in \mathbb{R}
\\ \end{aligned}\)}{}{}}
\LANGcs{
\Question{
Rovina \(\alpha \) je zadaná
obecnou rovnicí: \(2x + y - z - 5 = 0\).
Určete parametrické vyjádření přímky
\(k\), která je kolmá
na rovinu \(\alpha \) a
prochází bodem \(A = [0;0;1]\).}
{\(\begin{aligned}[t] x& =\phantom{ 1 -} 2t, &
\\y& =\phantom{ 1 -}\ t,
\\z& = 1 - t;\ t\in \mathbb{R}
\\ \end{aligned}\)}{\(\begin{aligned}[t] x& =\phantom{ -}2 + 2m, &
\\y& =\phantom{ -}1 +\phantom{ 2}m,
\\z& = -1 -\phantom{ 2}m;\ m\in \mathbb{R}
\\ \end{aligned}\)}{\(\begin{aligned}[t] x& =\phantom{ -}2k, &
\\y& =\phantom{ -2}k,
\\z& = -\phantom{2}k;\ k\in \mathbb{R}
\\ \end{aligned}\)}{\(\begin{aligned}[t] x& =\phantom{ -}2, &
\\y& =\phantom{ -}1,
\\z& = -1 + u;\ u\in \mathbb{R}
\\ \end{aligned}\)}{}{}}
\LANGsk{
\Question{
Rovina \(\alpha \) je daná
všeobecnou rovnicou: \(2x + y - z - 5 = 0\).
Určte parametrické vyjadrenie priamky
\(k\), ktorá je kolmá
na rovinu \(\alpha \) a
prechádza bodom \(A = [0;0;1]\).}
{\(\begin{aligned}[t] x& =\phantom{ 1 -} 2t, &
\\y& =\phantom{ 1 -}\ t,
\\z& = 1 - t;\ t\in \mathbb{R}
\\ \end{aligned}\)}{\(\begin{aligned}[t] x& =\phantom{ -}2 + 2m, &
\\y& =\phantom{ -}1 +\phantom{ 2}m,
\\z& = -1 -\phantom{ 2}m;\ m\in \mathbb{R}
\\ \end{aligned}\)}{\(\begin{aligned}[t] x& =\phantom{ -}2k, &
\\y& =\phantom{ -2}k,
\\z& = -\phantom{2}k;\ k\in \mathbb{R}
\\ \end{aligned}\)}{\(\begin{aligned}[t] x& =\phantom{ -}2, &
\\y& =\phantom{ -}1,
\\z& = -1 + u;\ u\in \mathbb{R}
\\ \end{aligned}\)}{}{}}
\LANGes{
\Question{
Halla la recta $k$ que es perpendicular al plano \(\alpha \) 
\[ 
 \alpha \colon 2x + y - z - 5 = 0
\]
y que pasa por el punto \(A = [0;0;1]\).}
{\(\begin{aligned}[t] x& =\phantom{ 1 -} 2t, &
\\y& =\phantom{ 1 -}\ t,
\\z& = 1 - t;\ t\in \mathbb{R}
\\ \end{aligned}\)}{\(\begin{aligned}[t] x& =\phantom{ -}2 + 2m, &
\\y& =\phantom{ -}1 +\phantom{ 2}m,
\\z& = -1 -\phantom{ 2}m;\ m\in \mathbb{R}
\\ \end{aligned}\)}{\(\begin{aligned}[t] x& =\phantom{ -}2k, &
\\y& =\phantom{ -2}k,
\\z& = -\phantom{2}k;\ k\in \mathbb{R}
\\ \end{aligned}\)}{\(\begin{aligned}[t] x& =\phantom{ -}2, &
\\y& =\phantom{ -}1,
\\z& = -1 + u;\ u\in \mathbb{R}
\\ \end{aligned}\)}{}{}}
\End

\Data\ID{9000106302}
\Updated{2020-07-15 03:07:37}
\Level{B}
\SubArea{Analytic geometry in a space}
\IMG{001063_02-figure0.pdf}
\LANGen{
\Question{
The plane \(\alpha \) 
has equation 
\[ 
 \alpha : 2x + y - z - 5 = 0.
\]
The line \(k\) passes
through the point \(A = [0;0;1]\) and
is perpendicular to \(\alpha \).
Find the intersection \(S\) 
of the line \(k\) and
the plane \(\alpha \).}
{\(S = [2;1;0]\)}{\(S = [2;0;1]\)}{\(S = [-2;1;0]\)}{\(S = [-2;0;1]\)}{}{}}
\LANGpl{
\Question{
Dana jest płaszczyzna  \(\alpha \) 
przedstawiona za pomocą równania
\[ 
 \alpha : 2x + y - z - 5 = 0.
\]
Prosta  \(k\) przechodząca przez
punkt \(A = [0;0;1]\) 
prostopadła do \(\alpha \).
Wyznacz punkt przecięcia  \(S\) 
prostej \(k\) i
płaszczyzny \(\alpha \).}
{\(S = [2;1;0]\)}{\(S = [2;0;1]\)}{\(S = [-2;1;0]\)}{\(S = [-2;0;1]\)}{}{}}
\LANGcs{
\Question{
Rovina \(\alpha \) je zadaná
obecnou rovnicí: \(2x + y - z - 5 = 0\).
Bodem \(A = [0;0;1]\) je vedena kolmice
\(k\) k této rovině. Určete
souřadnice bodu \(S\),
ve kterém kolmice \(k\) 
protíná danou rovinu.}
{\(S = [2;1;0]\)}{\(S = [2;0;1]\)}{\(S = [-2;1;0]\)}{\(S = [-2;0;1]\)}{}{}}
\LANGsk{
\Question{
Rovina \(\alpha \) je daná
všeobecnou rovnicou: \(2x + y - z - 5 = 0\).
Bodom \(A = [0;0;1]\) je vedená kolmica
\(k\) k tejto rovine. Určte
súradnice bodu \(S\),
v ktorom kolmica \(k\) 
pretína danú rovinu.}
{\(S = [2;1;0]\)}{\(S = [2;0;1]\)}{\(S = [-2;1;0]\)}{\(S = [-2;0;1]\)}{}{}}
\LANGes{
\Question{
El plano \(\alpha \) 
tiene la ecuación 
\[ 
 \alpha : 2x + y - z - 5 = 0.
\]
La recta \(k\) pasa
por el punto \(A = [0;0;1]\) y
es perpendicular al plano \(\alpha \).
Halla la intersección \(S\) 
de la recta \(k\) y
el plano \(\alpha \).}
{\(S = [2;1;0]\)}{\(S = [2;0;1]\)}{\(S = [-2;1;0]\)}{\(S = [-2;0;1]\)}{}{}}
\End

\Data\ID{9000106304}
\Updated{2020-07-15 03:07:37}
\Level{B}
\SubArea{Analytic geometry in a space}
\IMG{001063_04-figure0.pdf}
\LANGen{
\Question{
Find the third coordinate of the point
\(B = [2;0;?]\) using the fact that this
point is in the plane \(\alpha \) 
defined by the equation 
\[ 
 \alpha \colon 2x + y - z - 5 = 0.
\]
Use the point \(B\) to
find the angle \(\varphi \) 
between the plane \(\alpha \) 
and the line \(AB\),
where \(A = [0;0;1]\).}
{\(\varphi = 60^{\circ }\)}{\(\varphi = 45^{\circ }\)}{\(\varphi = 30^{\circ }\)}{\(\varphi = 75^{\circ }\)}{}{}}
\LANGpl{
\Question{
Wyznacz trzecią współrzędną punktu
\(B = [2;0;?]\) tak, aby punkt
leżał na płaszczyźnie  \(\alpha \) 
wyrażonej równaniem
\[ 
 \alpha \colon 2x + y - z - 5 = 0.
\]
Za pomocą punktu \(B\) 
wyznacz kąt  \(\varphi \) 
pomiędzy płaszczyzną \(\alpha \), 
a prostą \(AB\),
jeśli \(A = [0;0;1]\).}
{\(\varphi = 60^{\circ }\)}{\(\varphi = 45^{\circ }\)}{\(\varphi = 30^{\circ }\)}{\(\varphi = 75^{\circ }\)}{}{}}
\LANGcs{
\Question{
V rovině \(\alpha \) zadané
obecnou rovnicí \(2x + y - z - 5 = 0\) 
leží bod \(B = [2;0;?]\).
Určete odchylku \(\varphi \) 
přímky \(AB\),
kde \(A = [0;0;1]\), od
roviny \(\alpha \).}
{\(\varphi = 60^{\circ }\)}{\(\varphi = 45^{\circ }\)}{\(\varphi = 30^{\circ }\)}{\(\varphi = 75^{\circ }\)}{}{}}
\LANGsk{
\Question{
V rovine \(\alpha \) zadanej
všeobecnou rovnicou \(2x + y - z - 5 = 0\) 
leží bod \(B = [2;0;?]\).
Určte odchýlku \(\varphi \) 
priamky \(AB\),
kde \(A = [0;0;1]\), od
roviny \(\alpha \).}
{\(\varphi = 60^{\circ }\)}{\(\varphi = 45^{\circ }\)}{\(\varphi = 30^{\circ }\)}{\(\varphi = 75^{\circ }\)}{}{}}
\LANGes{
\Question{
Halla la tercera coordenada del punto 
\(B = [2;0;?]\) tomando en cuenta que este punto está en el plano \(\alpha \) 
definido por la ecuación 
\[ 
 \alpha \colon 2x + y - z - 5 = 0.
\]
Usa el punto \(B\) para encontrar el ángulo \(\varphi \) 
entre el plano \(\alpha \) 
y la recta \(AB\),
donde \(A = [0;0;1]\).}
{\(\varphi = 60^{\circ }\)}{\(\varphi = 45^{\circ }\)}{\(\varphi = 30^{\circ }\)}{\(\varphi = 75^{\circ }\)}{}{}}
\End

\Data\ID{9000106305}
\Updated{2020-07-15 03:07:37}
\Level{B}
\SubArea{Analytic geometry in a space}
\IMG{001063_05-figure0.pdf}
\LANGen{
\Question{
Find the area of the triangle \(ABS\).
Only first two coordinates of the point $B=[2;0;?]$ are given and $B$ lies in the plane $\alpha$
defined by the equation 
\[ 
 \alpha \colon 2x + y - z - 5 = 0.
\]
The point \(S\) is the
intersection point of the plane \(\alpha \) 
and the line \(k\) which is
perpendicular to \(\alpha \) and
passes through the point \(A = [0;0;1]\).}
{\(\sqrt{3}\)}{\(2\)}{\(4\)}{\(\sqrt{6}\)}{}{}}
\LANGpl{
\Question{
Oblicz pole trójkąta \(ABS\).
Podano dwie pierwsze współrzędne punktu \(B = [2;0;?]\) 
Punkt B leży na płaszczyźnie 
\(\alpha \) 
określonej równaniem
\[ 
 \alpha \colon 2x + y - z - 5 = 0.
\]
Punkt \(S\) jest punktem
przecięcia płaszczyzny \(\alpha \) 
i prostej  \(k\), która jest prostopadła do
 \(\alpha \) i
przechodzi przez punkt \(A = [0;0;1]\).}
{\(\sqrt{3}\)}{\(2\)}{\(4\)}{\(\sqrt{6}\)}{}{}}
\LANGcs{
\Question{
V rovině \(\alpha \) zadané
obecnou rovnicí \(2x + y - z - 5 = 0\) 
leží bod \(B = [2;0;?]\). Určete
obsah trojúhelníka \(ABS\),
kde \(A = [0;0;1]\) a
\(S\) je pata kolmice
\(k\) vedené
bodem \(A\) k
rovině \(\alpha \).}
{\(\sqrt{3}\)}{\(2\)}{\(4\)}{\(\sqrt{6}\)}{}{}}
\LANGsk{
\Question{
V rovine \(\alpha \) zadanej
všeobecnou rovnicou \(2x + y - z - 5 = 0\) 
leží bod \(B = [2;0;?]\). Určte
obsah trojuholníka \(ABS\),
kde \(A = [0;0;1]\) a
\(S\) je päta kolmice
\(k\) vedenej
bodom \(A\) k
rovine \(\alpha \).}
{\(\sqrt{3}\)}{\(2\)}{\(4\)}{\(\sqrt{6}\)}{}{}}
\LANGes{
\Question{
Halla la superficie del triángulo \(ABS\).
Dadas solo dos coordenadas del punto $B=[2;0;?]$. El punto $B$  está en el plano $\alpha$
definido por la ecuación siguiente 
\[ 
 \alpha \colon 2x + y - z - 5 = 0.
\]
El punto \(S\) es el punto de intersección 
del plano \(\alpha \) y la recta \(k\) que es 
perpendicular al plano \(\alpha \) y
pasa por el punto \(A = [0;0;1]\).}
{\(\sqrt{3}\)}{\(2\)}{\(4\)}{\(\sqrt{6}\)}{}{}}
\End

\Data\ID{9000106306}
\Updated{2020-07-15 03:07:37}
\Level{B}
\SubArea{Analytic geometry in a space}
\IMG{001063_06-figure0.pdf}
\LANGen{
\Question{
Find the general equation of the plane which is perpendicular to the plane
\(\alpha \) 
\[ 
 \alpha \colon 2x + y - z - 5 = 0
\]
and contains the line \(AB\),
where \(A = [0;0;1]\) 
and \(B\) is a
point in \(\alpha \) 
defined by it's first two coordinates 
\[ 
 B = [2;0;?].
\]}
{\(x - y + z - 1 = 0\)}{\(x + y - z + 1 = 0\)}{\(2x - y + z - 1 = 0\)}{\(- 2x + y - z + 1 = 0\)}{}{}}
\LANGpl{
\Question{
Wyznacz równanie skalarne płaszczyzny prostopadłej do płaszczyzny
\(\alpha \) 
\[ 
 \alpha \colon 2x + y - z - 5 = 0
\]
zawierającej prostą \(AB\),
gdzie \(A = [0;0;1]\) 
i \(B\) jest punktem na płaszczyźnie 
 \(\alpha \) 
określonym przez dwie pierwsze współrzędne
\[ 
 B = [2;0;?].
\]}
{\(x - y + z - 1 = 0\)}{\(x + y - z + 1 = 0\)}{\(2x - y + z - 1 = 0\)}{\(- 2x + y - z + 1 = 0\)}{}{}}
\LANGcs{
\Question{
Určete obecnou rovnici roviny, která je kolmá k rovině
\[ 
 \alpha \colon 2x + y - z - 5 = 0
\]
 a která prochází přímkou \(AB\),
je-li \(A = [0;0;1]\) a
víme-li, že \(B = [2;0;?]\in \alpha \).}
{\(x - y + z - 1 = 0\)}{\(x + y - z + 1 = 0\)}{\(2x - y + z - 1 = 0\)}{\(- 2x + y - z + 1 = 0\)}{}{}}
\LANGsk{
\Question{
Určte všeobecnú rovnicu roviny, ktorá je kolmá k rovine
\[ 
 \alpha \colon 2x + y - z - 5 = 0
\]
 a ktorá prechádza priamkou \(AB\),
ak \(A = [0;0;1]\) a
vieme, že \(B = [2;0;?]\in \alpha \).}
{\(x - y + z - 1 = 0\)}{\(x + y - z + 1 = 0\)}{\(2x - y + z - 1 = 0\)}{\(- 2x + y - z + 1 = 0\)}{}{}}
\LANGes{
\Question{
Halla la ecuación general del plano que es perpendicular al plano 
\(\alpha \) 
\[ 
 \alpha \colon 2x + y - z - 5 = 0
\]
y por el que pasa la recta \(AB\),
donde \(A = [0;0;1]\) 
y \(B\) es el 
punto en el plano \(\alpha \)
definido solo por las primeras dos coordenadas 
\[ 
 B = [2;0;?].
\]}
{\(x - y + z - 1 = 0\)}{\(x + y - z + 1 = 0\)}{\(2x - y + z - 1 = 0\)}{\(- 2x + y - z + 1 = 0\)}{}{}}
\End

\Data\ID{9000106308}
\Updated{2020-07-15 03:07:37}
\Level{B}
\SubArea{Analytic geometry in a space}
\LANGen{
\Question{
In the following list identify a pair of planes such that the distance of planes from the plane $\alpha$ is the same as the distance between the point $A=[0;0;1]$ and the plane \(\alpha \). 
\[ 
 \alpha \colon 2x + y - z - 5 = 0
\]}
{\(\begin{aligned}[t] 2x + y - z +\phantom{ 1}1& = 0&
\\2x + y - z - 11& = 0
\\ \end{aligned}\)}{\(\begin{aligned}[t] 2x + y - z +\phantom{ 1}1& = 0&
\\2x + y - z - 10& = 0
\\ \end{aligned}\)}{\(\begin{aligned}[t] 2x + y - z +\phantom{ 1}1& = 0&
\\2x + y - z - 12& = 0
\\ \end{aligned}\)}{\(\begin{aligned}[t] 2x + y - z + 1& = 0&
\\2x + y - z - 9& = 0
\\ \end{aligned}\)}{}{}}
\LANGpl{
\Question{
Wyznacz parę płaszczyzn tak, aby odległość pomiędzy nimi była taka sama jak odległość pomiędzy punktem
\(A = [0;0;1]\) a płaszczyzną 
 \(\alpha \) 
\[ 
 \alpha \colon 2x + y - z - 5 = 0.
\]}
{\(\begin{aligned}[t] 2x + y - z +\phantom{ 1}1& = 0&
\\2x + y - z - 11& = 0
\\ \end{aligned}\)}{\(\begin{aligned}[t] 2x + y - z +\phantom{ 1}1& = 0&
\\2x + y - z - 10& = 0
\\ \end{aligned}\)}{\(\begin{aligned}[t] 2x + y - z +\phantom{ 1}1& = 0&
\\2x + y - z - 12& = 0
\\ \end{aligned}\)}{\(\begin{aligned}[t] 2x + y - z + 1& = 0&
\\2x + y - z - 9& = 0
\\ \end{aligned}\)}{}{}}
\LANGcs{
\Question{
Vyberte dvojici rovin, jejichž vzdálenost od roviny
\[ 
 \alpha \colon 2x + y - z - 5 = 0
\]
 je stejná jako vzdálenost bodu \(A = [0;0;1]\) 
od roviny \(\alpha \).}
{\(\begin{aligned}[t] 2x + y - z +\phantom{ 1}1& = 0&
\\2x + y - z - 11& = 0
\\ \end{aligned}\)}{\(\begin{aligned}[t] 2x + y - z +\phantom{ 1}1& = 0&
\\2x + y - z - 10& = 0
\\ \end{aligned}\)}{\(\begin{aligned}[t] 2x + y - z +\phantom{ 1}1& = 0&
\\2x + y - z - 12& = 0
\\ \end{aligned}\)}{\(\begin{aligned}[t] 2x + y - z + 1& = 0&
\\2x + y - z - 9& = 0
\\ \end{aligned}\)}{}{}}
\LANGsk{
\Question{
Vyberte dvojicu rovín, ktorých vzdialenosť od roviny
\[ 
 \alpha \colon 2x + y - z - 5 = 0
\]
 je rovnaká ako vzdialenosť bodu \(A = [0;0;1]\) 
od roviny \(\alpha \).}
{\(\begin{aligned}[t] 2x + y - z +\phantom{ 1}1& = 0&
\\2x + y - z - 11& = 0
\\ \end{aligned}\)}{\(\begin{aligned}[t] 2x + y - z +\phantom{ 1}1& = 0&
\\2x + y - z - 10& = 0
\\ \end{aligned}\)}{\(\begin{aligned}[t] 2x + y - z +\phantom{ 1}1& = 0&
\\2x + y - z - 12& = 0
\\ \end{aligned}\)}{\(\begin{aligned}[t] 2x + y - z + 1& = 0&
\\2x + y - z - 9& = 0
\\ \end{aligned}\)}{}{}}
\LANGes{
\Question{
Identifica la pareja de planos cuya distancia del plano $\alpha$ es igual a la distancia entre el punto $A=[0;0;1]$ y el plano \(\alpha \). 
\[ 
 \alpha \colon 2x + y - z - 5 = 0
\]}
{\(\begin{aligned}[t] 2x + y - z +\phantom{ 1}1& = 0&
\\2x + y - z - 11& = 0
\\ \end{aligned}\)}{\(\begin{aligned}[t] 2x + y - z +\phantom{ 1}1& = 0&
\\2x + y - z - 10& = 0
\\ \end{aligned}\)}{\(\begin{aligned}[t] 2x + y - z +\phantom{ 1}1& = 0&
\\2x + y - z - 12& = 0
\\ \end{aligned}\)}{\(\begin{aligned}[t] 2x + y - z + 1& = 0&
\\2x + y - z - 9& = 0
\\ \end{aligned}\)}{}{}}
\End

\Data\ID{9000111801}
\Updated{2020-07-15 03:07:37}
\Level{B}
\SubArea{Analytic geometry in a space}
\LANGen{
\Question{
In the following list identify a point such that the distance from this point to the
plane 
\[ 
 2x + y - 2z + 2 = 0
\]
equals \(2\).}
{\([1;-2;4]\)}{\([1;0;1]\)}{\([1;2;3]\)}{}{}{}}
\LANGpl{
\Question{
Wskaż punkt tak, aby odległość od tego punktu do płaszczyzny
\[ 
 2x + y - 2z + 2 = 0
\]
wynosiła \(2\).}
{\([1;-2;4]\)}{\([1;0;1]\)}{\([1;2;3]\)}{}{}{}}
\LANGcs{
\Question{
Pro který z následujících bodů platí, že jeho vzdálenost od roviny dané obecnou rovnicí
\[2x + y - 2z + 2 = 0\] je rovna
\(2\)?}
{\([1;-2;4]\)}{\([1;0;1]\)}{\([1;2;3]\)}{}{}{}}
\LANGsk{
\Question{
Pre ktorý z nasledujúcich bodov platí, že jeho vzdialenosť od roviny danej všeobecnou rovnicou
\[2x + y - 2z + 2 = 0\] je rovná
\(2\)?}
{\([1;-2;4]\)}{\([1;0;1]\)}{\([1;2;3]\)}{}{}{}}
\LANGes{
\Question{
Identifica para cuál de los siguientes puntos, la distancia entre el punto y el plano definido por la ecuación  
\[ 
 2x + y - 2z + 2 = 0
\]
es igual a \(2\).}
{\([1;-2;4]\)}{\([1;0;1]\)}{\([1;2;3]\)}{}{}{}}
\End

\Data\ID{9000111802}
\Updated{2020-07-15 03:07:37}
\Level{B}
\SubArea{Analytic geometry in a space}
\LANGen{
\Question{
In the following list identify a line parallel to the plane
\(\rho \) 
such that the distance between the line and the plane equals
\(1\).
\[ 
\begin{aligned}[t] \rho \colon x& = 1 + r, &
\\y& = 1 + 2s,
\\z& = 1 + r + s;\ r,s\in \mathbb{R}
\\ \end{aligned}
\]}
{\(\begin{aligned}[t] o\colon x& = t, &
\\y & = 2 + 2t,
\\z & = -1 + 2t;\ t\in \mathbb{R}
\\ \end{aligned}\)}{\(\begin{aligned}[t] p\colon x& = 1 - 2t, &
\\y & = -3 - t,
\\z & = 2 + 2t;\ t\in \mathbb{R}
\\ \end{aligned}\)}{\(\begin{aligned}[t] q\colon x& = 1 - 2t, &
\\y & = -3 - t,
\\z & = 1 + 2t;\ t\in \mathbb{R}
\\ \end{aligned}\)}{}{}{}}
\LANGpl{
\Question{
Wskaż prostą równoległą do płaszczyzny
\(\rho \) 
tak, aby odległość pomiędzy prostą a płaszczyzną wynosiła
\(1\).
\[ 
\begin{aligned}[t] \rho \colon x& = 1 + r, &
\\y& = 1 + 2s,
\\z& = 1 + r + s;\ r,s\in \mathbb{R}
\\ \end{aligned}
\]}
{\(\begin{aligned}[t] o\colon x& = t, &
\\y & = 2 + 2t,
\\z & = -1 + 2t;\ t\in \mathbb{R}
\\ \end{aligned}\)}{\(\begin{aligned}[t] p\colon x& = 1 - 2t, &
\\y & = -3 - t,
\\z & = 2 + 2t;\ t\in \mathbb{R}
\\ \end{aligned}\)}{\(\begin{aligned}[t] q\colon x& = 1 - 2t, &
\\y & = -3 - t,
\\z & = 1 + 2t;\ t\in \mathbb{R}
\\ \end{aligned}\)}{}{}{}}
\LANGcs{
\Question{
Pro kterou z následujících přímek platí, že její vzdálenost od roviny
\(\rho \) je rovna \(1\)?
\[ 
\begin{aligned}[t] \rho \colon x& = 1 + r, &
\\y& = 1 + 2s,
\\z& = 1 + r + s;\ r,s\in \mathbb{R}
\\ \end{aligned}
\]}
{\(\begin{aligned}[t] o\colon x& = t, &
\\y & = 2 + 2t,
\\z & = -1 + 2t;\ t\in \mathbb{R}
\\ \end{aligned}\)}{\(\begin{aligned}[t] p\colon x& = 1 - 2t, &
\\y & = -3 - t,
\\z & = 2 + 2t;\ t\in \mathbb{R}
\\ \end{aligned}\)}{\(\begin{aligned}[t] q\colon x& = 1 - 2t, &
\\y & = -3 - t,
\\z & = 1 + 2t;\ t\in \mathbb{R}
\\ \end{aligned}\)}{}{}{}}
\LANGsk{
\Question{
Pre ktorú z nasledujúcich priamok platí, že jej vzdialenosť od roviny
\(\rho \) je rovná \(1\)?
\[ 
\begin{aligned}[t] \rho \colon x& = 1 + r, &
\\y& = 1 + 2s,
\\z& = 1 + r + s;\ r,s\in \mathbb{R}
\\ \end{aligned}
\]}
{\(\begin{aligned}[t] o\colon x& = t, &
\\y & = 2 + 2t,
\\z & = -1 + 2t;\ t\in \mathbb{R}
\\ \end{aligned}\)}{\(\begin{aligned}[t] p\colon x& = 1 - 2t, &
\\y & = -3 - t,
\\z & = 2 + 2t;\ t\in \mathbb{R}
\\ \end{aligned}\)}{\(\begin{aligned}[t] q\colon x& = 1 - 2t, &
\\y & = -3 - t,
\\z & = 1 + 2t;\ t\in \mathbb{R}
\\ \end{aligned}\)}{}{}{}}
\LANGes{
\Question{
Identifica para cuál de las rectas paralelas, la distancia entre el plano 
\(\rho \) 
es igual a 
\(1\).
\[ 
\begin{aligned}[t] \rho \colon x& = 1 + r, &
\\y& = 1 + 2s,
\\z& = 1 + r + s;\ r,s\in \mathbb{R}
\\ \end{aligned}
\]}
{\(\begin{aligned}[t] o\colon x& = t, &
\\y & = 2 + 2t,
\\z & = -1 + 2t;\ t\in \mathbb{R}
\\ \end{aligned}\)}{\(\begin{aligned}[t] p\colon x& = 1 - 2t, &
\\y & = -3 - t,
\\z & = 2 + 2t;\ t\in \mathbb{R}
\\ \end{aligned}\)}{\(\begin{aligned}[t] q\colon x& = 1 - 2t, &
\\y & = -3 - t,
\\z & = 1 + 2t;\ t\in \mathbb{R}
\\ \end{aligned}\)}{}{}{}}
\End

\Data\ID{9000111803}
\Updated{2020-07-15 03:07:37}
\Level{B}
\SubArea{Analytic geometry in a space}
\LANGen{
\Question{
In the following list identify a point such that the distance from this point to the line
\(p\) is
\(\sqrt{3}\).
\[ 
\begin{aligned}[t] p\colon x& = 2 - t, &
\\y & = -1 + 2t,
\\z & = t;\ t\in \mathbb{R}
\\ \end{aligned}
\]}
{\([2;2;0]\)}{\([5;-1;-3]\)}{\([1;1;1]\)}{}{}{}}
\LANGpl{
\Question{
Wyznacz punkt tak, aby odległość od tego punktu do prostej
\(p\) wynosiła
\(\sqrt{3}\).
\[ 
\begin{aligned}[t] p\colon x& = 2 - t, &
\\y & = -1 + 2t,
\\z & = t;\ t\in \mathbb{R}
\\ \end{aligned}
\]}
{\([2;2;0]\)}{\([5;-1;-3]\)}{\([1;1;1]\)}{}{}{}}
\LANGcs{
\Question{
Pro který z následujících bodů platí, že jeho vzdálenost od přímky
\(p\) je rovna \(\sqrt{3}\)?
\[ 
\begin{aligned}[t] p\colon x& = 2 - t, &
\\y & = -1 + 2t,
\\z & = t;\ t\in \mathbb{R}
\\ \end{aligned}
\]}
{\([2;2;0]\)}{\([5;-1;-3]\)}{\([1;1;1]\)}{}{}{}}
\LANGsk{
\Question{
Pre ktorý z nasledujúcich bodov platí, že jeho vzdialenosť od priamky
\(p\) je rovná \(\sqrt{3}\)?
\[ 
\begin{aligned}[t] p\colon x& = 2 - t, &
\\y & = -1 + 2t,
\\z & = t;\ t\in \mathbb{R}
\\ \end{aligned}
\]}
{\([2;2;0]\)}{\([5;-1;-3]\)}{\([1;1;1]\)}{}{}{}}
\LANGes{
\Question{
Identifica para cuál de los puntos la distancia entre el punto y la recta 
\(p\) es igual a
\(\sqrt{3}\).
\[ 
\begin{aligned}[t] p\colon x& = 2 - t, &
\\y & = -1 + 2t,
\\z & = t;\ t\in \mathbb{R}
\\ \end{aligned}
\]}
{\([2;2;0]\)}{\([5;-1;-3]\)}{\([1;1;1]\)}{}{}{}}
\End

\Data\ID{9000111804}
\Updated{2020-07-15 03:07:37}
\Level{B}
\SubArea{Analytic geometry in a space}
\LANGen{
\Question{
In the following list identify a line such that the line is parallel to
\(s\) and the distance
between both lines is \(\sqrt{5}\).
\[ 
\begin{aligned}[t] s\colon x& = -1 + t,&
\\y & = 2t,
\\z & = 2 - t;\ t\in \mathbb{R}
\\ \end{aligned}
\]}
{\(\begin{aligned}[t] r\colon x& = 3 - 2t,&
\\y & = 3 - 4t,
\\z & = 2t;\ t\in \mathbb{R}
\\ \end{aligned}\)}{\(\begin{aligned}[t] q\colon x& = 1, &
\\y & = -1 + 5t,
\\z & = 2 - 2t;\ t\in \mathbb{R}
\\ \end{aligned}\)}{\(\begin{aligned}[t] p\colon x& = -5 - t,&
\\y & = 2 - 2t,
\\z & = 2 + t;\ t\in \mathbb{R}
\\ \end{aligned}\)}{}{}{}}
\LANGpl{
\Question{
Wyznacz prostą tak, aby była równoległa do
\(s\), a odległość pomiędzy obiema prostymi wynosiła
 \(\sqrt{5}\).
\[ 
\begin{aligned}[t] s\colon x& = -1 + t,&
\\y & = 2t,
\\z & = 2 - t;\ t\in \mathbb{R}
\\ \end{aligned}
\]}
{\(\begin{aligned}[t] r\colon x& = 3 - 2t,&
\\y & = 3 - 4t,
\\z & = 2t;\ t\in \mathbb{R}
\\ \end{aligned}\)}{\(\begin{aligned}[t] q\colon x& = 1, &
\\y & = -1 + 5t,
\\z & = 2 - 2t;\ t\in \mathbb{R}
\\ \end{aligned}\)}{\(\begin{aligned}[t] p\colon x& = -5 - t,&
\\y & = 2 - 2t,
\\z & = 2 + t;\ t\in \mathbb{R}
\\ \end{aligned}\)}{}{}{}}
\LANGcs{
\Question{
Pro kterou z následujících přímek platí, že se jedná o přímku rovnoběžnou s
přímkou \(s\) a vzdálenost mezi oběma přímkami je \(\sqrt{5}\)?
\[ 
\begin{aligned}[t] s\colon x& = -1 + t,&
\\y & = 2t,
\\z & = 2 - t;\ t\in \mathbb{R}
\\ \end{aligned}
\]}
{\(\begin{aligned}[t] r\colon x& = 3 - 2t,&
\\y & = 3 - 4t,
\\z & = 2t;\ t\in \mathbb{R}
\\ \end{aligned}\)}{\(\begin{aligned}[t] q\colon x& = 1, &
\\y & = -1 + 5t,
\\z & = 2 - 2t;\ t\in \mathbb{R}
\\ \end{aligned}\)}{\(\begin{aligned}[t] p\colon x& = -5 - t,&
\\y & = 2 - 2t,
\\z & = 2 + t;\ t\in \mathbb{R}
\\ \end{aligned}\)}{}{}{}}
\LANGsk{
\Question{
Pre ktorú z nasledujúcich priamok platí, že sa jedná o priamku rovnobežnú s
priamkou \(s\) a vzdialenosť medzi oboma priamkami je \(\sqrt{5}\)?
\[ 
\begin{aligned}[t] s\colon x& = -1 + t,&
\\y & = 2t,
\\z & = 2 - t;\ t\in \mathbb{R}
\\ \end{aligned}
\]}
{\(\begin{aligned}[t] r\colon x& = 3 - 2t,&
\\y & = 3 - 4t,
\\z & = 2t;\ t\in \mathbb{R}
\\ \end{aligned}\)}{\(\begin{aligned}[t] q\colon x& = 1, &
\\y & = -1 + 5t,
\\z & = 2 - 2t;\ t\in \mathbb{R}
\\ \end{aligned}\)}{\(\begin{aligned}[t] p\colon x& = -5 - t,&
\\y & = 2 - 2t,
\\z & = 2 + t;\ t\in \mathbb{R}
\\ \end{aligned}\)}{}{}{}}
\LANGes{
\Question{
Identifica la recta paralela con la recta \(s\), y cuya distancia entre las dos paralelas es igual a \(\sqrt{5}\).
\[ 
\begin{aligned}[t] s\colon x& = -1 + t,&
\\y & = 2t,
\\z & = 2 - t;\ t\in \mathbb{R}
\\ \end{aligned}
\]}
{\(\begin{aligned}[t] r\colon x& = 3 - 2t,&
\\y & = 3 - 4t,
\\z & = 2t;\ t\in \mathbb{R}
\\ \end{aligned}\)}{\(\begin{aligned}[t] q\colon x& = 1, &
\\y & = -1 + 5t,
\\z & = 2 - 2t;\ t\in \mathbb{R}
\\ \end{aligned}\)}{\(\begin{aligned}[t] p\colon x& = -5 - t,&
\\y & = 2 - 2t,
\\z & = 2 + t;\ t\in \mathbb{R}
\\ \end{aligned}\)}{}{}{}}
\End

\Data\ID{9000111805}
\Updated{2020-07-15 03:07:37}
\Level{B}
\SubArea{Analytic geometry in a space}
\LANGen{
\Question{
In the following list identify a plane which is parallel to the plane
\(\delta \) and the distance
between both planes is \(2\).
\[ 
 \delta \colon x - 2y + 2y - 2 = 0
\]}
{\(\begin{aligned}[t] \beta \colon x& = -4 + 2s, &
\\y& = 1 + r + s,
\\z& = 1 + r;\ r,s\in \mathbb{R}
\\ \end{aligned}\)}{\(\gamma \colon - x + 2y - 2z - 2 = 0\)}{\(\alpha \colon 2x - 4y + z - 4 = 0\)}{}{}{}}
\LANGpl{
\Question{
Wskaż płaszczyznę tak, aby była równoległa do płaszczyzny
\(\delta \), a odległość pomiędzy obiema 
płaszczyznami wynosiła \(2\).
\[ 
 \delta \colon x - 2y + 2y - 2 = 0
\]}
{\(\begin{aligned}[t] \beta \colon x& = -4 + 2s, &
\\y& = 1 + r + s,
\\z& = 1 + r;\  r,s\in \mathbb{R}
\\ \end{aligned}\)}{\(\gamma \colon - x + 2y - 2z - 2 = 0\)}{\(\alpha \colon 2x - 4y + z - 4 = 0\)}{}{}{}}
\LANGcs{
\Question{
Pro kterou z následujících rovin platí, že její vzdálenost od roviny 
dané obecnou rovnicí 
\[ 
 \delta \colon x - 2y + 2y - 2 = 0
\]
je rovna \(2\)?}
{\(\begin{aligned}[t] \beta \colon x& = -4 + 2s, &
\\y& = 1 + r + s,
\\z& = 1 + r;\ r,s\in \mathbb{R}
\\ \end{aligned}\)}{\(\gamma \colon - x + 2y - 2z - 2 = 0\)}{\(\alpha \colon 2x - 4y + z - 4 = 0\)}{}{}{}}
\LANGsk{
\Question{
Pre ktorú z nasledujúcich rovín platí, že jej vzdialenosť od roviny 
danej všeobecnou rovnicou 
\[ 
 \delta \colon x - 2y + 2y - 2 = 0
\]
je rovná \(2\)?}
{\(\begin{aligned}[t] \beta \colon x& = -4 + 2s, &
\\y& = 1 + r + s,
\\z& = 1 + r;\ r,s\in \mathbb{R}
\\ \end{aligned}\)}{\(\gamma \colon - x + 2y - 2z - 2 = 0\)}{\(\alpha \colon 2x - 4y + z - 4 = 0\)}{}{}{}}
\LANGes{
\Question{
Identifica el plano que es paralelo con el plano 
\(\delta \) y cuya distancia
entre los dos planos es igual a \(2\).
\[ 
 \delta \colon x - 2y + 2y - 2 = 0
\]}
{\(\begin{aligned}[t] \beta \colon x& = -4 + 2s, &
\\y& = 1 + r + s,
\\z& = 1 + r;\ r,s\in \mathbb{R}
\\ \end{aligned}\)}{\(\gamma \colon - x + 2y - 2z - 2 = 0\)}{\(\alpha \colon 2x - 4y + z - 4 = 0\)}{}{}{}}
\End

\Data\ID{9000111806}
\Updated{2020-07-15 03:07:37}
\Level{B}
\SubArea{Analytic geometry in a space}
\LANGen{
\Question{
In the following list identify the line such that the angle between this line and the
line \(s\) is
\(60^{\circ }\).
\[ 
\begin{aligned}[t] s\colon x& = 2 + t, &
\\y & = -1 - 2t,
\\z & = 3 - t;\ t\in \mathbb{R}
\\ \end{aligned}
\]}
{\(\begin{aligned}[t] r\colon x& = t, &
\\y & = -3 + t,
\\z & = 1 + 2t;\ t\in \mathbb{R}
\\ \end{aligned}\)}{\(\begin{aligned}[t] q\colon x& = 1, &
\\y & = -1 - t,
\\z & = 3 + 2t;\ t\in \mathbb{R}
\\ \end{aligned}\)}{\(\begin{aligned}[t] p\colon x& = -5 - 2t,&
\\y & = 2 + 4t,
\\z & = 2 + 2t;\ t\in \mathbb{R}
\\ \end{aligned}\)}{}{}{}}
\LANGpl{
\Question{
Wyznacz prostą tak, aby kąt pomiędzy tą prostą a prostą  
 \(s\) wynosił
\(60^{\circ }\).
\[ 
\begin{aligned}[t] s\colon x& = 2 + t, &
\\y & = -1 - 2t,
\\z & = 3 - t;\ t\in \mathbb{R}
\\ \end{aligned}
\]}
{\(\begin{aligned}[t] r\colon x& = t, &
\\y & = -3 + t,
\\z & = 1 + 2t;\  t\in \mathbb{R}
\\ \end{aligned}\)}{\(\begin{aligned}[t] q\colon x& = 1, &
\\y & = -1 - t,
\\z & = 3 + 2t;\ t\in \mathbb{R}
\\ \end{aligned}\)}{\(\begin{aligned}[t] p\colon x& = -5 - 2t,&
\\y & = 2 + 4t,
\\z & = 2 + 2t;\ t\in \mathbb{R}
\\ \end{aligned}\)}{}{}{}}
\LANGcs{
\Question{
Pro kterou z následujících přímek platí, že její odchylka od přímky \(s\) je 
rovna \(60^{\circ}\)?
\[ 
\begin{aligned}[t] s\colon x& = 2 + t, &
\\y & = -1 - 2t,
\\z & = 3 - t;\ t\in \mathbb{R}
\\ \end{aligned}
\]}
{\(\begin{aligned}[t] r\colon x& = t, &
\\y & = -3 + t,
\\z & = 1 + 2t;\ t\in \mathbb{R}
\\ \end{aligned}\)}{\(\begin{aligned}[t] q\colon x& = 1, &
\\y & = -1 - t,
\\z & = 3 + 2t;\ t\in \mathbb{R}
\\ \end{aligned}\)}{\(\begin{aligned}[t] p\colon x& = -5 - 2t,&
\\y & = 2 + 4t,
\\z & = 2 + 2t;\ t\in \mathbb{R}
\\ \end{aligned}\)}{}{}{}}
\LANGsk{
\Question{
Pre ktorú z nasledujúcich priamok platí, že jej odchýlka od priamky \(s\) je 
rovná \(60^{\circ}\)?
\[ 
\begin{aligned}[t] s\colon x& = 2 + t, &
\\y & = -1 - 2t,
\\z & = 3 - t;\ t\in \mathbb{R}
\\ \end{aligned}
\]}
{\(\begin{aligned}[t] r\colon x& = t, &
\\y & = -3 + t,
\\z & = 1 + 2t;\ t\in \mathbb{R}
\\ \end{aligned}\)}{\(\begin{aligned}[t] q\colon x& = 1, &
\\y & = -1 - t,
\\z & = 3 + 2t;\ t\in \mathbb{R}
\\ \end{aligned}\)}{\(\begin{aligned}[t] p\colon x& = -5 - 2t,&
\\y & = 2 + 4t,
\\z & = 2 + 2t;\ t\in \mathbb{R}
\\ \end{aligned}\)}{}{}{}}
\LANGes{
\Question{
Identifica la recta cuyo ángulo entre la recta \(s\) es igual a
\(60^{\circ }\).
\[ 
\begin{aligned}[t] s\colon x& = 2 + t, &
\\y & = -1 - 2t,
\\z & = 3 - t;\ t\in \mathbb{R}
\\ \end{aligned}
\]}
{\(\begin{aligned}[t] r\colon x& = t, &
\\y & = -3 + t,
\\z & = 1 + 2t;\ t\in \mathbb{R}
\\ \end{aligned}\)}{\(\begin{aligned}[t] q\colon x& = 1, &
\\y & = -1 - t,
\\z & = 3 + 2t;\ t\in \mathbb{R}
\\ \end{aligned}\)}{\(\begin{aligned}[t] p\colon x& = -5 - 2t,&
\\y & = 2 + 4t,
\\z & = 2 + 2t;\ t\in \mathbb{R}
\\ \end{aligned}\)}{}{}{}}
\End

\Data\ID{9000111807}
\Updated{2020-07-15 03:07:37}
\Level{B}
\SubArea{Analytic geometry in a space}
\LANGen{
\Question{
In the following list identify a line such that the angle between this line and the
plane 
\[ 
 2x - y + 3z - 5 = 0
\]
is \(30^{\circ }\).}
{\(\begin{aligned}[t] p\colon x& = 2 + t, &
\\y & = 1 + 3t,
\\z & = -2t;\ t\in \mathbb{R}
\\ \end{aligned}\)}{\(\begin{aligned}[t] r\colon x& = -2t, &
\\y & = -3 + t,
\\z & = 1 - 3t;\ t\in \mathbb{R}
\\ \end{aligned}\)}{\(\begin{aligned}[t] q\colon x& = 2 + 3t, &
\\y & = 3 - 2t,
\\z & = 3 + t;\ t\in \mathbb{R}
\\ \end{aligned}\)}{}{}{}}
\LANGpl{
\Question{
Wyznacz prostą tak, aby kąt pomiędzy tą prostą a płaszczyzną 
\[ 
 2x - y + 3z - 5 = 0
\]
wynosił  \(30^{\circ }\).}
{\(\begin{aligned}[t] p\colon x& = 2 + t, &
\\y & = 1 + 3t,
\\z & = -2t;\ t\in \mathbb{R}
\\ \end{aligned}\)}{\(\begin{aligned}[t] r\colon x& = -2t, &
\\y & = -3 + t,
\\z & = 1 - 3t;\ t\in \mathbb{R}
\\ \end{aligned}\)}{\(\begin{aligned}[t] q\colon x& = 2 + 3t, &
\\y & = 3 - 2t,
\\z & = 3 + t;\ t\in \mathbb{R}
\\ \end{aligned}\)}{}{}{}}
\LANGcs{
\Question{
Pro kterou z následujících přímek 
platí, že její odchylka od roviny dané obecnou rovnicí 
\[ 
 2x - y + 3z - 5 = 0
\]
je rovna \(30^{\circ }\)?}
{\(\begin{aligned}[t] p\colon x& = 2 + t, &
\\y & = 1 + 3t,
\\z & = -2t;\ t\in \mathbb{R}
\\ \end{aligned}\)}{\(\begin{aligned}[t] r\colon x& = -2t, &
\\y & = -3 + t,
\\z & = 1 - 3t;\ t\in \mathbb{R}
\\ \end{aligned}\)}{\(\begin{aligned}[t] q\colon x& = 2 + 3t, &
\\y & = 3 - 2t,
\\z & = 3 + t;\ t\in \mathbb{R}
\\ \end{aligned}\)}{}{}{}}
\LANGsk{
\Question{
Pre ktorú z nasledujúcich priamok 
platí, že jej odchýlka od roviny danej všeobecnou rovnicou 
\[ 
 2x - y + 3z - 5 = 0
\]
je rovná \(30^{\circ }\)?}
{\(\begin{aligned}[t] p\colon x& = 2 + t, &
\\y & = 1 + 3t,
\\z & = -2t;\ t\in \mathbb{R}
\\ \end{aligned}\)}{\(\begin{aligned}[t] r\colon x& = -2t, &
\\y & = -3 + t,
\\z & = 1 - 3t;\ t\in \mathbb{R}
\\ \end{aligned}\)}{\(\begin{aligned}[t] q\colon x& = 2 + 3t, &
\\y & = 3 - 2t,
\\z & = 3 + t;\ t\in \mathbb{R}
\\ \end{aligned}\)}{}{}{}}
\LANGes{
\Question{
Identifica la recta cuyo ángulo al plano  
\[ 
 2x - y + 3z - 5 = 0
\]
es igual a \(30^{\circ }\).}
{\(\begin{aligned}[t] p\colon x& = 2 + t, &
\\y & = 1 + 3t,
\\z & = -2t;\ t\in \mathbb{R}
\\ \end{aligned}\)}{\(\begin{aligned}[t] r\colon x& = -2t, &
\\y & = -3 + t,
\\z & = 1 - 3t;\ t\in \mathbb{R}
\\ \end{aligned}\)}{\(\begin{aligned}[t] q\colon x& = 2 + 3t, &
\\y & = 3 - 2t,
\\z & = 3 + t;\ t\in \mathbb{R}
\\ \end{aligned}\)}{}{}{}}
\End

\Data\ID{9000111808}
\Updated{2020-07-15 03:07:37}
\Level{B}
\SubArea{Analytic geometry in a space}
\LANGen{
\Question{
In the following list identify a plane such that the angle between this plane and the
plane \(\rho \) 
is \(45^{\circ }\).
\[ 
\rho \colon \begin{aligned}[t] x& = 1 + r - 2s, &
\\y& = 3 - r + 2s,
\\z& = -5 - 4r;\  r,\; s\in \mathbb{R}
\\ \end{aligned}
\]}
{\(\gamma \colon 3x - 2 = 0\)}{\(\beta \colon 2z - 2 = 0\)}{\(\alpha \colon x + y - 2 = 0\)}{}{}{}}
\LANGpl{
\Question{
Wyznacz płaszczyznę tak, aby kąt pomiędzy tą płaszczyzną a płaszczyzną  
\(\rho \) 
wynosił \(45^{\circ }\).
\[ 
\rho \colon \begin{aligned}[t] x& = 1 + r - 2s, &
\\y& = 3 - r + 2s,
\\z& = -5 - 4r;\  r,\; s\in \mathbb{R}
\\ \end{aligned}
\]}
{\(\gamma \colon 3x - 2 = 0\)}{\(\beta \colon 2z - 2 = 0\)}{\(\alpha \colon x + y - 2 = 0\)}{}{}{}}
\LANGcs{
\Question{
Pro kterou z následujících rovin platí, že její odchylka od roviny 
\[ 
\rho \colon \begin{aligned}[t] x& = 1 + r - 2s, &
\\y& = 3 - r + 2s,
\\z& = -5 - 4r;\  r,\; s\in \mathbb{R}
\\ \end{aligned}
\]
je rovna \(45^{\circ }\)?}
{\(\gamma \colon 3x - 2 = 0\)}{\(\beta \colon 2z - 2 = 0\)}{\(\alpha \colon x + y - 2 = 0\)}{}{}{}}
\LANGsk{
\Question{
Pre ktorú z nasledujúcich rovín platí, že jej odchýlka od roviny 
\[ 
\rho \colon \begin{aligned}[t] x& = 1 + r - 2s, &
\\y& = 3 - r + 2s,
\\z& = -5 - 4r;\  r,\; s\in \mathbb{R}
\\ \end{aligned}
\]
je rovná \(45^{\circ }\)?}
{\(\gamma \colon 3x - 2 = 0\)}{\(\beta \colon 2z - 2 = 0\)}{\(\alpha \colon x + y - 2 = 0\)}{}{}{}}
\LANGes{
\Question{
Identifica el plano cuyo ángulo entre el plano \(\rho \) 
es igual a \(45^{\circ }\).
\[ 
\rho \colon \begin{aligned}[t] x& = 1 + r - 2s, &
\\y& = 3 - r + 2s,
\\z& = -5 - 4r;\  r,\; s\in \mathbb{R}
\\ \end{aligned}
\]}
{\(\gamma \colon 3x - 2 = 0\)}{\(\beta \colon 2z - 2 = 0\)}{\(\alpha \colon x + y - 2 = 0\)}{}{}{}}
\End

\Data\ID{9000117401}
\Updated{2020-07-15 03:07:37}
\Level{B}
\SubArea{Analytic geometry in a space}
\LANGen{
\Question{
Find the intersection of the planes \(\rho \) 
and \(\sigma \).
\[\begin{aligned}
 \rho \colon 2x - 5y + 4z - 10 = 0,\qquad \sigma \colon x - y - z - 2 = 0 & &
\end{aligned}\]}
{\(\begin{aligned}[t] p\colon x& = 3t, &
\\y & = -2 + 2t,
\\z & = t;\  t\in \mathbb{R}
\\ \end{aligned}\)}{\(\begin{aligned}[t] q\colon x& = 2s - 10,&
\\y & = 5s - 10,
\\z & = s;\  s\in \mathbb{R}
\\ \end{aligned}\)}{\(\begin{aligned}[t] a\colon x& = 2u - 4,&
\\y & = 2u - 4,
\\z & = u;\  u\in \mathbb{R}
\\ \end{aligned}\)}{\(\begin{aligned}[t] b\colon x& = 3v + 1,&
\\y & = v - 2,
\\z & = v;\ v\in \mathbb{R}
\\ \end{aligned}\)}{}{}}
\LANGpl{
\Question{
Określ część wspólną przecięcia płaszczyzn \(\rho \) 
i \(\sigma \).
\[\begin{aligned}
 \rho \colon 2x - 5y + 4z - 10 = 0,\qquad \sigma \colon x - y - z - 2 = 0 & &
\end{aligned}\]}
{\(\begin{aligned}[t] p\colon x& = 3t, &
\\y & = -2 + 2t,
\\z & = t;\  t\in \mathbb{R}
\\ \end{aligned}\)}{\(\begin{aligned}[t] q\colon x& = 2s - 10,&
\\y & = 5s - 10,
\\z & = s;\  s\in \mathbb{R}
\\ \end{aligned}\)}{\(\begin{aligned}[t] a\colon x& = 2u - 4,&
\\y & = 2u - 4,
\\z & = u;\ u\in \mathbb{R}
\\ \end{aligned}\)}{\(\begin{aligned}[t] b\colon x& = 3v + 1,&
\\y & = v - 2,
\\z & = v;\ v\in \mathbb{R}
\\ \end{aligned}\)}{}{}}
\LANGcs{
\Question{
Jsou dány roviny
\[\begin{aligned}
 \rho \colon 2x - 5y + 4z - 10 = 0,\quad \sigma \colon x - y - z - 2 = 0. & &
\end{aligned}\]
Která z uvedených přímek je průsečnicí zadaných rovin?}
{\(p = \{[3t;-2 + 2t;t],\ t\in \mathbb{R}\}\)}{\(q = \{[2s - 10;5s - 10;s],\ s\in \mathbb{R}\}\)}{\(a = \{[2u - 4;2u - 4;u],\ u\in \mathbb{R}\}\)}{\(b = \{[3v + 1;v - 2;v],\ v\in \mathbb{R}\}\)}{}{}}
\LANGsk{
\Question{
Sú dané roviny
\[\begin{aligned}
 \rho \colon 2x - 5y + 4z - 10 = 0,\quad \sigma \colon x - y - z - 2 = 0. & &
\end{aligned}\]
Ktorá z uvedených priamok je priesečnicou zadaných rovín?}
{\(p = \{[3t;-2 + 2t;t],\ t\in \mathbb{R}\}\)}{\(q = \{[2s - 10;5s - 10;s],\ s\in \mathbb{R}\}\)}{\(a = \{[2u - 4;2u - 4;u],\ u\in \mathbb{R}\}\)}{\(b = \{[3v + 1;v - 2;v],\ v\in \mathbb{R}\}\)}{}{}}
\LANGes{
\Question{
Halla la intersección de los planos \(\rho \) 
y \(\sigma \).
\[\begin{aligned}
 \rho \colon 2x - 5y + 4z - 10 = 0,\qquad \sigma \colon x - y - z - 2 = 0 & &
\end{aligned}\]}
{\(\begin{aligned}[t] p\colon x& = 3t, &
\\y & = -2 + 2t,
\\z & = t;\  t\in \mathbb{R}
\\ \end{aligned}\)}{\(\begin{aligned}[t] q\colon x& = 2s - 10,&
\\y & = 5s - 10,
\\z & = s;\  s\in \mathbb{R}
\\ \end{aligned}\)}{\(\begin{aligned}[t] a\colon x& = 2u - 4,&
\\y & = 2u - 4,
\\z & = u;\  u\in \mathbb{R}
\\ \end{aligned}\)}{\(\begin{aligned}[t] b\colon x& = 3v + 1,&
\\y & = v - 2,
\\z & = v;\ v\in \mathbb{R}
\\ \end{aligned}\)}{}{}}
\End

\Data\ID{9000117407}
\Updated{2020-07-15 03:07:37}
\Level{B}
\SubArea{Analytic geometry in a space}
\LANGen{
\Question{
Determine the value of the real parameter
\(p\) which
ensures that the following planes are perpendicular.
\[\begin{aligned}
 \rho \colon 2x - 4y + 5z - 4 = 0,\qquad \sigma \colon - 3x + py - 2z + 4 = 0 & &
\end{aligned}\]}
{\(p = -4\)}{\(p = 4\)}{\(p = 0\)}{\(p = -3\)}{}{}}
\LANGpl{
\Question{
Wyznacz wartość rzeczywistą parametru 
\(p\) tak, aby
dane płaszczyzny były prostopadłe.
\[\begin{aligned}
 \rho \colon 2x - 4y + 5z - 4 = 0,\qquad \sigma \colon - 3x + py - 2z + 4 = 0 & &
\end{aligned}\]}
{\(p = -4\)}{\(p = 4\)}{\(p = 0\)}{\(p = -3\)}{}{}}
\LANGcs{
\Question{
Jaká musí být hodnota reálného parametru
\(p\), aby
roviny
\[\begin{aligned}
 \rho \colon 2x - 4y + 5z - 4 = 0,\quad \sigma \colon - 3x + py - 2z + 4 = 0 & &
\end{aligned}\]
byly navzájem kolmé?}
{\(p = -4\)}{\(p = 4\)}{\(p = 0\)}{\(p = -3\)}{}{}}
\LANGsk{
\Question{
Aká musí byť hodnota reálneho parametra
\(p\), aby
roviny
\[\begin{aligned}
 \rho \colon 2x - 4y + 5z - 4 = 0,\quad \sigma \colon - 3x + py - 2z + 4 = 0 & &
\end{aligned}\]
boli navzájom kolmé?}
{\(p = -4\)}{\(p = 4\)}{\(p = 0\)}{\(p = -3\)}{}{}}
\LANGes{
\Question{
Determina el valor del parametro real 
\(p\) para que los siguientes planos sean perpendiculares.
\[\begin{aligned}
 \rho \colon 2x - 4y + 5z - 4 = 0,\qquad \sigma \colon - 3x + py - 2z + 4 = 0 & &
\end{aligned}\]}
{\(p = -4\)}{\(p = 4\)}{\(p = 0\)}{\(p = -3\)}{}{}}
\End

\Data\ID{9000117408}
\Updated{2020-07-15 03:07:37}
\Level{B}
\SubArea{Analytic geometry in a space}
\LANGen{
\Question{
In the following list find the plane perpendicular to the plane
\(\rho \).
\[\begin{aligned}
 \rho \colon 2x - 3y + 7z - 2 = 0 & &
\end{aligned}\]}
{\(\omega \colon x + 3y + z + 7 = 0\)}{\(\tau \colon - 2x + 3y - 7z + 2 = 0\)}{\(\nu \colon - 2x - 3y + 7z + 2 = 0\)}{\(\sigma \colon 7x - 3y + 2z - 2 = 0\)}{}{}}
\LANGpl{
\Question{
Wskaż płaszczyznę prostopadłą do płaszczyzny
\(\rho \).
\[\begin{aligned}
 \rho \colon 2x - 3y + 7z - 2 = 0 & &
\end{aligned}\]}
{\(\omega \colon x + 3y + z + 7 = 0\)}{\(\tau \colon - 2x + 3y - 7z + 2 = 0\)}{\(\nu \colon - 2x - 3y + 7z + 2 = 0\)}{\(\sigma \colon 7x - 3y + 2z - 2 = 0\)}{}{}}
\LANGcs{
\Question{
Je dána rovina
\[\begin{aligned}
 \rho \colon 2x - 3y + 7z - 2 = 0. & &
\end{aligned}\]
Která z uvedených rovin je kolmá k rovině
\(\rho \)?}
{\(\omega \colon x + 3y + z + 7 = 0\)}{\(\tau \colon - 2x + 3y - 7z + 2 = 0\)}{\(\nu \colon - 2x - 3y + 7z + 2 = 0\)}{\(\sigma \colon 7x - 3y + 2z - 2 = 0\)}{}{}}
\LANGsk{
\Question{
Je daná rovina
\[\begin{aligned}
 \rho \colon 2x - 3y + 7z - 2 = 0. & &
\end{aligned}\]
Ktorá z uvedených rovín je kolmá k rovine
\(\rho \)?}
{\(\omega \colon x + 3y + z + 7 = 0\)}{\(\tau \colon - 2x + 3y - 7z + 2 = 0\)}{\(\nu \colon - 2x - 3y + 7z + 2 = 0\)}{\(\sigma \colon 7x - 3y + 2z - 2 = 0\)}{}{}}
\LANGes{
\Question{
Halla el plano perpendicular al plano \(\rho \).
\[\begin{aligned}
 \rho \colon 2x - 3y + 7z - 2 = 0 & &
\end{aligned}\]}
{\(\omega \colon x + 3y + z + 7 = 0\)}{\(\tau \colon - 2x + 3y - 7z + 2 = 0\)}{\(\nu \colon - 2x - 3y + 7z + 2 = 0\)}{\(\sigma \colon 7x - 3y + 2z - 2 = 0\)}{}{}}
\End

\Data\ID{9000117409}
\Updated{2020-07-15 03:07:37}
\Level{B}
\SubArea{Analytic geometry in a space}
\LANGen{
\Question{
Find the plane parallel to \(\rho \) 
passing through the point \(M\).
\[\begin{aligned}
 \rho \colon x - 2y + 5z - 3 = 0,\qquad M = [3;-1;1] & &
\end{aligned}\]}
{\(\tau \colon x - 2y + 5z - 10 = 0\)}{\(\sigma \colon 3x - y + z - 3 = 0\)}{\(\nu \colon x - 2y + 5z + 1 = 0\)}{\(\omega \colon 3x - y + z - 11 = 0\)}{}{}}
\LANGpl{
\Question{
Wskaż płaszczyznę równoległą do \(\rho \) 
i przechodzącą przez punkt \(M\).
\[\begin{aligned}
 \rho \colon x - 2y + 5z - 3 = 0,\qquad M = [3;-1;1] & &
\end{aligned}\]}
{\(\tau \colon x - 2y + 5z - 10 = 0\)}{\(\sigma \colon 3x - y + z - 3 = 0\)}{\(\nu \colon x - 2y + 5z + 1 = 0\)}{\(\omega \colon 3x - y + z - 11 = 0\)}{}{}}
\LANGcs{
\Question{
Je dána rovina
\[\begin{aligned}
 \rho \colon x - 2y + 5z - 3 = 0 & &
\end{aligned}\]
a bod \(M = [3;-1;1]\). Vyberte, která z uvedených rovin
prochází bodem \(M\) a je
rovnoběžná s rovinou \(\rho \).}
{\(\tau \colon x - 2y + 5z - 10 = 0\)}{\(\sigma \colon 3x - y + z - 3 = 0\)}{\(\nu \colon x - 2y + 5z + 1 = 0\)}{\(\omega \colon 3x - y + z - 11 = 0\)}{}{}}
\LANGsk{
\Question{
Je daná rovina
\[\begin{aligned}
 \rho \colon x - 2y + 5z - 3 = 0 & &
\end{aligned}\]
a bod \(M = [3;-1;1]\). Vyberte, ktorá z uvedených rovín 
prechádza bodom \(M\) a je
rovnobežná s rovinou \(\rho \).}
{\(\tau \colon x - 2y + 5z - 10 = 0\)}{\(\sigma \colon 3x - y + z - 3 = 0\)}{\(\nu \colon x - 2y + 5z + 1 = 0\)}{\(\omega \colon 3x - y + z - 11 = 0\)}{}{}}
\LANGes{
\Question{
Halla el plano paralelo al plano \(\rho \) 
y que pasa por el punto \(M\).
\[\begin{aligned}
 \rho \colon x - 2y + 5z - 3 = 0,\qquad M = [3;-1;1] & &
\end{aligned}\]}
{\(\tau \colon x - 2y + 5z - 10 = 0\)}{\(\sigma \colon 3x - y + z - 3 = 0\)}{\(\nu \colon x - 2y + 5z + 1 = 0\)}{\(\omega \colon 3x - y + z - 11 = 0\)}{}{}}
\End

\Data\ID{9000117410}
\Updated{2020-07-15 03:07:37}
\Level{B}
\SubArea{Analytic geometry in a space}
\LANGen{
\Question{
Adjust real parameters \(p\) 
and \(q\) to
ensure that \(\rho \) 
and \(\sigma \) 
are parallel but not identical planes.
\[\begin{aligned}
 \rho \colon 2x - 3y + 5z + 6 = 0,\qquad \sigma \colon 4x + py + qz - 2 = 0 & &
\end{aligned}\]}
{\(p = -6;\ q = 10\)}{\(p = 6;\ q = 10\)}{\(p = 6;\ q = -10\)}{\(p = -6;\ q = -10\)}{}{}}
\LANGpl{
\Question{
Określ wartość rzeczywistą parametrów \(p\) 
i \(q\) tak, 
aby \(\rho \) 
i \(\sigma \) 
były płaszczyznami równoległymi, nie pokrywającymi się.
\[\begin{aligned}
 \rho \colon 2x - 3y + 5z + 6 = 0,\qquad \sigma \colon 4x + py + qz - 2 = 0 & &
\end{aligned}\]}
{\(p = -6;\ q = 10\)}{\(p = 6;\ q = 10\)}{\(p = 6;\ q = -10\)}{\(p = -6;\ q = -10\)}{}{}}
\LANGcs{
\Question{
Jaká musí být hodnota reálných parametrů
\(p,\ q\), aby
roviny
\[\begin{aligned}
 \rho \colon 2x - 3y + 5z + 6 = 0,\quad \sigma \colon 4x + py + qz - 2 = 0 & &
\end{aligned}\]
byly rovnoběžné různé?}
{\(- 6;10\)}{\(6;10\)}{\(6;-10\)}{\(- 6;-10\)}{}{}}
\LANGsk{
\Question{
Aká musí byť hodnota reálnych parametrov
\(p,\ q\), aby
roviny
\[\begin{aligned}
 \rho \colon 2x - 3y + 5z + 6 = 0,\quad \sigma \colon 4x + py + qz - 2 = 0 & &
\end{aligned}\]
boli rovnobežné rôzne?}
{\(- 6;10\)}{\(6;10\)}{\(6;-10\)}{\(- 6;-10\)}{}{}}
\LANGes{
\Question{
Ajusta los parametros reales \(p\) 
y \(q\) para que los planos \(\rho \) 
y \(\sigma \) 
sean paralelas no idénticas.
\[\begin{aligned}
 \rho \colon 2x - 3y + 5z + 6 = 0,\qquad \sigma \colon 4x + py + qz - 2 = 0 & &
\end{aligned}\]}
{\(p = -6;\ q = 10\)}{\(p = 6;\ q = 10\)}{\(p = 6;\ q = -10\)}{\(p = -6;\ q = -10\)}{}{}}
\End

\Data\ID{1003233605}
\Updated{2020-07-15 04:07:40}
\Level{C}
\SubArea{Analytic geometry in a space}
\LANGen{
\Question{
We are given skew lines $p$ and $q$.
\begin{align*}
p\colon  x&= 1-t, & q\colon  x&= 1-2s, \\
y&= 1+t, & y&=s, \\
z&= 3+2t;\  t\in\mathbb{R}, & z&= 3+3s;\  s\in\mathbb{R}.
\end{align*}
Find parametric equations of a straight line $r$, that is intersecting both lines $p$ and $q$ and lying in the plane $x+2y-z+2=0$.}
{$\begin{aligned}
r\colon x&=-1+2m, \\
y&=3-3m, \\
z&=7-4m;\ m\in\mathbb{R}
\end{aligned}$}{$\begin{aligned}
r\colon x&=-1+m, \\
y&=3+3m, \\
z&=7-m;\ m\in\mathbb{R}
\end{aligned}$}{$\begin{aligned}
r\colon x&=-1+3m, \\
y&=3+2m, \\
z&=7+5m;\ m\in\mathbb{R}
\end{aligned}$}{$\begin{aligned}
r\colon x&=-1+m, \\
y&=3-m, \\
z&=7+m;\ m\in\mathbb{R}
\end{aligned}$}{}{}}
\LANGpl{
\Question{
Dane są proste skośne $p$ i $q$.

\begin{align*}
p\colon  x&= 1-t, & q\colon  x&= 1-2s, \\
y&= 1+t, & y&=s, \\
z&= 3+2t;\  t\in\mathbb{R}, & z&= 3+3s;\  s\in\mathbb{R}.
\end{align*}
Wskaż równanie parametryczne prostej $r$, przecinającej obie proste $p$ i $q$ leżącej na płaszczyźnie $x+2y-z+2=0$.}
{$\begin{aligned}
r\colon x&=-1+2m, \\
y&=3-3m, \\
z&=7-4m;\ m\in\mathbb{R}
\end{aligned}$}{$\begin{aligned}
r\colon x&=-1+m, \\
y&=3+3m, \\
z&=7-m;\ m\in\mathbb{R}
\end{aligned}$}{$\begin{aligned}
r\colon x&=-1+3m, \\
y&=3+2m, \\
z&=7+5m;\ m\in\mathbb{R}
\end{aligned}$}{$\begin{aligned}
r\colon x&=-1+m, \\
y&=3-m, \\
z&=7+m;\ m\in\mathbb{R}
\end{aligned}$}{}{}}
\LANGcs{
\Question{
Jsou dány přímky $p$ a $q$.

\begin{align*}
p\colon  x&= 1-t, & q\colon  x&= 1-2s, \\
y&= 1+t, & y&=s, \\
z&= 3+2t;\  t\in\mathbb{R}, & z&= 3+3s;\  s\in\mathbb{R}.
\end{align*}
Určete parametrické rovnice přímky $r$, která obě přímky $p$ a $q$ protíná a leží přitom v rovině $x+2y-z+2=0$.}
{$\begin{aligned}
r\colon x&=-1+2m, \\
y&=3-3m, \\
z&=7-4m;\ m\in\mathbb{R}
\end{aligned}$}{$\begin{aligned}
r\colon x&=-1+m, \\
y&=3+3m, \\
z&=7-m;\ m\in\mathbb{R}
\end{aligned}$}{$\begin{aligned}
r\colon x&=-1+3m, \\
y&=3+2m, \\
z&=7+5m;\ m\in\mathbb{R}
\end{aligned}$}{$\begin{aligned}
r\colon x&=-1+m, \\
y&=3-m, \\
z&=7+m;\ m\in\mathbb{R}
\end{aligned}$}{}{}}
\LANGsk{
\Question{
Dané sú mimobežné priamky $p$ a $q$.

\begin{align*}
p\colon  x&= 1-t, & q\colon  x&= 1-2s, \\
y&= 1+t, & y&=s, \\
z&= 3+2t;\  t\in\mathbb{R}, & z&= 3+3s;\  s\in\mathbb{R}.
\end{align*}
Nájdite parametrické vyjadrenie priamky r, ktorá pretína obe priamky $p$ a $q$  a leží v rovine $x+2y-z+2=0$.}
{$\begin{aligned}
r\colon x&=-1+2m, \\
y&=3-3m, \\
z&=7-4m;\ m\in\mathbb{R}
\end{aligned}$}{$\begin{aligned}
r\colon x&=-1+m, \\
y&=3+3m, \\
z&=7-m;\ m\in\mathbb{R}
\end{aligned}$}{$\begin{aligned}
r\colon x&=-1+3m, \\
y&=3+2m, \\
z&=7+5m;\ m\in\mathbb{R}
\end{aligned}$}{$\begin{aligned}
r\colon x&=-1+m, \\
y&=3-m, \\
z&=7+m;\ m\in\mathbb{R}
\end{aligned}$}{}{}}
\LANGes{
\Question{
Dadas las rectas $p$ and $q$.
\begin{align*}
p\colon  x&= 1-t, & q\colon  x&= 1-2s, \\
y&= 1+t, & y&=s, \\
z&= 3+2t;\  t\in\mathbb{R}, & z&= 3+3s;\  s\in\mathbb{R}.
\end{align*}
Determina las ecuaciones paramétricas de la recta $r$, que pasa por las dos rectas  $p$ y $q$, al mismo tiempo, está en el plano $x+2y-z+2=0$.}
{$\begin{aligned}
r\colon x&=-1+2m, \\
y&=3-3m, \\
z&=7-4m;\ m\in\mathbb{R}
\end{aligned}$}{$\begin{aligned}
r\colon x&=-1+m, \\
y&=3+3m, \\
z&=7-m;\ m\in\mathbb{R}
\end{aligned}$}{$\begin{aligned}
r\colon x&=-1+3m, \\
y&=3+2m, \\
z&=7+5m;\ m\in\mathbb{R}
\end{aligned}$}{$\begin{aligned}
r\colon x&=-1+m, \\
y&=3-m, \\
z&=7+m;\ m\in\mathbb{R}
\end{aligned}$}{}{}}
\End

\Data\ID{1003233606}
\Updated{2020-07-15 04:07:40}
\Level{C}
\SubArea{Analytic geometry in a space}
\LANGen{
\Question{
We are given a triangle $ABC$, where $A=[1;2;3]$, $B= [3;6;2]$, and $C=[-1;10;-2]$. Find the  height to the side $BC$.}
{$3\sqrt2$}{$2\sqrt3$}{$2\sqrt6$}{$3\sqrt3$}{}{}}
\LANGpl{
\Question{
Dany jest trójkąt $ABC$, gdzie $A=[1;2;3]$, $B= [3;6;2]$, i $C=[-1;10;-2]$. Wyznacz wysokość do boku $BC$.}
{$3\sqrt2$}{$2\sqrt3$}{$2\sqrt6$}{$3\sqrt3$}{}{}}
\LANGcs{
\Question{
Je dán trojúhelník $ABC$, kde $A=[1;2;3]$, $B= [3;6;2]$ a $C=[-1;10;-2]$. Vypočtěte velikost výšky na stranu $BC$.}
{$3\sqrt2$}{$2\sqrt3$}{$2\sqrt6$}{$3\sqrt3$}{}{}}
\LANGsk{
\Question{
Je daný trojuholník $ABC$, kde $A=[1;2;3]$, $B= [3;6;2]$, a $C=[-1;10;-2]$. Vypočítaj výšku na stranu $BC$.}
{$3\sqrt2$}{$2\sqrt3$}{$2\sqrt6$}{$3\sqrt3$}{}{}}
\LANGes{
\Question{
Dado el triángulo $ABC$, donde $A=[1;2;3]$, $B= [3;6;2]$, y $C=[-1;10;-2]$. Determina la altura del lado $BC$.}
{$3\sqrt2$}{$2\sqrt3$}{$2\sqrt6$}{$3\sqrt3$}{}{}}
\End

\Data\ID{1003233607}
\Updated{2020-07-15 04:07:44}
\Level{C}
\SubArea{Analytic geometry in a space}
\LANGen{
\Question{
Determine the relative position of three planes:
\begin{align*}
\alpha\colon\ &2x+y+9z-18=0, \\
\beta\colon\ &x+3y+2z+16=0, \\
\gamma\colon\ &x+2y+3z+6=0.
\end{align*}}
{Planes $\alpha$, $\beta$ and $\gamma$ intersect in a straight line.}{Each of the two planes are intersecting and the lines of intersection are three different lines parallel  to each other.}{All three planes intersect at just one point.}{}{}{}}
\LANGpl{
\Question{
Określ wzajemne położenie trzech płaszczyzn:
\begin{align*}
\alpha\colon\ &2x+y+9z-18=0, \\
\beta\colon\ &x+3y+2z+16=0, \\
\gamma\colon\ &x+2y+3z+6=0.
\end{align*}}
{Płaszczyzny $\alpha$, $\beta$ i $\gamma$ przecinają się na prostej.}{Każda z dwóch płaszczyzn przecina się, a proste przecięcia to trzy różne proste równoległe do siebie.}{Trzy płaszczyzny przecinają się w jednym punkcie.}{}{}{}}
\LANGcs{
\Question{
Určete vzájemnou polohu tří rovin:
\begin{align*}
\alpha\colon\ &2x+y+9z-18=0, \\
\beta\colon\ &x+3y+2z+16=0, \\
\gamma\colon\ &x+2y+3z+6=0.
\end{align*}}
{Roviny $\alpha$, $\beta$ a $\gamma$ se protínají v přímce.}{Každá dvojice rovin se protíná v přímce. Tyto tři průsečnice jsou navzájem různé rovnoběžky.}{Všechny tři roviny se protínají v právě jednom bodě.}{}{}{}}
\LANGsk{
\Question{
Určte polohu troch rovín:
\begin{align*}
\alpha\colon\ &2x+y+9z-18=0, \\
\beta\colon\ &x+3y+2z+16=0, \\
\gamma\colon\ &x+2y+3z+6=0.
\end{align*}}
{Roviny $\alpha$, $\beta$ a $\gamma$ sa pretínajú v jednej priamke.}{Každá z týchto dvoch rovín sa pretína a priesečníkmi sú tri rôzne priamky, ktoré sú navzájom rovnobežné.}{Všetky tri roviny sa pretínajú len v jednom bode.}{}{}{}}
\LANGes{
\Question{
Determina la posición relativa de tres planos:
\begin{align*}
\alpha\colon\ &2x+y+9z-18=0, \\
\beta\colon\ &x+3y+2z+16=0, \\
\gamma\colon\ &x+2y+3z+6=0.
\end{align*}}
{Los planos $\alpha$, $\beta$ y $\gamma$ se cortan en una recta.}{Cada pareja de los planos se corta en la recta. Estas tres líneas de intersección son rectas no paralelas entre sí.}{Los tres planos se cortan justo en un punto.}{}{}{}}
\End

\Data\ID{1103212201}
\Updated{2020-07-15 03:07:21}
\Level{C}
\SubArea{Analytic geometry in a space}
\IMG{1103212201.pdf}
\LANGen{
\Question{
A straight line \( p \) is given by the points \( M=[4;2;0] \) and \( N=[6;6;7] \) (see the picture). Find the parametric equations of the line \( p' \) that is symmetrical to the line \( p \) in the plane symmetry across the coordinate \( xy \)-plane.}
{\( \begin{aligned}
p'\colon x&=4+2t, \\
y&=2+4t, \\
z&=-7t;\ t\in\mathbb{R}
\end{aligned} \)}{\( \begin{aligned}
p'\colon x&=4+6, \\
y&=2+6t, \\
z&=-7t;\ t\in\mathbb{R}
\end{aligned} \)}{\( \begin{aligned}
p'\colon x&=4+2t, \\
y&=2+4t, \\
z&=7t;\ t\in\mathbb{R}
\end{aligned} \)}{\( \begin{aligned}
p'\colon x&=4+6t, \\
y&=2+6t, \\
z&=7t;\ t\in\mathbb{R}
\end{aligned} \)}{}{}}
\LANGpl{
\Question{
Prosta \( p \) określona punktami \( M=[4;2;0] \) i \( N=[6;6;7] \) (spójrz na rysunek). Wskaż równanie parametryczne prostej  \( p' \), która jest symetryczna do prostej  \( p \) w płaszczyźnie symetrii w układzie współrzędnych  \( xy \).}
{\( \begin{aligned}
p'\colon x&=4+2t, \\
y&=2+4t, \\
z&=-7t;\ t\in\mathbb{R}
\end{aligned} \)}{\( \begin{aligned}
p'\colon x&=4+6t, \\
y&=2+6t, \\
z&=-7t;\ t\in\mathbb{R}
\end{aligned} \)}{\( \begin{aligned}
p'\colon x&=4+2t, \\
y&=2+4t, \\
z&=7t;\ t\in\mathbb{R}
\end{aligned} \)}{\( \begin{aligned}
p'\colon x&=4+6t, \\
y&=2+6t, \\
z&=7t;\ t\in\mathbb{R}
\end{aligned} \)}{}{}}
\LANGcs{
\Question{
Přímka \( p \) je zadána body \( M=[4;2;0] \) a \( N=[6;6;7] \) (viz obrázek). Určete parametrické rovnice přímky \( p' \), která je s přímkou \( p \) rovinově souměrná podle souřadné roviny \( (xy) \).}
{\( \begin{aligned}
p'\colon x&=4+2t, \\
y&=2+4t, \\
z&=-7t;\ t\in\mathbb{R}
\end{aligned} \)}{\( \begin{aligned}
p'\colon x&=4+6t, \\
y&=2+6t, \\
z&=-7t;\ t\in\mathbb{R}
\end{aligned} \)}{\( \begin{aligned}
p'\colon x&=4+2t, \\
y&=2+4t, \\
z&=7t;\ t\in\mathbb{R}
\end{aligned} \)}{\( \begin{aligned}
p'\colon x&=4+6t, \\
y&=2+6t, \\
z&=7t;\ t\in\mathbb{R}
\end{aligned} \)}{}{}}
\LANGsk{
\Question{
Priamka \( p \) je zadaná bodmi \( M=[4;2;0] \) a \( N=[6;6;7] \) (viď obrázok). Určte parametrické rovnice priamky \( p' \), ktorá je s priamkou \( p \) rovinovo súmerná podľa súradnicovej roviny \( (xy) \).}
{\( \begin{aligned}
p'\colon x&=4+2t, \\
y&=2+4t, \\
z&=-7t;\ t\in\mathbb{R}
\end{aligned} \)}{\( \begin{aligned}
p'\colon x&=4+6t, \\
y&=2+6t, \\
z&=-7t;\ t\in\mathbb{R}
\end{aligned} \)}{\( \begin{aligned}
p'\colon x&=4+2t, \\
y&=2+4t, \\
z&=7t;\ t\in\mathbb{R}
\end{aligned} \)}{\( \begin{aligned}
p'\colon x&=4+6t, \\
y&=2+6t, \\
z&=7t;\ t\in\mathbb{R}
\end{aligned} \)}{}{}}
\LANGes{
\Question{
Dada la recta \( p \) que pasa por los puntos \( M=[4;2;0] \) y \( N=[6;6;7] \) (mira la imagen). Determina las ecuaciones paramétricas de la recta \( p' \) que es simétrica a la recta \( p \) respecto al plano \( xy \).}
{\( \begin{aligned}
p'\colon x&=4+2t, \\
y&=2+4t, \\
z&=-7t;\ t\in\mathbb{R}
\end{aligned} \)}{\( \begin{aligned}
p'\colon x&=4+6, \\
y&=2+6t, \\
z&=-7t;\ t\in\mathbb{R}
\end{aligned} \)}{\( \begin{aligned}
p'\colon x&=4+2t, \\
y&=2+4t, \\
z&=7t;\ t\in\mathbb{R}
\end{aligned} \)}{\( \begin{aligned}
p'\colon x&=4+6t, \\
y&=2+6t, \\
z&=7t;\ t\in\mathbb{R}
\end{aligned} \)}{}{}}
\End

\Data\ID{1103212202}
\Updated{2020-07-15 03:07:21}
\Level{C}
\SubArea{Analytic geometry in a space}
\IMG{1103212202.pdf}
\LANGen{
\Question{
A straight line \( p \) is given by the points \( M=[4;3;2] \) and \( N=[0;6;7] \) (see the picture). Find the parametric equations of the line \( p' \) that is symmetrical to the line \( p \) in the plane symmetry across the coordinate \( yz \)-plane.}
{\( \begin{aligned}
p'\colon x&=4t, \\
y&=6+3t, \\
z&=7+5t;\ t\in\mathbb{R}
\end{aligned} \)}{\( \begin{aligned}
p'\colon x&=-4t, \\
y&=6+3t, \\
z&=7+5t;\ t\in\mathbb{R}
\end{aligned} \)}{\( \begin{aligned}
p'\colon x&=4t, \\
y&=6-3t, \\
z&=7+5t;\ t\in\mathbb{R}
\end{aligned} \)}{\( \begin{aligned}
p'\colon x&=-4t, \\
y&=6-3t, \\
z&=7+5t;\ t\in\mathbb{R}
\end{aligned} \)}{}{}}
\LANGpl{
\Question{
Dana jest prosta   \( p \) określona punktami \( M=[4;3;2] \) i \( N=[0;6;7] \) (spójrz na rysunek). Wskaż równanie parametryczne prostej \( p' \), która jest symetryczna do prostej  \( p \) w płaszczyźnie symetrii w układzie współrzędnych  \( yz \).}
{\( \begin{aligned}
p'\colon x&=4t, \\
y&=6+3t, \\
z&=7+5t;\ t\in\mathbb{R}
\end{aligned} \)}{\( \begin{aligned}
p'\colon x&=-4t, \\
y&=6+3t, \\
z&=7+5t;\ t\in\mathbb{R}
\end{aligned} \)}{\( \begin{aligned}
p'\colon x&=4t, \\
y&=6-3t, \\
z&=7+5t;\ t\in\mathbb{R}
\end{aligned} \)}{\( \begin{aligned}
p'\colon x&=-4t, \\
y&=6-3t, \\
z&=7+5t;\ t\in\mathbb{R}
\end{aligned} \)}{}{}}
\LANGcs{
\Question{
Přímka \( p \) je zadána body \( M=[4;3;2] \) a \( N=[0;6;7] \) (viz obrázek). Určete parametrické rovnice přímky \( p' \) která je souměrná s přímkou \( p \) v rovinové souměrnosti podle souřadné roviny \( (yz) \).}
{\( \begin{aligned}
p'\colon x&=4t, \\
y&=6+3t, \\
z&=7+5t;\ t\in\mathbb{R}
\end{aligned} \)}{\( \begin{aligned}
p'\colon x&=-4t,\\
y&=6+3t, \\
z&=7+5t;\ t\in\mathbb{R}
\end{aligned} \)}{\( \begin{aligned}
p'\colon x&=4t, \\
y&=6-3t, \\
z&=7+5t;\ t\in\mathbb{R}
\end{aligned} \)}{\( \begin{aligned}
p'\colon x&=-4t, \\
y&=6-3t, \\
z&=7+5t;\ t\in\mathbb{R}
\end{aligned} \)}{}{}}
\LANGsk{
\Question{
Priamka \( p \) je zadaná bodmi \( M=[4;3;2] \) a \( N=[0;6;7] \) (viď obrázok). Určte parametrické rovnice priamky \( p' \) ktorá je súmerná s priamkou \( p \) v rovinovej súmernosti podľa súradnicovej roviny \( (yz) \).}
{\( \begin{aligned}
p'\colon x&=4t, \\
y&=6+3t, \\
z&=7+5t;\ t\in\mathbb{R}
\end{aligned} \)}{\( \begin{aligned}
p'\colon x&=-4t, \\
y&=6+3t, \\
z&=7+5t;\ t\in\mathbb{R}
\end{aligned} \)}{\( \begin{aligned}
p'\colon x&=4t, \\
y&=6-3t, \\
z&=7+5t;\ t\in\mathbb{R}
\end{aligned} \)}{\( \begin{aligned}
p'\colon x&=-4t, \\
y&=6-3t, \\
z&=7+5t;\ t\in\mathbb{R}
\end{aligned} \)}{}{}}
\LANGes{
\Question{
Dada la recta  \( p \) que pasa por los puntos \( M=[4;3;2] \) y \( N=[0;6;7] \) (mira la imagen). Determina las ecuaciones paramétricas de la recta \( p' \) que es simétrica a la recta \( p \) respecto al \( yz \).}
{\( \begin{aligned}
p'\colon x&=4t, \\
y&=6+3t, \\
z&=7+5t;\ t\in\mathbb{R}
\end{aligned} \)}{\( \begin{aligned}
p'\colon x&=-4t, \\
y&=6+3t, \\
z&=7+5t;\ t\in\mathbb{R}
\end{aligned} \)}{\( \begin{aligned}
p'\colon x&=4t, \\
y&=6-3t, \\
z&=7+5t;\ t\in\mathbb{R}
\end{aligned} \)}{\( \begin{aligned}
p'\colon x&=-4t, \\
y&=6-3t, \\
z&=7+5t;\ t\in\mathbb{R}
\end{aligned} \)}{}{}}
\End

\Data\ID{1103212203}
\Updated{2020-07-15 03:07:21}
\Level{C}
\SubArea{Analytic geometry in a space}
\IMG{1103212203.pdf}
\LANGen{
\Question{
A straight line \( p \) is given by the points \( M=[4;3;2] \) and \( N=[8;0;5] \) (see the picture). Find the parametric equations of the line \( p' \) that is symmetrical to the line \( p \) in the plane symmetry across the coordinate \( xz \)-plane.}
{\( \begin{aligned}
p'\colon x&=8+4t, \\
y&=3t, \\
z&=5+3t;\ t\in\mathbb{R}
\end{aligned} \)}{\( \begin{aligned}
p'\colon x&=8+4t, \\
y&=0, \\
z&=5+3t;\ t\in\mathbb{R}
\end{aligned} \)}{\( \begin{aligned}
p'\colon x&=8+4t, \\
y&=-3t, \\
z&=5+3t;\ t\in\mathbb{R}
\end{aligned} \)}{\( \begin{aligned}
p'\colon x&=8-4t, \\
y&=3t, \\
z&=5-3t;\ t\in\mathbb{R}
\end{aligned} \)}{}{}}
\LANGpl{
\Question{
Dana jest prosta \( p \) określona punktami \( M=[4;3;2] \) i \( N=[8;0;5] \) (spójrz na rysunek). Wskaż równanie parametryczne prostej \( p' \), która jest symetryczna do prostej \( p \) w płaszczyźnie symetrii w układzie współrzędnych\( xz \).}
{\( \begin{aligned}
p'\colon x&=8+4t, \\
y&=3t, \\
z&=5+3t;\ t\in\mathbb{R}
\end{aligned} \)}{\( \begin{aligned}
p'\colon x&=8+4t, \\
y&=0, \\
z&=5+3t;\ t\in\mathbb{R}
\end{aligned} \)}{\( \begin{aligned}
p'\colon x&=8+4t, \\
y&=-3t, \\
z&=5+3t;\ t\in\mathbb{R}
\end{aligned} \)}{\( \begin{aligned}
p'\colon x&=8-4t, \\
y&=3t, \\
z&=5-3t;\ t\in\mathbb{R}
\end{aligned} \)}{}{}}
\LANGcs{
\Question{
Přímka \( p \) je zadána body \( M=[4;3;2] \) a \( N=[8;0;5] \) (viz obrázek). Určete parametrické rovnice přímky \( p' \), která je souměrná s přímkou \( p \) v rovinové souměrnosti podle souřadné roviny \( (xz) \).}
{\( \begin{aligned}
p'\colon x&=8+4t, \\
y&=3t, \\
z&=5+3t;\ t\in\mathbb{R}
\end{aligned} \)}{\( \begin{aligned}
p'\colon x&=8+4t, \\
y&=0, \\
z&=5+3t;\ t\in\mathbb{R}
\end{aligned} \)}{\( \begin{aligned}
p'\colon x&=8+4t, \\
y&=-3t, \\
z&=5+3t;\ t\in\mathbb{R}
\end{aligned} \)}{\( \begin{aligned}
p'\colon x&=8-4t, \\
y&=3t, \\
z&=5-3t;\ t\in\mathbb{R}
\end{aligned} \)}{}{}}
\LANGsk{
\Question{
Priamka \( p \) je zadaná bodmi \( M=[4;3;2] \) a \( N=[8;0;5] \) (viď obrázok). Určte parametrické rovnice priamky \( p' \), ktorá je súmerná s priamkou \( p \) v rovinovej súmernosti podľa súradnicovej roviny \( (xz) \).}
{\( \begin{aligned}
p'\colon x&=8+4t, \\
y&=3t, \\
z&=5+3t;\ t\in\mathbb{R}
\end{aligned} \)}{\( \begin{aligned}
p'\colon x&=8+4t, \\
y&=0, \\
z&=5+3t;\ t\in\mathbb{R}
\end{aligned} \)}{\( \begin{aligned}
p'\colon x&=8+4t, \\
y&=-3t, \\
z&=5+3t;\ t\in\mathbb{R}
\end{aligned} \)}{\( \begin{aligned}
p'\colon x&=8-4t, \\
y&=3t, \\
z&=5-3t;\ t\in\mathbb{R}
\end{aligned} \)}{}{}}
\LANGes{
\Question{
Dada la recta \( p \) que pasa por los puntos \( M=[4;3;2] \) y \( N=[8;0;5] \) (mira la imagen). Determina las ecuaciones paramétricas de la recta \( p' \) que es simétrica a la recta \( p \) respecto al \( xz \).}
{\( \begin{aligned}
p'\colon x&=8+4t, \\
y&=3t, \\
z&=5+3t;\ t\in\mathbb{R}
\end{aligned} \)}{\( \begin{aligned}
p'\colon x&=8+4t, \\
y&=0, \\
z&=5+3t;\ t\in\mathbb{R}
\end{aligned} \)}{\( \begin{aligned}
p'\colon x&=8+4t, \\
y&=-3t, \\
z&=5+3t;\ t\in\mathbb{R}
\end{aligned} \)}{\( \begin{aligned}
p'\colon x&=8-4t, \\
y&=3t, \\
z&=5-3t;\ t\in\mathbb{R}
\end{aligned} \)}{}{}}
\End

\Data\ID{1103212204}
\Updated{2020-07-15 03:07:21}
\Level{C}
\SubArea{Analytic geometry in a space}
\IMG{1103212204.pdf}
\LANGen{
\Question{
A cube \( ABCDEFGH \) with an edge length of \( 2 \) is placed in a coordinate system (see the picture). Let the point \( M \) be the centre of the edge \( EF \). Find the general form equation of the plane \( \rho \) passing through the points \( B \), \( D \), and \( G \) and calculate the distance of \( M \) from the plane \( \rho \).}
{\( \rho\colon x-y+z=0;\ |M\rho|=\sqrt3 \)}{\( \rho\colon x-y+z+2=0;\ |M\rho|=\sqrt3 \)}{\( \rho\colon x-y+z+2=0;\ |M\rho|=2\sqrt3 \)}{\( \rho\colon x-y+z=0;\ |M\rho|=2\sqrt3 \)}{}{}}
\LANGpl{
\Question{
Sześcian \( ABCDEFGH \) o krawędzi równej \( 2 \) znajduje się w układzie współrzędnych  (spójrz na rysunek). Punkt \( M \) to środek krawędzi \( EF \). Wskaż równanie ogólne płaszczyzny  \( \rho \) przechodzącej przez punkty  \( B \), \( D \), i \( G \) oraz oblicz odległość punktu  \( M \) do płaszczyzny \( \rho \).}
{\( \rho\colon x-y+z=0;\ |M\rho|=\sqrt3 \)}{\( \rho\colon x-y+z+2=0;\ |M\rho|=\sqrt3 \)}{\( \rho\colon x-y+z+2=0;\ |M\rho|=2\sqrt3 \)}{\( \rho\colon x-y+z=0;\ |M\rho|=2\sqrt3 \)}{}{}}
\LANGcs{
\Question{
Krychle \( ABCDEFGH \) s délkou hrany \( 2 \) je umístěna v souřadném systému (viz obrázek). Bod \( M \) je střed hrany \( EF \). Určete obecnou rovnici roviny \( \rho \) procházející body \( B \), \( D \) a \( G \) a vypočtěte vzdálenost bodu \( M \) od roviny \( \rho \).}
{\( \rho\colon x-y+z=0;\ |M\rho|=\sqrt3 \)}{\( \rho\colon x-y+z+2=0;\ |M\rho|=\sqrt3 \)}{\( \rho\colon x-y+z+2=0;\ |M\rho|=2\sqrt3 \)}{\( \rho\colon x-y+z=0;\ |M\rho|=2\sqrt3 \)}{}{}}
\LANGsk{
\Question{
Kocka \( ABCDEFGH \) s dĺžkou hrany \( 2 \) je umiestnená v súradnicovom systému (viď obrázok). Bod \( M \) je stred hrany \( EF \). Určte všeobecnú rovnicu roviny \( \rho \) prechádzajúcu bodmi \( B \), \( D \) a \( G \) a vypočítajte vzdialenosť bodu \( M \) od roviny \( \rho \).}
{\( \rho\colon x-y+z=0;\ |M\rho|=\sqrt3 \)}{\( \rho\colon x-y+z+2=0;\ |M\rho|=\sqrt3 \)}{\( \rho\colon x-y+z+2=0;\ |M\rho|=2\sqrt3 \)}{\( \rho\colon x-y+z=0;\ |M\rho|=2\sqrt3 \)}{}{}}
\LANGes{
\Question{
Sea \( ABCDEFGH \) el cubo cuya longitud de arista es \( 2 \) (mira la imagen) y sea el punto \( M \) el punto medio de la arista \( EF \). Determina la ecuación general del plano \( \rho \) que pasa por los puntos \( B \), \( D \), y \( G \) y calcula la distancia del punto \( M \) a dicho plano \( \rho \).}
{\( \rho\colon x-y+z=0;\ |M\rho|=\sqrt3 \)}{\( \rho\colon x-y+z+2=0;\ |M\rho|=\sqrt3 \)}{\( \rho\colon x-y+z+2=0;\ |M\rho|=2\sqrt3 \)}{\( \rho\colon x-y+z=0;\ |M\rho|=2\sqrt3 \)}{}{}}
\End

\Data\ID{1103212205}
\Updated{2020-07-15 03:07:21}
\Level{C}
\SubArea{Analytic geometry in a space}
\IMG{1103212205.pdf}
\LANGen{
\Question{
A cube \( ABCDEFGH \) with an edge length of \( 2 \) is placed in a coordinate system (see the picture). Find the distance between parallel planes \( \alpha \) and \( \beta \), where \( \alpha \) is passing through \( B \), \( D \) and \( G \) and \( \beta \) is passing through \( A \), \( F \) and \( H \).}
{\( |\alpha\beta|=\frac{2\sqrt3}3 \)}{\( |\alpha\beta|=\frac{4\sqrt3}3 \)}{\( |\alpha\beta|=\frac{3\sqrt3}2 \)}{\( |\alpha\beta|=\frac{3\sqrt3}4 \)}{}{}}
\LANGpl{
\Question{
Sześcian \( ABCDEFGH \) o krawędzi równej  \( 2 \) znajduje się w układzie współrzędnych  (spójrz na rysunek). Oblicz odległość pomiędzy równoległymi płaszczyznami \( \alpha \) i \( \beta \), gdzie \( \alpha \) przechodzi przez\( B \), \( D \) i \( G \) a \( \beta \) przechodzi przez \( A \), \( F \) i \( H \).}
{\( |\alpha\beta|=\frac{2\sqrt3}3 \)}{\( |\alpha\beta|=\frac{4\sqrt3}3 \)}{\( |\alpha\beta|=\frac{3\sqrt3}2 \)}{\( |\alpha\beta|=\frac{3\sqrt3}4 \)}{}{}}
\LANGcs{
\Question{
Krychle \( ABCDEFGH \) s délkou hrany \( 2 \) je umístěna v souřadném systému (viz obrázek). Vypočtěte vzdálenost rovnoběžných rovin \( \alpha \) a \( \beta \), kde \( \alpha \) je určena body \( B \), \( D \),  \( G \) a \( \beta \) je určena body \( A \), \( F \), \( H \).}
{\( |\alpha\beta|=\frac{2\sqrt3}3 \)}{\( |\alpha\beta|=\frac{4\sqrt3}3 \)}{\( |\alpha\beta|=\frac{3\sqrt3}2 \)}{\( |\alpha\beta|=\frac{3\sqrt3}4 \)}{}{}}
\LANGsk{
\Question{
Kocka \( ABCDEFGH \) s dĺžkou hrany \( 2 \) je umiestnená v súradnicovom systéme (viď obrázok). Vypočítajte vzdialenosť rovnobežných rovín \( \alpha \) a \( \beta \), kde \( \alpha \) je určená bodmi \( B \), \( D \),  \( G \) a \( \beta \) je určená bodmi \( A \), \( F \), \( H \).}
{\( |\alpha\beta|=\frac{2\sqrt3}3 \)}{\( |\alpha\beta|=\frac{4\sqrt3}3 \)}{\( |\alpha\beta|=\frac{3\sqrt3}2 \)}{\( |\alpha\beta|=\frac{3\sqrt3}4 \)}{}{}}
\LANGes{
\Question{
Sea \( ABCDEFGH \) el cubo cuya arista mide \( 2 \) (mira la imagen). Calcula la distancia de los planos paralelos \( \alpha \) y \( \beta \), en los cuales \( \alpha \) está determinado por los puntos \( B \), \( D \) y \( G \) y el plano \( \beta \) está determinado por los puntos \( A \), \( F \) y \( H \).}
{\( |\alpha\beta|=\frac{2\sqrt3}3 \)}{\( |\alpha\beta|=\frac{4\sqrt3}3 \)}{\( |\alpha\beta|=\frac{3\sqrt3}2 \)}{\( |\alpha\beta|=\frac{3\sqrt3}4 \)}{}{}}
\End

\Data\ID{1103212206}
\Updated{2020-07-15 03:07:21}
\Level{C}
\SubArea{Analytic geometry in a space}
\IMG{1103212206.pdf}
\LANGen{
\Question{
A cube \( ABCDEFGH \) with an edge length of \( 2 \) is placed in a coordinate system (see the picture). Let \( p \) be a line of intersection of planes \( \alpha \) and \( \beta \), where \( \alpha \) is passing through \( C \), \( F \) and \( H \) and \( \beta \) is passing through \( A \), \( F \) and \( H \). Find the parametric equations of the line \( p \) and calculate the angle \( \varphi \) between planes \( \alpha \) and \( \beta \) . Round \( \varphi \) to the nearest minute.}
{\( \begin{aligned}
p\colon x&=t, & \varphi&\,\dot{=}\,70^{\circ}32' \\
y&=t, & &\\
z&=2;\ t\in\mathbb{R},  & &
\end{aligned} \)}{\( \begin{aligned}
p\colon x&=2t, & \varphi&\,\dot{=}\,90^{\circ} \\
y&=2t, & & \\
z&=2+2t;\ t\in\mathbb{R},  & &
\end{aligned} \)}{\( \begin{aligned}
p\colon x&=t, &  \varphi&\,\dot{=}\,90^{\circ}\\
y&=t, & & \\
z&=2;\ t\in\mathbb{R}, & &
\end{aligned} \)}{\( \begin{aligned}
p\colon x&=2t, & \varphi&\,\dot{=}\,70^{\circ}32' \\
y&=2t, & & \\
z&=2t;\ t\in\mathbb{R}, & & 
\end{aligned} \)}{}{}}
\LANGpl{
\Question{
Sześcian \( ABCDEFGH \) o długości krawędzi równej  \( 2 \) jest umieszczony w układzie współrzędnych  (spójrz na rysunek). Prosta \( p \) to prosta przecięcia płaszczyzn  \( \alpha \) i \( \beta \), gdzie \( \alpha \) przechodzi przez \( C \), \( F \) i \( H \) a \( \beta \) przechodzi przez  \( A \), \( F \) i \( H \). Wskaż równanie parametryczne prostej \( p \) i oblicz miarę kąta \( \varphi \) pomiędzy płaszczyznami  \( \alpha \) i \( \beta \) . Zaokrągli \( \varphi \) do pełnych minut.}
{\( \begin{aligned}
p\colon x&=t, & \varphi&\,\dot{=}\,70^{\circ}32' \\
y&=t, & & \\
z&=2;\ t\in\mathbb{R},
\end{aligned} \)}{\( \begin{aligned}
p\colon x&=2t, & \varphi&\,\dot{=}\,90^{\circ} \\
y&=2t, & & \\
z&=2+2t;\ t\in\mathbb{R},
\end{aligned} \)}{\( \begin{aligned}
p\colon x&=t, &  \varphi&\,\dot{=}\,90^{\circ}\\
y&=t, & & \\
z&=2;\ t\in\mathbb{R}, & &
\end{aligned} \)}{\( \begin{aligned}
p\colon x&= 2t, & \varphi&\,\dot{=}\,70^{\circ}32' \\
y&=2t, & & \\
z&=2t;\ t\in\mathbb{R}, & &
\end{aligned} \)}{}{}}
\LANGcs{
\Question{
Krychle \( ABCDEFGH \) s délkou hrany \( 2 \) je umístěna v souřadném systému (viz obrázek). Přímka \( p \) je průsečnicí rovin \( \alpha \) a \( \beta \), kde \( \alpha \) je určena body \( C \), \( F \), \( H \) a \( \beta \) je určena body \( A \), \( F \), \( H \). Určete parametrické vyjádření přímky \( p \) a vypočtěte odchylku \( \varphi \) rovin \( \alpha \) a \( \beta \) . Odchylku \( \varphi \) zaokrouhlete na minuty.}
{\( \begin{aligned}
p\colon x&=t, & \varphi&\,\dot{=}\,70^{\circ}32'\\
y&=t, & &\\
z&=2;\ t\in\mathbb{R}, & &
\end{aligned} \)}{\( \begin{aligned}
p\colon x&=2t &  \varphi&\,\dot{=}\,90^{\circ} \\
y&=2t & & \\
z&=2+2t;\ t\in\mathbb{R}, & &
\end{aligned} \)}{\( \begin{aligned}
p\colon x&=t, &  \varphi&\,\dot{=}\,90^{\circ}\\
y&=t, & & \\
z&=2;\ t\in\mathbb{R}, & &
\end{aligned} \)}{\( \begin{aligned}
p\colon x&=2t, & \varphi&\,\dot{=}\,70^{\circ}32' \\
y&=2t & & \\
z&=2t;\ t\in\mathbb{R}, & &
\end{aligned} \)}{}{}}
\LANGsk{
\Question{
Kocka \( ABCDEFGH \) s dĺžkou hrany \( 2 \) je umiestnená v súradnicovom systéme (viď obrázok). Priamka \( p \) je priesečnica rovín \( \alpha \) a \( \beta \), kde \( \alpha \) je určená bodmi \( C \), \( F \), \( H \) a \( \beta \) je určená bodmi \( A \), \( F \), \( H \). Určte parametrické vyjadrenie priamky \( p \) a vypočítajte odchýlku \( \varphi \) rovín \( \alpha \) a \( \beta \) . Odchýlku \( \varphi \) zaokrúhlite na minúty.}
{\( \begin{aligned}
p\colon x&=t, & \varphi&\,\dot{=}\,70^{\circ}32' \\
y&=t, & & \\
z&=2;\ t\in\mathbb{R},
\end{aligned} \)}{\( \begin{aligned}
p\colon x&=2t, & \varphi&\,\dot{=}\,90^{\circ} \\
y&=2t, & & \\
z&=2+2t;\ t\in\mathbb{R},
\end{aligned} \)}{\( \begin{aligned}
p\colon x&=t, &  \varphi&\,\dot{=}\,90^{\circ}\\
y&=t, & & \\
z&=2;\ t\in\mathbb{R}, & &
\end{aligned} \)}{\( \begin{aligned}
p\colon x&= 2t, & \varphi&\,\dot{=}\,70^{\circ}32' \\
y&=2t, & & \\
z&=2t;\ t\in\mathbb{R}, & &
\end{aligned} \)}{}{}}
\LANGes{
\Question{
Sea \( ABCDEFGH \) el cubo cuya longitud de arista es \( 2 \) (mira la imagen). La línea \( p \) es línea de intersección de los planos \( \alpha \) y \( \beta \), donde el plano \( \alpha \) está determinado por los puntos \( C \), \( F \) y \( H \) y el plano \( \beta \) está determinado por los puntos \( A \), \( F \) y \( H \). Halla las ecuaciones paramétricas para la recta \( p \) y calcula el ángulo \( \varphi \) de los planos \( \alpha \) y \( \beta \) . Aproxima el ángulo \( \varphi \) a minutos.}
{\( \begin{aligned}
p\colon x&=t, & \varphi&\,\dot{=}\,70^{\circ}32' \\
y&=t, & &\\
z&=2;\ t\in\mathbb{R},  & &
\end{aligned} \)}{\( \begin{aligned}
p\colon x&=2t, & \varphi&\,\dot{=}\,90^{\circ} \\
y&=2t, & & \\
z&=2+2t;\ t\in\mathbb{R},  & &
\end{aligned} \)}{\( \begin{aligned}
p\colon x&=t, &  \varphi&\,\dot{=}\,90^{\circ}\\
y&=t, & & \\
z&=2;\ t\in\mathbb{R}, & &
\end{aligned} \)}{\( \begin{aligned}
p\colon x&=2t, & \varphi&\,\dot{=}\,70^{\circ}32' \\
y&=2t, & & \\
z&=2t;\ t\in\mathbb{R}, & & 
\end{aligned} \)}{}{}}
\End

\Data\ID{1103212901}
\Updated{2020-08-04 02:08:09}
\Level{C}
\SubArea{Analytic geometry in a space}
\IMG{1103212901.pdf}
\LANGen{
\Question{
A cube \( ABCDEFGH \) with an edge length of \( 2 \) units is placed in a coordinate system (see the picture). Find the distance of parallel lines \( p=KL \) and \( q=MN \), where points \( K \), \( L \), \( M \) and \( N \) are midpoints of edges \( CD \), \( BC \), \( EH \) and \( EF \) respectively.}
{\( |pq|=\sqrt6 \)}{\( |pq|=2\sqrt3 \)}{\( |pq|=3\sqrt2 \)}{\( |pq|=2\sqrt2 \)}{}{}}
\LANGpl{
\Question{
Sześcian  \( ABCDEFGH \) o długości krawędzi równej \( 2 \) jednostki znajduje się w układzie współrzędnych  (spójrz na rysunek). Oblicz odległość pomiędzy prostymi równoległymi  \( p=KL \) i \( q=MN \), gdzie punkty \( K \), \( L \), \( M \) i \( N \) to środki krawędzi  \( CD \), \( BC \), \( EH \) i \( EF \).}
{\( |pq|=\sqrt6 \)}{\( |pq|=2\sqrt3 \)}{\( |pq|=3\sqrt2 \)}{\( |pq|=2\sqrt2 \)}{}{}}
\LANGcs{
\Question{
Krychle \( ABCDEFGH \) s délkou hrany \( 2 \) je umístěna v souřadném systému (viz obrázek). Vypočtěte vzdálenost rovnoběžných přímek \( p=KL \) a \( q=MN \), kde body \( K \), \( L \), \( M \) a \( N \) jsou po řadě středy hran \( CD \), \( BC \), \( EH \) a \( EF \) .}
{\( |pq|=\sqrt6 \)}{\( |pq|=2\sqrt3 \)}{\( |pq|=3\sqrt2 \)}{\( |pq|=2\sqrt2 \)}{}{}}
\LANGsk{
\Question{
Kocka \( ABCDEFGH \) s dĺžkou hrany \( 2 \) je umiestnená v súradnicovom systéme (viď obrázok). Vypočítajte vzdialenosť rovnobežných priamok \( p=KL \) a \( q=MN \), kde body \( K \), \( L \), \( M \) a \( N \) sú po rade stredy hrán \( CD \), \( BC \), \( EH \) a \( EF \) .}
{\( |pq|=\sqrt6 \)}{\( |pq|=2\sqrt3 \)}{\( |pq|=3\sqrt2 \)}{\( |pq|=2\sqrt2 \)}{}{}}
\LANGes{
\Question{
Sea \( ABCDEFGH \) el cubo cuya arista mide \( 2 \) (mira la imagen).  Determina la distancia de las rectas paralelas \( p=KL \) y \( q=MN \) en las cuales los puntos \( K \), \( L \), \( M \) y \( N \) son puntos medios de aristas  \( CD \), \( BC \), \( EH \) y \( EF \) (siguiendo el orden).}
{\( |pq|=\sqrt6 \)}{\( |pq|=2\sqrt3 \)}{\( |pq|=3\sqrt2 \)}{\( |pq|=2\sqrt2 \)}{}{}}
\End

\Data\ID{1103212902}
\Updated{2020-07-15 03:07:21}
\Level{C}
\SubArea{Analytic geometry in a space}
\IMG{1103212902.pdf}
\LANGen{
\Question{
A cube \( ABCDEFGH \) with an edge length of \( 2 \) units is placed in a coordinate system (see the picture). Let \( S \) be the midpoint of the face \( ABFE \), and let  \( K \) and \( L \) be the midpoints of edges \( DH \) and \( CG \) consecutively. Find the standard equation of a plane \( \alpha \) passing through the points \( A \), \( B \) and \( L \), and calculate the distance of the point \( S \) from \( \alpha \).}
{\( \alpha\colon x+2z-2=0;\ |S\alpha|=\frac{2\sqrt5}{5} \)}{\( \alpha\colon x+2z-2=0;\ |S\alpha|=\frac{2\sqrt3}{3} \)}{\( \alpha\colon x+2y-2=0;\ |S\alpha|=\frac{2\sqrt5}{5} \)}{\( \alpha\colon x+2y-2=0;\ |S\alpha|=\frac{2\sqrt3}{3} \)}{}{}}
\LANGpl{
\Question{
Sześcian \( ABCDEFGH \) o krawędzi równej  \( 2 \) jednostki jest umieszczony w układzie współrzędnych (spójrz na rysunek). \( S \) to środek ściany  \( ABFE \),   \( K \) i  \( L \) to środki krawędzi  \( DH \) i \( CG \). Wskaż równanie płaszczyzny  \( \alpha \) przechodzącej przez punkty  \( A \), \( B \) i \( L \), oraz oblicz odległość punktu \( S \) od \( \alpha \).}
{\( \alpha\colon x+2z-2=0;\ |S\alpha|=\frac{2\sqrt5}{5} \)}{\( \alpha\colon x+2z-2=0;\ |S\alpha|=\frac{2\sqrt3}{3} \)}{\( \alpha\colon x+2y-2=0;\ |S\alpha|=\frac{2\sqrt5}{5} \)}{\( \alpha\colon x+2y-2=0;\ |S\alpha|=\frac{2\sqrt3}{3} \)}{}{}}
\LANGcs{
\Question{
Krychle \( ABCDEFGH \) s délkou hrany \( 2 \) je umístěna v souřadném systému (viz obrázek). Bod \( S \) je středem stěny \( ABFE \) a body  \( K \) a \( L \) jsou po řadě středy hran \( DH \) a \( CG \) . Určete obecnou rovnici roviny \( \alpha \) procházející body \( A \), \( B \) a \( L \) a vypočtěte vzdálenost bodu \( S \) od roviny \( \alpha \).}
{\( \alpha\colon x+2z-2=0;\ |S\alpha|=\frac{2\sqrt5}{5} \)}{\( \alpha\colon x+2z-2=0;\ |S\alpha|=\frac{2\sqrt3}{3} \)}{\( \alpha\colon x+2y-2=0;\ |S\alpha|=\frac{2\sqrt5}{5} \)}{\( \alpha\colon x+2y-2=0;\ |S\alpha|=\frac{2\sqrt3}{3} \)}{}{}}
\LANGsk{
\Question{
Kocka \( ABCDEFGH \) s dĺžkou hrany \( 2 \) je umiestnená v súradnicovom systéme (viď obrázok). Bod \( S \) je stredom steny \( ABFE \) a body  \( K \) a \( L \) sú po rade stredy hrán \( DH \) a \( CG \) . Určte všeobecnú rovnicu roviny \( \alpha \) prechádzajúcimi bodmi \( A \), \( B \) a \( L \) a vypočítajte vzdialenosť bodu \( S \) od roviny \( \alpha \).}
{\( \alpha\colon x+2z-2=0;\ |S\alpha|=\frac{2\sqrt5}{5} \)}{\( \alpha\colon x+2z-2=0;\ |S\alpha|=\frac{2\sqrt3}{3} \)}{\( \alpha\colon x+2y-2=0;\ |S\alpha|=\frac{2\sqrt5}{5} \)}{\( \alpha\colon x+2y-2=0;\ |S\alpha|=\frac{2\sqrt3}{3} \)}{}{}}
\LANGes{
\Question{
Sea \( ABCDEFGH \) el cubo cuya arista mide \( 2 \) (mira la imagen). El punto \( S \) es el centro de la cara \( ABFE \) y los puntos \( K \) y \( L \) son los puntos medios de aristas \( DH \) y \( CG \) consecutivamente. Determina la ecuación general del plano \( \alpha \) que pase por los puntos \( A \), \( B \) y \( L \), y calcula la distancia del punto \( S \) al plano \( \alpha \).}
{\( \alpha\colon x+2z-2=0;\ |S\alpha|=\frac{2\sqrt5}{5} \)}{\( \alpha\colon x+2z-2=0;\ |S\alpha|=\frac{2\sqrt3}{3} \)}{\( \alpha\colon x+2y-2=0;\ |S\alpha|=\frac{2\sqrt5}{5} \)}{\( \alpha\colon x+2y-2=0;\ |S\alpha|=\frac{2\sqrt3}{3} \)}{}{}}
\End

\Data\ID{1103212903}
\Updated{2020-07-15 03:07:21}
\Level{C}
\SubArea{Analytic geometry in a space}
\IMG{1103212903.pdf}
\LANGen{
\Question{
A cube \( ABCDEFGH \) with an edge length of \( 2 \) units is placed in a coordinate system (see the picture). Find an angle \( \varphi \) between the plane \( \alpha \) passing through the points \( E \), \( D \) and \( C \) and the straight line \( AF \).

Hint: An angle between a line and a plane is an angle between the line and its orthogonal projection into this plane.}
{\( \varphi = 30^{\circ} \)}{\( \varphi = 15^{\circ} \)}{\( \varphi = 45^{\circ} \)}{\( \varphi = 60^{\circ} \)}{}{}}
\LANGpl{
\Question{
Sześcian \( ABCDEFGH \) o długości krawędzi równej \( 2 \) jednostki został umieszczony w układzie współrzędnych (spójrz na rysunek). Wyznacz kąt \( \varphi \) pomiędzy płaszczyzną \( \alpha \) przechodzącą przez punkty \( E \), \( D \) i \( C \) a prostą  \( AF \).

Wskazówka: Kąt pomiędzy prostą a płaszczyzną to kąt pomiędzy prostą a jej rzutem prostopadłym na płaszczyznę.}
{\( \varphi = 30^{\circ} \)}{\( \varphi = 15^{\circ} \)}{\( \varphi = 45^{\circ} \)}{\( \varphi = 60^{\circ} \)}{}{}}
\LANGcs{
\Question{
Krychle \( ABCDEFGH \) s délkou hrany \( 2 \) je umístěna v souřadném systému (viz obrázek). Vypočtěte odchylku \( \varphi \) přímky \( AF \) od roviny \( \alpha \) procházející body \( E \), \( D \) a \( C \).

Nápověda: Odchylka přímky od roviny je odchylka přímky od jejího kolmého průmětu do této roviny.}
{\( \varphi = 30^{\circ} \)}{\( \varphi = 15^{\circ} \)}{\( \varphi = 45^{\circ} \)}{\( \varphi = 60^{\circ} \)}{}{}}
\LANGsk{
\Question{
Kocka \( ABCDEFGH \) s dĺžkou hrany \( 2 \) je umiestnená v súradnicovom systéme (viď obrázok). Vypočítajte odchýlku \( \varphi \) priamky \( AF \) od roviny \( \alpha \) prechádzajúcej bodmi \( E \), \( D \) a \( C \).

Nápoveda: Odchýlka priamky od roviny je odchýlka priamky od jej kolmého priemetu do tejto roviny.}
{\( \varphi = 30^{\circ} \)}{\( \varphi = 15^{\circ} \)}{\( \varphi = 45^{\circ} \)}{\( \varphi = 60^{\circ} \)}{}{}}
\LANGes{
\Question{
Sea \( ABCDEFGH \) el cubo cuya longitud de la arista es \( 2 \) (mira la imagen). Determina el ángulo \( \varphi \) formado entre la recta \( AF \) y el plano \( \alpha \) que pasa por los puntos \( E \), \( D \) y \( C \).

Pista: El ángulo formado entre una recta y un plano es igual al ángulo formado entre la proyección perpendicular de la recta en el plano.}
{\( \varphi = 30^{\circ} \)}{\( \varphi = 15^{\circ} \)}{\( \varphi = 45^{\circ} \)}{\( \varphi = 60^{\circ} \)}{}{}}
\End

\Data\ID{1103212904}
\Updated{2020-07-15 03:07:21}
\Level{C}
\SubArea{Analytic geometry in a space}
\IMG{1103212904.pdf}
\LANGen{
\Question{
A rectangle-based right pyramid \( ABCDV \) with a bottom edge length of \( 6 \) units and the perpendicular height of \( 6 \) units is placed in a coordinate system (see the picture). Let \( S \) be the midpoint of the edge \( AD \). Find the standard equation of the plane \( \alpha \) passing through the points \( B \), \( V \) and \( C \), and calculate the distance of the point \( S \) from \( \alpha \).}
{\( \alpha\colon 2y+z-12=0;\ d=|S\alpha|=\frac{12\sqrt5}{5} \)}{\( \alpha\colon 2x+z-12=0;\ d=|S\alpha|=\frac{12\sqrt5}{5} \)}{\( \alpha\colon 2y+z-12=0;\ d=|S\alpha|=\frac{6\sqrt5}{5} \)}{\( \alpha\colon 2x+z-12=0;\ d=|S\alpha|=\frac{6\sqrt5}{5} \)}{}{}}
\LANGpl{
\Question{
Dany jest ostrosłup  \( ABCDV \), którego podstawą jest prostokąt, długość podstawy jest równa   \( 6 \) jednostek, jego wysokość jest równa \( 6 \) jednostek. Ostrosłup został umieszczony w układzie współrzędnych  (spójrz na rysunek). Punkt \( S \) to środek krawędzi \( AD \). Wyznacz równanie płaszczyzny  \( \alpha \) przechodzącej przez punkty  \( B \), \( V \) i \( C \) oraz oblicz odległość punktu \( S \) od \( \alpha \).}
{\( \alpha\colon 2y+z-12=0;\ d=|S\alpha|=\frac{12\sqrt5}{5} \)}{\( \alpha\colon 2x+z-12=0;\ d=|S\alpha|=\frac{12\sqrt5}{5} \)}{\( \alpha\colon 2y+z-12=0;\ d=|S\alpha|=\frac{6\sqrt5}{5} \)}{\( \alpha\colon 2x+z-12=0;\ d=|S\alpha|=\frac{6\sqrt5}{5} \)}{}{}}
\LANGcs{
\Question{
Pravidelný čtyřboký jehlan \( ABCDV \) s délkou podstavné hrany \( 6 \) a tělesovou výškou \( 6 \) je umístěn v souřadném systému (viz obrázek). Bod \( S \) je středem hrany \( AD \). Určete obecnou rovnici roviny \( \alpha \) procházející body \( B \), \( V \) a \( C \) a vypočtěte vzdálenost bodu \( S \) od této roviny.}
{\( \alpha\colon 2y+z-12=0;\ d=|S\alpha|=\frac{12\sqrt5}{5} \)}{\( \alpha\colon 2x+z-12=0;\ d=|S\alpha|=\frac{12\sqrt5}{5} \)}{\( \alpha\colon 2y+z-12=0;\ d=|S\alpha|=\frac{6\sqrt5}{5} \)}{\( \alpha\colon 2x+z-12=0;\ d=|S\alpha|=\frac{6\sqrt5}{5} \)}{}{}}
\LANGsk{
\Question{
Pravidelný štvorboký ihlan \( ABCDV \) s dĺžkou hrany podstavy \( 6 \) a telesovou výškou \( 6 \) je umiestnený v súradnicovom systéme (viď obrázok). Bod \( S \) je stredom hrany \( AD \). Určte všeobecnú rovnicu roviny \( \alpha \) prechádzajúcou bodmi \( B \), \( V \) a \( C \) a vypočítajte vzdialenosť bodu \( S \) od tejto roviny.}
{\( \alpha\colon 2y+z-12=0;\ d=|S\alpha|=\frac{12\sqrt5}{5} \)}{\( \alpha\colon 2x+z-12=0;\ d=|S\alpha|=\frac{12\sqrt5}{5} \)}{\( \alpha\colon 2y+z-12=0;\ d=|S\alpha|=\frac{6\sqrt5}{5} \)}{\( \alpha\colon 2x+z-12=0;\ d=|S\alpha|=\frac{6\sqrt5}{5} \)}{}{}}
\LANGes{
\Question{
Sea \( ABCDV \) la pirámide cuadrangular regular cuya arista de base es \( 6 \) y la altura de todo el sólido es \( 6 \) (mira la imagen). El punto \( S \) es el punto medio de la arista \( AD \). Determina la ecuación general del plano \( \alpha \)  que pasa por los puntos \( B \), \( V \) y \( C \), y calcula la distancia del punto \( S \) al plano \( \alpha \).}
{\( \alpha\colon 2y+z-12=0;\ d=|S\alpha|=\frac{12\sqrt5}{5} \)}{\( \alpha\colon 2x+z-12=0;\ d=|S\alpha|=\frac{12\sqrt5}{5} \)}{\( \alpha\colon 2y+z-12=0;\ d=|S\alpha|=\frac{6\sqrt5}{5} \)}{\( \alpha\colon 2x+z-12=0;\ d=|S\alpha|=\frac{6\sqrt5}{5} \)}{}{}}
\End

\Data\ID{1103212905}
\Updated{2020-07-15 03:07:21}
\Level{C}
\SubArea{Analytic geometry in a space}
\IMG{1103212905.pdf}
\LANGen{
\Question{
A rectangle-based right pyramid \( ABCDV \) with its bottom edge length of \( 6 \) units and the perpendicular height of \( 6 \) units is placed in a coordinate system (see the picture). Find the parametric equations of an intersection line \( p \) of planes \( \alpha \) and \( \beta \), where \( \alpha \) passes through the points \( B \), \( C \) and \( V \), and \( \beta \) passes through the points \( A \), \( D \) and \( V \). What is the measure of an angle \( \varphi \) between the planes \( \alpha \) and \( \beta \). Round \( \varphi \) to the nearest minute.}
{\(\begin{aligned}
p\colon x&=3+t, & \varphi\,\dot{=}\,53^{\circ}8'\\
y&=3, &\\
z&=6;\ t\in\mathbb{R} & 
\end{aligned}\)}{\(\begin{aligned}
p\colon x&=3+t, & \varphi\,\dot{=}\,63^{\circ}8'\\
y&=3, &\\
z&=0;\ t\in\mathbb{R} & 
\end{aligned}\)}{\(\begin{aligned}
p\colon x&=3+t, & \varphi\,\dot{=}\,53^{\circ}8'\\
y&=3+t, &\\
z&=6+2t;\ t\in\mathbb{R} & 
\end{aligned}\)}{\(\begin{aligned}
p\colon x&=3+t, & \varphi\,\dot{=}\,63^{\circ}8'\\
y&=3, &\\
z&=6;\ t\in\mathbb{R} & 
\end{aligned}\)}{}{}}
\LANGpl{
\Question{
Dany jest ostrosłup \( ABCDV \), którego podstawą jest prostokąt, długość krawędzi jest równa  \( 6 \) jednostek a jego wysokość \( 6 \) jednostek. Ostrosłup jest umieszczony w układzie współrzędnych  (spójrz na rysunek). Wyznacz równanie parametryczne prostej przecięcia \( p \) płaszczyzny  \( \alpha \) i \( \beta \), gdzie \( \alpha \) przechodzi przez punkty \( B \), \( C \) i \( V \) oraz \( \beta \) przechodzi przez punkty \( A \), \( D \) i \( V \). Jaka jest miara kąta \( \varphi \) pomiędzy płaszczyznami \( \alpha \) i \( \beta \). Zaokrągli \( \varphi \) do pełnych minut.}
{\(\begin{aligned}
p\colon x&=3+t, & \varphi\,\dot{=}\,53^{\circ}8'\\
y&=3, &\\
z&=6;\ t\in\mathbb{R}, & 
\end{aligned}\)}{\(\begin{aligned}
p\colon x&=3+t, & \varphi\,\dot{=}\,63^{\circ}8'\\
y&=3, &\\
z&=0;\ t\in\mathbb{R}, & 
\end{aligned}\)}{\(\begin{aligned}
p\colon x&=3+t, & \varphi\,\dot{=}\,53^{\circ}8'\\
y&=3+t, &\\
z&=6+2t;\ t\in\mathbb{R}, & 
\end{aligned}\)}{\(\begin{aligned}
p\colon x&=3+t, & \varphi\,\dot{=}\,63^{\circ}8'\\
y&=3, &\\
z&=6;\ t\in\mathbb{R}, & 
\end{aligned}\)}{}{}}
\LANGcs{
\Question{
Pravidelný čtyřboký jehlan \( ABCDV \) s délkou podstavné hrany \( 6 \) a tělesovou výškou \( 6 \) je umístěn v souřadném systému (viz obrázek). Určete parametrické vyjádření průsečnice \( p \) rovin \( \alpha \) a \( \beta \), kde \( \alpha \) je rovina procházející body \( B \), \( C \) a \( V \) a \( \beta \) je rovina procházející body \( A \), \( D \) a \( V \). Dále vypočtěte velikost odchylky \( \varphi \) mezi rovinami \( \alpha \) a \( \beta \). Odchylku \( \varphi \) zaokrouhlete na minuty.}
{\(\begin{aligned}
p\colon x&=3+t, & \varphi\,\dot{=}\,53^{\circ}8'\\
y&=3, &\\
z&=6;\ t\in\mathbb{R}, & 
\end{aligned}\)}{\(\begin{aligned}
p\colon x&=3+t, & \varphi\,\dot{=}\,63^{\circ}8'\\
y&=3, &\\
z&=0;\ t\in\mathbb{R}, & 
\end{aligned}\)}{\(\begin{aligned}
p\colon x&=3+t, & \varphi\,\dot{=}\,53^{\circ}8'\\
y&=3+t, &\\
z&=6+2t;\ t\in\mathbb{R}, & 
\end{aligned}\)}{\(\begin{aligned}
p\colon x&=3+t, & \varphi\,\dot{=}\,63^{\circ}8'\\
y&=3, &\\
z&=6;\ t\in\mathbb{R}, & 
\end{aligned}\)}{}{}}
\LANGsk{
\Question{
Pravidelný štvorboký ihlan \( ABCDV \) s dĺžkou hrany podstavy \( 6 \) a telesovou výškou \( 6 \) je umiestený v súradnicovom systéme (viď obrázok). Určte parametrické vyjadrenie priesečnice \( p \) rovín \( \alpha \) a \( \beta \), kde \( \alpha \) je rovina prechádzajúce bodmi \( B \), \( C \) a \( V \) a \( \beta \) je rovina prechádzajúce bodmi \( A \), \( D \) a \( V \). Ďalej vypočítajte veľkosť odchýlky \( \varphi \) medzi rovinami \( \alpha \) a \( \beta \). Odchýlku \( \varphi \) zaokrúhlite na minúty.}
{\(\begin{aligned}
p\colon x&=3+t, & \varphi\,\dot{=}\,53^{\circ}8'\\
y&=3, &\\
z&=6;\ t\in\mathbb{R}, & 
\end{aligned}\)}{\(\begin{aligned}
p\colon x&=3+t, & \varphi\,\dot{=}\,63^{\circ}8'\\
y&=3, &\\
z&=0;\ t\in\mathbb{R}, & 
\end{aligned}\)}{\(\begin{aligned}
p\colon x&=3+t, & \varphi\,\dot{=}\,53^{\circ}8'\\
y&=3+t, &\\
z&=6+2t;\ t\in\mathbb{R}, & 
\end{aligned}\)}{\(\begin{aligned}
p\colon x&=3+t, & \varphi\,\dot{=}\,63^{\circ}8'\\
y&=3, &\\
z&=6;\ t\in\mathbb{R}, & 
\end{aligned}\)}{}{}}
\LANGes{
\Question{
Sea \( ABCDV \) la pirámide cuadrangular regular cuya arista de base es \( 6 \) y la altura de todo el sólido es \( 6 \) (mira la imagen). Determina las ecuaciones paramétricas de la línea de intersección \( p \) y los planos \( \alpha \) y \( \beta \), donde \( \alpha \) es el plano que pasa por los puntos \( B \), \( C \) y \( V \), y \( \beta \) es el plano que pasa por los puntos \( A \), \( D \) y \( V \). Calcula también el ángulo \( \varphi \) formado entre los planos \( \alpha \) y \( \beta \). Aproxima el ángulo \( \varphi \) a minutos.}
{\(\begin{aligned}
p\colon x&=3+t, & \varphi\,\dot{=}\,53^{\circ}8'\\
y&=3, &\\
z&=6;\ t\in\mathbb{R} & 
\end{aligned}\)}{\(\begin{aligned}
p\colon x&=3+t, & \varphi\,\dot{=}\,63^{\circ}8'\\
y&=3, &\\
z&=0;\ t\in\mathbb{R} & 
\end{aligned}\)}{\(\begin{aligned}
p\colon x&=3+t, & \varphi\,\dot{=}\,53^{\circ}8'\\
y&=3+t, &\\
z&=6+2t;\ t\in\mathbb{R} & 
\end{aligned}\)}{\(\begin{aligned}
p\colon x&=3+t, & \varphi\,\dot{=}\,63^{\circ}8'\\
y&=3, &\\
z&=6;\ t\in\mathbb{R} & 
\end{aligned}\)}{}{}}
\End

\Data\ID{1103233601}
\Updated{2020-07-15 04:07:40}
\Level{C}
\SubArea{Analytic geometry in a space}
\IMG{1103233601.pdf}
\LANGen{
\Question{
Let $ABCDEFGH$ be a cube with an edge length of $1$ placed in the rectangular coordinate system.  In the cube a regular tetrahedron $ACHF$ is highlighted (see the picture). Find its perpendicular height. \[ \] Hint: Find e.g. the distance between the point $F$ and the plane $ACH$.}
{$\frac{2\sqrt3}3$}{$\frac{\sqrt3}3$}{$\frac{2\sqrt6}3$}{$\frac23$}{}{}}
\LANGpl{
\Question{
Dany jest sześcian $ABCDEFGH$ długość jego krawędzi jest równa $1$, sześcian umieszczono w układzie współrzędnych.  W sześcianie podświetlono czworościan foremny $ACHF$  (spójrz na rysunek). Oblicz jego prostopadłą wysokość. \[ \]  Wskazówka: Oblicz odległość pomiędzy punktem $F$ a płaszczyzną  $ACH$.}
{$\frac{2\sqrt3}3$}{$\frac{\sqrt3}3$}{$\frac{2\sqrt6}3$}{$\frac23$}{}{}}
\LANGcs{
\Question{
V krychli $ABCDEFGH$ s hranou délky $1$, která je umístěna v souřadném systému, je vyznačen pravidelný čtyřstěn $ACHF$ (viz obrázek). Určete velikost jeho tělesové výšky. 
\[ \] Nápověda: Vypočtěte např. vzdálenost bodu $F$ od roviny $ACH$.}
{$\frac{2\sqrt3}3$}{$\frac{\sqrt3}3$}{$\frac{2\sqrt6}3$}{$\frac23$}{}{}}
\LANGsk{
\Question{
Nech $ABCDEFGH$ je kocka s dĺžkou hrany $1$ umiestnená v pravouhlon súradnicovo systéme. V kocke je zvýraznený pravidelný tetraéder $ACHF$ (viď obrázok). Vypočítajte jeho výšku.\[ \]   Nápoveda: Nájdite napr. vzdialenosť medzi bodom $F$ a rovinou $ACH$.}
{$\frac{2\sqrt3}3$}{$\frac{\sqrt3}3$}{$\frac{2\sqrt6}3$}{$\frac23$}{}{}}
\LANGes{
\Question{
En el cubo $ABCDEFGH$ cuya arista mide $1$ hay un tetraedro regular $ACHF$ (mira la imagen). Determina la altura del sólido.
\[ \] Pista: Calcula por ejemplo la distancia entre el punto $F$ y el plano $ACH$.}
{$\frac{2\sqrt3}3$}{$\frac{\sqrt3}3$}{$\frac{2\sqrt6}3$}{$\frac23$}{}{}}
\End

\Data\ID{1103233602}
\Updated{2020-07-15 04:07:40}
\Level{C}
\SubArea{Analytic geometry in a space}
\IMG{1103233602.pdf}
\LANGen{
\Question{
Let $ABCDEFGH$ be a cube with an edge length of $1$ placed in the rectangular coordinate system.  In the cube a regular tetrahedron $ACHF$ is highlighted (see the picture). Find the distance between the opposite edges of this tetrahedron.\[ \]   Hint: A tetrahedron’s  opposite edges lie on skew lines. Their distance is the same as the distance of the midpoint of one edge from  the opposite edge.}
{$1$}{$\sqrt3$}{$\frac{\sqrt3}2$}{$\frac{\sqrt5}2$}{}{}}
\LANGpl{
\Question{
Dany jest sześcian $ABCDEFGH$ długość jego krawędzi jest równa $1$, sześcian umieszczono w układzie współrzędnych. W sześcianie podświetlono czworościan foremny $ACHF$ (spójrz na rysunek).  Oblicz odległość pomiędzy przeciwległymi krawędziami czworościanu. \[ \]  Wskazówka: Przeciwległe krawędzie czworościanu leżą na prostych skośnych. Ich odległość jest taka sama jak odległość punktu środkowego jednej krawędzi do przeciwległej krawędzi.}
{$1$}{$\sqrt3$}{$\frac{\sqrt3}2$}{$\frac{\sqrt5}2$}{}{}}
\LANGcs{
\Question{
V krychli $ABCDEFGH$ s hranou délky $1$, která je umístěna v souřadném systému, je vyznačen pravidelný čtyřstěn $ACHF$ (viz obrázek). Vypočtěte vzdálenost jeho protilehlých hran. 
\[ \]  Nápověda: Protilehlé hrany čtyřstěnu leží na mimoběžných přímkách. Jejich vzdálenost je rovna vzdálenosti středu jedné hrany od hrany k ní protilehlé.}
{$1$}{$\sqrt3$}{$\frac{\sqrt3}2$}{$\frac{\sqrt5}2$}{}{}}
\LANGsk{
\Question{
V kocke $ABCDEFGH$ s hranou dĺžky $1$, ktorá je umiestnená v súradnicovom systéme, je vyznačený pravidelný štvorsten $ACHF$ (viď obrázok). Vypočítajte vzdialenosť jeho protiľahlých hrán. 
\[ \]  Nápoveda: Protiľahlé hrany štvorstena ležia na mimobežných priamkach. Ich vzdialenosť je rovná vzdialenosti stredu jednej hrany od hrany k nej protiľahlej.}
{$1$}{$\sqrt3$}{$\frac{\sqrt3}2$}{$\frac{\sqrt5}2$}{}{}}
\LANGes{
\Question{
En el cubo $ABCDEFGH$ cuya arista mide $1$ está marcado un tetraedro regular $ACHF$ (mira la imagen). Determina la distancia de sus aristas opuestas. \[ \] Pista: Las aristas opuestas de un tetraedro están en las rectas sesgadas. Su distancia es igual a la distancia de puntos medios de una arista y la arista opuesta a dicha arista.}
{$1$}{$\sqrt3$}{$\frac{\sqrt3}2$}{$\frac{\sqrt5}2$}{}{}}
\End

\Data\ID{1103233603}
\Updated{2020-07-15 04:07:40}
\Level{C}
\SubArea{Analytic geometry in a space}
\IMG{1103233603.pdf}
\LANGen{
\Question{
Let $ABCDEFGH$ be a cube with an edge length of $1$ placed in the rectangular coordinate system.  In the cube a regular tetrahedron $ACHF$ is highlighted (see the picture).
Find the angle between its faces and round the number to the nearest minute.}
{$70^{\circ}32'$}{$54^{\circ}44'$}{$45^{\circ}$}{$51^{\circ}4'$}{}{}}
\LANGpl{
\Question{
Dany jest sześcian $ABCDEFGH$, którego krawędź jest równa $1$, sześcian umieszczono  w układzie współrzędnych.  W sześcianie podświetlono foremny czworościan $ACHF$  (spójrz na rysunek).
 Wskaż miarę kąta pomiędzy ścianami oraz zaokrągli wynik do pełnych minut.}
{$70^{\circ}32'$}{$54^{\circ}44'$}{$45^{\circ}$}{$51^{\circ}4'$}{}{}}
\LANGcs{
\Question{
V krychli $ABCDEFGH$ s hranou délky $1$, která je umístěna v souřadném systému, je vyznačen pravidelný čtyřstěn $ACHF$ (viz obrázek). Vypočtěte odchylku jeho stěn a zaokrouhlete ji na minuty.}
{$70^{\circ}32'$}{$54^{\circ}44'$}{$45^{\circ}$}{$51^{\circ}4'$}{}{}}
\LANGsk{
\Question{
V kocke $ABCDEFGH$ s hranou dĺžky $1$, ktorá je umiestnená v súradnicovom systéme, je vyznačený pravidelný štvorsten $ACHF$ (viď obrázok). Vypočítajte odchýlku jeho stien a zaokrúhlite ju na minúty.}
{$70^{\circ}32'$}{$54^{\circ}44'$}{$45^{\circ}$}{$51^{\circ}4'$}{}{}}
\LANGes{
\Question{
En el cubo $ABCDEFGH$ cuya arista mide $1$ hay un tetraedro regular marcado $ACHF$ (mira la imagen). Determina el ángulo de sus caras y aproxímalo a minutos.}
{$70^{\circ}32'$}{$54^{\circ}44'$}{$45^{\circ}$}{$51^{\circ}4'$}{}{}}
\End

\Data\ID{1103233604}
\Updated{2020-07-15 04:07:40}
\Level{C}
\SubArea{Analytic geometry in a space}
\IMG{1103233604.pdf}
\LANGen{
\Question{
Reflect the point $A=[1;10;-8]$ over the straight line $p$:
\begin{align*}
p\colon x&= 1-2t, \\
      y&= 3+t, \\
      z&= -1+3t;\  t\in\mathbb{R}.
\end{align*}
Hint: See the picture.}
{$A'=[5;-6;0]$}{$A'=[3;2;-4]$}{$A'=[-1;11;-5]$}{$A'=[-2;10;-24]$}{}{}}
\LANGpl{
\Question{
Określ obraz punktu $A=[1;10;-8]$ w symetrii osiowej względem prostej $p$:
\begin{align*}
p\colon x&= 1-2t, \\
      y&= 3+t, \\
      z&= -1+3t;\  t\in\mathbb{R}.
\end{align*}
Wskazówka: Spójrz na rysunek.}
{$A'=[5;-6;0]$}{$A'=[3;2;-4]$}{$A'=[-1;11;-5]$}{$A'=[-2;10;-24]$}{}{}}
\LANGcs{
\Question{
Určete obraz bodu $A=[1;10;-8]$ v osové souměrnosti podle přímky $p$:
\begin{align*}
p\colon x&= 1-2t, \\
      y&= 3+t, \\
      z&= -1+3t;\  t\in\mathbb{R}.
\end{align*}
Nápověda: viz obrázek}
{$A'=[5;-6;0]$}{$A'=[3;2;-4]$}{$A'=[-1;11;-5]$}{$A'=[-2;10;-24]$}{}{}}
\LANGsk{
\Question{
Určte obraz bodu $A=[1;10;-8]$ v osovej súmernosti podľa priamky $p$:
\begin{align*}
p\colon x&= 1-2t, \\
      y&= 3+t, \\
      z&= -1+3t;\  t\in\mathbb{R}.
\end{align*}
Nápoveda: viď obrázok}
{$A'=[5;-6;0]$}{$A'=[3;2;-4]$}{$A'=[-1;11;-5]$}{$A'=[-2;10;-24]$}{}{}}
\LANGes{
\Question{
Refleja el punto $A=[1;10;-8]$ en la simetría de eje según la recta $p$:
\begin{align*}
p\colon x&= 1-2t, \\
      y&= 3+t, \\
      z&= -1+3t;\  t\in\mathbb{R}.
\end{align*}
Pista: mira la imagen.}
{$A'=[5;-6;0]$}{$A'=[3;2;-4]$}{$A'=[-1;11;-5]$}{$A'=[-2;10;-24]$}{}{}}
\End

\Data\ID{9000106303}
\Updated{2020-07-15 03:07:37}
\Level{C}
\SubArea{Analytic geometry in a space}
\IMG{001063_03-figure0.pdf}
\LANGen{
\Question{
In the plane symmetry (reflection) defined by the plane
\(\alpha \),
\[ 
 \alpha \colon 2x + y - z - 5 = 0,
\]
find the image \(A'\) 
of the point \(A = [0;0;1]\).}
{\(A' = [4;2;-1]\)}{\(A' = [6;3;-2]\)}{\(A' = [4;2;1]\)}{\(A' = [0;0;1]\)}{}{}}
\LANGpl{
\Question{
Na płaszczyźnie symetrii określonej płaszczyzną 
\(\alpha \),
\[ 
 \alpha \colon 2x + y - z - 5 = 0,
\]
wyznacz  współrzędne punktu \(A'\) 
będącego obrazem \(A = [0;0;1]\).}
{\(A' = [4;2;-1]\)}{\(A' = [6;3;-2]\)}{\(A' = [4;2;1]\)}{\(A' = [0;0;1]\)}{}{}}
\LANGcs{
\Question{
Rovina \(\alpha \) je zadaná
obecnou rovnicí: \(2x + y - z - 5 = 0\). Určete
souřadnice bodu \(A'\), který je
obrazem bodu \(A = [0;0;1]\) v rovinové
souměrnosti podle roviny \(\alpha \).}
{\(A' = [4;2;-1]\)}{\(A' = [6;3;-2]\)}{\(A' = [4;2;1]\)}{\(A' = [0;0;1]\)}{}{}}
\LANGsk{
\Question{
Rovina \(\alpha \) je daná
všeobecnou rovnicou: \(2x + y - z - 5 = 0\). Určte
súradnice bodu \(A'\), ktorý je
obrazom bodu \(A = [0;0;1]\) v rovinovej
súmernosti podľa roviny \(\alpha \).}
{\(A' = [4;2;-1]\)}{\(A' = [6;3;-2]\)}{\(A' = [4;2;1]\)}{\(A' = [0;0;1]\)}{}{}}
\LANGes{
\Question{
En el plano de simetría definido por el plano \(\alpha \),
\[ 
 \alpha \colon 2x + y - z - 5 = 0,
\]
halla la imagen \(A'\) 
del punto \(A = [0;0;1]\).}
{\(A' = [4;2;-1]\)}{\(A' = [6;3;-2]\)}{\(A' = [4;2;1]\)}{\(A' = [0;0;1]\)}{}{}}
\End

\Data\ID{9000106307}
\Updated{2020-07-15 03:07:37}
\Level{C}
\SubArea{Analytic geometry in a space}
\IMG{001063_07-figure0.pdf}
\LANGen{
\Question{
Given points \(A = [0;0;1]\),
\(B = [2;0;-1]\) and
\(S = [2;1;0]\),
find the parametric equations of the image of the line
\(AB\) in a point reflection
about the point \(S\).}
{\(\begin{aligned}[t] x& =\phantom{ -}4 + t, &
\\y& =\phantom{ -}2,
\\z& = -1 - t;\ t\in \mathbb{R}
\\ \end{aligned}\)}{\(\begin{aligned}[t] x& = 2 + 2m, &
\\y& = 2 +\phantom{ 2}m,
\\z& = 1 -\phantom{ 2}m;\ m\in \mathbb{R}
\\ \end{aligned}\)}{\(\begin{aligned}[t] x& =\phantom{ -}4 + 2k, &
\\y& =\phantom{ -}2 +\phantom{ 2}k,
\\z& = -1 -\phantom{ 2}k;\ k\in \mathbb{R}
\\ \end{aligned}\)}{\(\begin{aligned}[t] x& = -2 + 2u, &
\\y& =\phantom{ -}2,
\\z& =\phantom{ -}1 - 2u;\ u\in \mathbb{R}
\\ \end{aligned}\)}{}{}}
\LANGpl{
\Question{
Dane są punkty \(A = [0;0;1]\),
\(B = [2;0;-1]\) i
\(S = [2;1;0]\),
wyznacz równanie parametryczne obrazu prostej 
\(AB\) w symetrii względem 
punktu \(S\).}
{\(\begin{aligned}[t] x& =\phantom{ -}4 + t, &
\\y& =\phantom{ -}2,
\\z& = -1 - t;\ t\in \mathbb{R}
\\ \end{aligned}\)}{\(\begin{aligned}[t] x& = 2 + 2m, &
\\y& = 2 +\phantom{ 2}m,
\\z& = 1 -\phantom{ 2}m;\ m\in \mathbb{R}
\\ \end{aligned}\)}{\(\begin{aligned}[t] x& =\phantom{ -}4 + 2k, &
\\y& =\phantom{ -}2 +\phantom{ 2}k,
\\z& = -1 -\phantom{ 2}k;\ k\in \mathbb{R}
\\ \end{aligned}\)}{\(\begin{aligned}[t] x& = -2 + 2u, &
\\y& =\phantom{ -}2,
\\z& =\phantom{ -}1 - 2u;\ u\in \mathbb{R}
\\ \end{aligned}\)}{}{}}
\LANGcs{
\Question{
Jsou dány body \(A = [0;0;1]\) ;
\(B = [2;0;-1]\) a
\(S = [2;1;0]\).
Určete parametrické vyjádření přímky, která je obrazem přímky \(AB\) ve středové souměrnosti se středem v bodě \(S\).}
{\(\begin{aligned}[t] x& =\phantom{ -}4 + t, &
\\y& =\phantom{ -}2,
\\z& = -1 - t;\ t\in \mathbb{R}
\\ \end{aligned}\)}{\(\begin{aligned}[t] x& = 2 + 2m, &
\\y& = 2 +\phantom{ 2}m,
\\z& = 1 -\phantom{ 2}m;\ m\in \mathbb{R}
\\ \end{aligned}\)}{\(\begin{aligned}[t] x& =\phantom{ -}4 + 2k, &
\\y& =\phantom{ -}2 +\phantom{ 2}k,
\\z& = -1 -\phantom{ 2}k;\  k\in \mathbb{R}
\\ \end{aligned}\)}{\(\begin{aligned}[t] x& = -2 + 2u, &
\\y& =\phantom{ -}2,
\\z& =\phantom{ -}1 - 2u;\ u\in \mathbb{R}
\\ \end{aligned}\)}{}{}}
\LANGsk{
\Question{
Sú dané body \(A = [0;0;1]\) ;
\(B = [2;0;-1]\) a
\(S = [2;1;0]\).
Určte parametrické vyjadrenie priamky, ktorá je s priamkou
\(AB\) stredovo súmerná
podľa bodu \(S\).}
{\(\begin{aligned}[t] x& =\phantom{ -}4 + t, &
\\y& =\phantom{ -}2,
\\z& = -1 - t;\ t\in \mathbb{R}
\\ \end{aligned}\)}{\(\begin{aligned}[t] x& = 2 + 2m, &
\\y& = 2 +\phantom{ 2}m,
\\z& = 1 -\phantom{ 2}m;\ m\in \mathbb{R}
\\ \end{aligned}\)}{\(\begin{aligned}[t] x& =\phantom{ -}4 + 2k, &
\\y& =\phantom{ -}2 +\phantom{ 2}k,
\\z& = -1 -\phantom{ 2}k;\  k\in \mathbb{R}
\\ \end{aligned}\)}{\(\begin{aligned}[t] x& = -2 + 2u, &
\\y& =\phantom{ -}2,
\\z& =\phantom{ -}1 - 2u;\ u\in \mathbb{R}
\\ \end{aligned}\)}{}{}}
\LANGes{
\Question{
Dados los puntos \(A = [0;0;1]\),
\(B = [2;0;-1]\) y
\(S = [2;1;0]\),
halla las ecuaciones paramátricas de la imagen de la recta
\(AB\) en la simetría central con respecto al punto \(S\).}
{\(\begin{aligned}[t] x& =\phantom{ -}4 + t, &
\\y& =\phantom{ -}2,
\\z& = -1 - t;\ t\in \mathbb{R}
\\ \end{aligned}\)}{\(\begin{aligned}[t] x& = 2 + 2m, &
\\y& = 2 +\phantom{ 2}m,
\\z& = 1 -\phantom{ 2}m;\ m\in \mathbb{R}
\\ \end{aligned}\)}{\(\begin{aligned}[t] x& =\phantom{ -}4 + 2k, &
\\y& =\phantom{ -}2 +\phantom{ 2}k,
\\z& = -1 -\phantom{ 2}k;\ k\in \mathbb{R}
\\ \end{aligned}\)}{\(\begin{aligned}[t] x& = -2 + 2u, &
\\y& =\phantom{ -}2,
\\z& =\phantom{ -}1 - 2u;\ u\in \mathbb{R}
\\ \end{aligned}\)}{}{}}
\End
